% =====================================================================
%  Harald Høffding, ENGLISH PHILOSOPHY IN OUR TIME.
%  English translation of: Den engelske Philosophi i vor Tid.  Kjøbenhavn: P. G. Philipsens Forlag, 1874
%  (Trykt hos J. H. Schultz).  [IV], 190 pp.
%
%  GENERATED; DO NOT EDIT BY HAND.  Built by ebook-method/engelske-philosophi/tr_parts.py build --write from
%  tr_template.tex and the parts in ebook-method/engelske-philosophi/tr/ (translated 2026-10-15 from
%  texts/hoeffding/engelske-philosophi/transcription.tex, following TRANSLATOR-BRIEF.md and GLOSSARY.md there).
%  Page numbers in the margin [N] are those of the Danish print.  Footnote marks *) **) ***) as printed.
%  Terms (GLOSSARY.md): Udvikling = development, evolution for Spencer's and Darwin's theory (Udviklingsphilosophien =
%  the Philosophy of Evolution); Forestilling = idea; Aand = mind (spirit for Hegel and religion); Jeget = the self;
%  Tillempelse = adjustment; det Ubetingede = the Unconditioned.  Høffding's Danish renderings of English authors are
%  translated from his Danish, not replaced by the originals.
%  ODD SPOTS: translating was also a second reading of the Danish.  The translators listed 208 places that looked
%  wrong or unclear (tr/odd-*.tsv); all were read blind at the image (the transcription's header, (g)).  71 were
%  transcription faults, mended in the Danish (e.g. Øvelse p. 178, Paavirkninger p. 180, Plante- eller p. 138); 53 are
%  the print's own misprints, kept in the Danish (PRINTED AS IS) and translated here as the sense requires; the rest
%  were found to stand as printed, or were questions of grammar or reference settled in translating.
% =====================================================================
\documentclass[12pt,a4paper,openany]{book}
\usepackage[utf8]{inputenc}
\usepackage[T1]{fontenc}
\usepackage{libertinus}
\usepackage[english]{babel}
\usepackage{setspace}
\usepackage[margin=1.2in]{geometry}
\usepackage{fancyhdr}
\setlength{\headheight}{28pt}
\addtolength{\topmargin}{-13.5pt}
\usepackage{hyperref}
\hypersetup{pdftitle={English Philosophy in Our Time (1874)}, pdfauthor={Harald Høffding}, hidelinks}
\pagestyle{fancy}
\fancyhf{}
\fancyhead[LE,RO]{\thepage}
\fancyhead[RE]{\textit{Harald Høffding: English Philosophy in Our Time}}
\fancyhead[LO]{\textit{1874}}
\fancypagestyle{plain}{\fancyhf{}\fancyhead[LE,RO]{\thepage}}
\setstretch{1.3}
\usepackage{marginnote}
\usepackage{graphicx}
\DeclareUnicodeCharacter{0254}{\reflectbox{c}}
\renewcommand{\thefootnote}{\ifcase\value{footnote}\or *)\or **)\or ***)\else \arabic{footnote})\fi}
\newcommand{\opage}[1]{\marginnote{\footnotesize\textit{[#1]}}}
\newcommand{\chapterhead}[2]{\par\bigskip\begin{center}{\Large #1}\\[0.5em]\rule{2.5em}{0.4pt}\\[0.6em]{\large\bfseries\sffamily #2}\end{center}\medskip\addcontentsline{toc}{chapter}{#2}}
\newcommand{\closinghead}[1]{\par\bigskip\begin{center}{\Large\bfseries\sffamily #1}\\[0.4em]\rule{2.5em}{0.4pt}\end{center}\medskip\addcontentsline{toc}{chapter}{#1}}
\pdfstringdefDisableCommands{\renewcommand{\footnote}[2][]{}\renewcommand{\opage}[1]{}}

\begin{document}

\begin{titlepage}
\begin{center}
\vspace*{3em}
{\LARGE English Philosophy}\\[0.8em]
{\LARGE in Our Time.}\\[3em]
{\small By}\\[1em]
{\large Dr. \textbf{H. Høffding.}}\\[1.5em]
\rule{3em}{0.4pt}
\vfill
{\bfseries Copenhagen.}\\[0.4em]
P. G. \emph{Philipsen's Publishing House.}\\[0.4em]
{\footnotesize Printed by J. H. \emph{Schultz.}}\\[0.4em]
1874.\\[2em]
{\footnotesize English translation of \textit{Den engelske Philosophi i vor Tid} (1874).}
\end{center}
\end{titlepage}
\clearpage
\thispagestyle{empty}
\begin{center}{\large\bfseries \emph{Contents.}}\end{center}
\begin{center}\rule{3em}{0.4pt}\end{center}
\hfill{\small Page.}\par
\noindent\emph{Chapter One:} \textbf{General Characterization}\dotfill\ 1.\par
\begin{quote}\small The Conditions of Philosophical Knowledge. --- The Earlier English Philosophy. --- The English Philosophy of Our Time and its Relation to Other Schools.\end{quote}
\noindent\emph{Chapter Two:} \textbf{Pure Empiricism}\dotfill\ 18.\par
\noindent\hspace*{2em}I. \textit{John Stuart Mill}\dotfill\ 18.\par
\begin{quote}\small Historical Presuppositions (Hartley and James Mill). --- Logic and Psychology. --- The Idea of the Absolute. --- Ethics and Sociology.\end{quote}
\noindent\hspace*{2em}II. \textit{Alexander Bain}\dotfill\ 66.\par
\begin{quote}\small Sense-perception and Thought. --- The Classification of the Feelings. --- Ethics. --- Will and Belief.\end{quote}
\noindent\emph{Chapter Three:} \textbf{The Critical School in England}\dotfill\ 88.\par
\noindent\hspace*{2em}I. \textit{William Whewell}\dotfill\ 90.\par
\begin{quote}\small The Philosophy of the Inductive Sciences. --- Relation to Kant and to Theology.\end{quote}
\noindent\hspace*{2em}II. \textit{William Hamilton}\dotfill\ 99.\par
\begin{quote}\small Polemic against the Philosophy of the Unconditioned. --- Hamilton's Disciple \textit{Mansel} and his Critique of Theology. --- Hamilton's Logic and Psychology. --- The Doctrine of Causality.\end{quote}
\noindent\emph{Chapter Four:} \textbf{The Philosophy of Evolution}\dotfill\ 116.\par
\begin{quote}\small \textit{Herbert Spencer}. --- General Characterization. --- The Idea of the Absolute (the Relation between Religion and Science). --- The First Principles. --- Biology. --- Psychology and Theory of Knowledge. --- Ethics and Sociology.\end{quote}
\noindent\textbf{Conclusion}\dotfill\ 184.\par
\begin{center}\rule{3em}{0.4pt}\end{center}
\clearpage

% --- p. 1 ---
\setcounter{footnote}{0}%
\chapterhead{Chapter One.}{General Characterization.}
\opage{1}


Human development does not proceed in a straight line.
In every domain we find that times full of
fortunate and magnificent progress alternate with apparent
stagnation or retrogression, in which the forces stir in
concealment or work in other domains. And this holds
to a higher degree, the more composite an endeavor is, the
more conditions it presupposes in order to succeed. Of
none, therefore, does it hold more than of the philosophical
striving. To philosophize is not the same as to give
oneself up to immediate intuitions and brilliant attempts
to grasp and present the Ideal. If philosophizing
is really to have scientific significance, a clear and certain
consciousness of the nature of truth and of reality must
be laid at the foundation of the whole endeavor. The
True and the Real must be powers before which thought
bows, because in them it recognizes its own law and in
submission to them its true freedom. The striving
of thought must aim at letting itself be bound step by step;
only by being bound throughout does thought have
validity. All philosophizing therefore presupposes a double
system of conditions. Thought must be awake, the drive
toward a coherent and all-embracing view of the world
must be alive; but on the other hand the factual
foundation, the mass of particulars, must also lie so
prepared that it becomes really possible for thought to get
beyond its subjective and accidental ideas and to assume
% --- p. 2 ---
\setcounter{footnote}{0}%
\opage{2}the forms of truth and of reality. According as these
conditions alternate, the epochs of philosophical inquiry alternate also.
This lies in the matter itself and is confirmed by history. Recent
German philosophy, which was introduced in so energetic and brilliant a
manner by Kant, and which in him, its greatest
representative, still stood in close and fruitful interaction with
experience, exhausted all its possibilities in systems in which
the given fund of determinations of reality was so far absorbed
that the law of existence was finally made one with
the dialectical movement of thought itself. Then came a
time in which men were inclined to throw all philosophy overboard;
interest turned toward reality; it was a
brilliant period for exact, experimental natural science;
material was gathered in all domains. But the old
problems could not be kept out; they press forward again
with irresistible force, and we find, in part even precisely
on the side of such men as took part in experimental
research, a series of new attempts to treat the
philosophical tasks. A similar course of development we find
in England. After a magnificent philosophical activity
in the 17th and the first part of the 18th century,
England was poor in philosophical attempts, until interest
was awakened again toward the middle of our century.
``There was a time,'' says Herbert Spencer (Essays. 1868.1.
p. 300), ``when our countrymen speculated --- and certainly
to as great an extent as any other people --- on
all the high questions that present themselves to the human
mind; and a glance at the philosophical systems
that are or have been dominant on the Continent suffices
to show how much other nations owe
to the discoveries of our forefathers. But for a generation or
two these more abstract subjects have been neglected, indeed
even (by those who boast of being `practical') despised.
Partly perhaps a consequence of our rapid material
growth, this intellectual phase is due in great part
to the exhaustion of material and the necessity of better
data. Not so much with conscious recognition of the
% --- p. 3 ---
\setcounter{footnote}{0}%
\opage{3}end to be served, as with unconscious subordination
to the rhythm that can be traced in social changes
as well as in other things, a period of theorizing
without observing has been followed by a period of observing
without theorizing. During the long-continued devotion
to concrete science an immense
quantity of raw material has been accumulated for abstract science, and now a period
is evidently about to begin in which this accumulated
raw material will be organized into a coherent theory.
At all points, in the inorganic sciences, in the science
of life and in the science of society, we may observe
a tendency to pass from the superficial and
empirical to the deeper and rational.'' When we speak of
the English philosophy of our time, it is thus more than
an accidental series of phenomena that we bring forward; it is
a series of endeavors that arise in consequence of a
natural law of development and therefore with a peculiar character
conditioned by their place in the development.

The double presupposition of philosophical
inquiry, which has the effect that its course of development becomes rhythmical
or almost leaping, also contains the possibility
of a decisive opposition within philosophy. Two
different points of departure become possible. The general
task is that thought shall appropriate the True and the Real.
But, it may then be said, only by and for thought
is anything real and true; the essence of thought is therefore an essential
and self-valid element in reality; what actually
exists is only a means for the development of thought, the knowledge
of it only the way by which thought becomes conscious
of itself, and in this consciousness the highest
truth is contained. Here, then, the emphasis is laid on the ideal
element, on the original activity of thought as the last presupposition
for all knowledge. But, it is said from the other
side, a knowledge without real, objectively valid content
is merely subjective and illusory. When thought wishes to be
self-active, it keeps within its own self-created ideas,
in so far as it does not follow traditional prejudices. If
% --- p. 4 ---
\setcounter{footnote}{0}%
\opage{4}reality is to be known, thought must therefore behave
altogether passively and receptively. All knowledge must come about by
the gathering of impressions and experiences, and only that can lay
claim to truth which can be traced directly or indirectly
back to this source. --- The opposition between the two
main tendencies within philosophy rests, then, on the
different answer to the question: Wherein does the criterion
of truth and of reality ultimately consist?
The one answer refers us to the necessity of thought, the
other to the reality of experience. This is the fundamental opposition;
but on the basis of it there unfolds a multitude of
derivative oppositions, of extraordinary importance for the whole view of the world. The question must then present itself whether
it is possible to reconcile the two principles and with them the two
views of the world.

However, this question can come forward only when
an interaction, a real debate, between
the two philosophical schools has been initiated. But this may in part be prevented
by their belonging to different nations. An opposition as
decisive as the one that has here developed
with regard to the whole direction of philosophical inquiry
points back to more comprehensive psychological and national
differences. Mode of thought and view of the world are not
anything isolated; they are factors that stand in living connection
with the whole character of the individual and of the people. The
idealistic school has in modern times had its home chiefly
in Germany, the realistic in England. Germans
and Englishmen are the real philosophizing nations
of modern times, they stand in the history of philosophy beside
the Greeks. National limitation has here acted
as a hindrance, but not merely as a hindrance. For in philosophical
inquiry what matters essentially is clarity and consistent
carrying through of the line of thought, and to this is required
an exact and thorough working out of the fundamental ideas, which
is possible only on the condition of a certain one-sidedness.
On this rests the significance of the different philosophical schools,
and on this rests also the significance of the influence of the individual
% --- p. 5 ---
\setcounter{footnote}{0}%
\opage{5}and the national character upon philosophical
thought. Important and decisive fundamental concepts would never
have been thought through without such limitations and
one-sidednesses. Thus German and English
philosophy in particular, by their consistent development proceeding
step by step, have cast a clear light on great circles of thought
that would otherwise not have been measurable. And only
thereby can that fundamental opposition between the main
philosophical tendencies present itself in so clear and precise
a form that work toward overcoming it becomes
possible. ---

From what has been set forth so far it follows that there are three essential
conditions for progress in philosophical knowledge.
Philosophy presupposes the reality given in experience as its material. It aims at a consistent
carrying through of the principle. But since this principle, according
to the nature of the case, is developed one-sidedly in the different schools,
it strives to correct this one-sidedness by a union
of the apparently opposed standpoints. --- It shall now
be shown that the English philosophy of our time is of great significance
in all three respects.

English philosophy has at all times had an
empirical and realistic character. In the Middle Ages we find
Englishmen as the most prominent representatives of
Nominalism, the tendency within Scholasticism that corresponds
to modern Realism. Duns Scotus († 1308) polemizes
against Thomas Aquinas, the representative of ecclesiastical
speculation, and asserts in particular the significance of individuality.
According to Thomas the universal idea was true
reality, and the individual came about only through a limitation
that was due to matter. But, says Duns Scotus,
that which makes a thing this determinate thing must be
something positive; the individual is the perfect, the
highest reality. Roger Bacon († 1292) scoffed at Albertus
Magnus and Thomas Aquinas because they wanted to philosophize
without having learned languages, mathematics and physics,
and he had to languish in his prison because his zealous
% --- p. 6 ---
\setcounter{footnote}{0}%
\opage{6}occupation with, and praise of, scientia experimentalis made
him suspected of magical arts. William Occam
(† 1347) carried Nominalism further. Only the particular
is real; it cannot be explained and dissolved by universal
determinations of concepts; the universal concepts have no reality
outside the understanding; they are signs and symbols that
express the sum of the individual phenomena. All knowledge
therefore rests on intuition, which in turn points back
to inner or outer experience. Since we cannot experience or
intuit the essence of God, we cannot therefore attain any
knowledge of God; theology is no science, but
rests solely and alone on supernatural authority. Occam
repeats what Duns Scotus had already expressed, that
something may be true in theology that is not so in
philosophy. --- The nominalistic or empirical tendency
in this way bursts Scholasticism asunder. The opposition
that appears in more recent times between English philosophy
and that of the Continent we can thus trace back to
the Middle Ages.

Bacon of Verulam († 1626) is the founder of modern
English philosophy. From speculative and theological
studies he calls the human mind back to nature and
its knowledge. Only there is a secure foundation to be found
for science. Left to itself, thought strays into
illusory and arbitrary presuppositions (idola mentis), which
arise partly from human nature, partly from social conditions.
From these unreliable presuppositions one frees
oneself by submitting to the law of experience, by examining
the particular by way of induction, in order thereby
to arrive at the true fundamental concepts (notiones rerum primæ).
Inductive natural science is for Bacon the mother of all sciences,
to which logic and moral philosophy too shall
be traced back. Its aim is not merely knowledge, but
progress in civilization, in that human life by its inventions
shall be more richly furnished with goods of every kind,
and man's power over nature extended. --- When Bacon
thus demands that man shall strip off all
% --- p. 7 ---
\setcounter{footnote}{0}%
\opage{7}self-activity in knowledge and behave as passively as possible,
he thereby paves the way for sensualism,
which derives all knowledge from sense-perception and regards all
ideas and thoughts as transformed sensations of sense.
This doctrine was developed by Thomas Hobbes († 1679) and
John Locke († 1704). The laws according to which sensations
are combined and compounded and thereby transformed
into thoughts could be none other than the laws of
the association of ideas; when mental activity is rejected, there remain only
the passive phenomena of the mind. David Hartley
(† 1757) was the first who sought to explain all phenomena of consciousness
as results of the association of ideas.
Finally David Hume († 1776) drew the last consequences
of empiricism and sensualism. Both the inner
and the outer world, both the self and the world, we know
only through our sensations of sense and ideas;
and the inner sense, the circle of ideas, is only a faded
image of the outer sense; in the last instance we know
then only the series of sensations of sense. But then it is
founded only on an illusion when we assume a self, an
inner mental substance as the bearer of the circle of ideas,
or an outer material substance as the cause of sensation of sense. It is altogether an illusion when we think we
know a causal connection, a relation between cause
and effect; we know only the succession of phenomena in
time, no inner and necessary connection between
them. It is only habit, which in turn is a product of the association of ideas,
that leads us to such assumptions; but habit
cannot ground any knowledge.

By sharp and consistent carrying through, empiricism had thus
passed over into skepticism and had thereby
dissolved itself. Nevertheless it was the principles of Hume and Hartley
on which English philosophy leaned
when, toward the middle of the 19th century, a period of flowering
began again for it. We must therefore now
state more precisely the relation between the English philosophy of our time
and the earlier. The fundamental thought is the same: the
% --- p. 8 ---
\setcounter{footnote}{0}%
\opage{8}nominalistic adherence to the particular and individual
as the only real, the empiricist assumption of
sense-perception and induction as the only sources of knowledge. But
this fundamental thought receives a peculiar carrying through
through the weight that is laid on the laws of the association of ideas as
decisive for the growth and development of the life of consciousness
and in particular of knowledge. Hume had said that we cannot
attain any true knowledge, because we have only sensations of sense
and their connections to build on. But,
the English philosophers of our time say to this, knowledge itself
is nothing but a firm and indissoluble association of ideas.
Thought is not a mystical force that in turn discovers a
mystical bond as necessary connection between things.
When certain ideas, in accordance with experience and the whole course of development,
have entered into a firm connection, then we must
regard them as belonging together. What are we ourselves,
our self, but a firm and intimate composition,
a chain of associated ideas? What we call necessary
is only an expression of firmness in the association.
But we can not only show psychologically how consciousness
actually develops according to the laws of the association of ideas;
we can also by logical analysis distinguish between justified
and unjustified associations. Bacon had in general terms
praised induction and combated the authority of the syllogism
as the highest form of knowledge. But only now
is inductive logic carried through, in that it is shown how
the experimental method leads to the elimination of
the accidental connections of thought and to the establishment and maintenance of those that are
universal and therefore justified. Only on the basis
of such a preceding, experimentally established linking together
of ideas do universal laws become possible, and only
thereby are starting points won for deduction, syllogism
and mathematical analysis. Hume polemicized against natural science,
but the English philosophy of our time places itself
precisely by virtue of Hume's principles on the ground of
natural science. Experimental natural science is in its essence
skeptical. It will find only the laws of the mutual
% --- p. 9 ---
\setcounter{footnote}{0}%
\opage{9}connections of phenomena, their succession and their coherence.
Only because Hume did not know the experimental
method could he turn his skepticism against natural science itself.\footnote[1]{Stuart Mill (Examination of Sir William Hamiltons philosophy p. 555) even holds that Hume's skepticism was only a disguise: ``he would rather be called a skeptic than be designated by a more odious name.''}


English philosophy in our time here recalls
French positivism. One finds even very often
positivism used as a common name, so that it also comprises
the English school. And yet the first representatives of this school
have with justice disclaimed the name of
positivists, when this word is to designate them as disciples of
Auguste Comte. We have already seen that the English
philosophy of our time is a further development of the standpoint and results
of the earlier English thinking, but in such a way that, in
accordance with its logical and psychological principles, it came to
enter into a close connection with experimental
research. This research and its results are the common
ground where it meets French positivism; what
is common to both tendencies thus derives from the source
from which they both draw. One must distinguish between positivism
as a particular school, and what the French call l'esprit
positif, a tendency of mind that has sprung from the whole of modern
culture and science, and of which positivism
itself is only a particular form. The opposition between the
French positivism and the English school may in its generality
be stated thus: positivism goes from the special sciences
to philosophy, the English school from
philosophy to the special sciences. French positivism
properly will not recognize any other problems than the
special-scientific. It collects in a classificatory encyclopedia
the principles and results of the individual sciences
and is content with that. Only one decisive problem
does it bring forward at last, and thereby involuntarily betrays
% --- p. 10 ---
\setcounter{footnote}{0}%
\opage{10}its own innermost motive: the problem of the development of humanity and
in particular of the human life of society (the sociological
problem). But it neglects the logical and
psychological investigations which alone can form a basis
for a philosophy, and which are also a necessary presupposition
for the treatment of that great problem. In any
case it does not recognize the independence of psychology
but lets it go wholly under physiology (so
Littré). English philosophy, on the contrary, takes its point of departure
in logic and psychology and finds from there the way
to the more special problems. ---

The consideration of the English philosophy of our time in its
relation to the earlier English thinking led us also
to touch on its relation to the special sciences. But the
most important contribution to the characterization of this philosophy remains
still to come. It consists in a statement of its relation
to the main philosophical tendency to which the
earlier English school had formed a decisive opposition,
the idealistic or speculative.

What we have for brevity's sake designated as the
English school has no exclusive right
to this name, in so far as, both earlier and in our days,
philosophical tendencies have asserted themselves in England that
stand in opposition to the empirical-realistic philosophy.
In the 17th century we find such a tendency
in the mystical Platonism that prevailed at the University
of Cambridge, and whose most notable representatives
were Henry More and Ralph Cudworth; Locke's polemic against
the ``innate ideas'' is perhaps directed rather against them
than against the Cartesians. In the 18th century there developed
at the Scottish universities, as a reaction against
English empiricism and its results, a particular Scottish
school, whose real founder was Thomas Reid († 1796).
This school had in part practical motives; it wished to
ward off the allegedly dangerous consequences of the doctrine of Locke and
Hume and was, as Cousin expresses himself, ``a protest
of sound human understanding and of conscience against
% --- p. 11 ---
\setcounter{footnote}{0}%
\opage{11}Locke's philosophy.'' Reid contests the empirical school's
doctrine of the sensations of sense as the original elements,
by whose association and combination the whole mental life
develops. He teaches that our consciousness immediately
apprehends external objects. We pass judgments
instinctively about things; and these involuntary assumptions contain
the first and natural principles of all our knowledge.
They are the original facts of consciousness,
whose sum makes up the so-called ``common human understanding''
(common sense). Of such ``natural principles''
Reid enumerates a great many, among others: the proposition
that every sensation points to an external object;
belief in our personal identity; further, mathematical
and metaphysical principles (e.g.: Everything has its
cause), and moral commands (e.g.: One shall set duty
above utility). All these principles are supposed to lie immediately
in consciousness. --- It is not by setting up ``the principles
of sound human understanding'' but by the investigation
of empiricism that they set in motion, and by the sense for
psychological investigations that they awakened, that Reid and his
school have exerted no small influence. In France
Royer Collard and Victor Cousin took the field against
Condillac with Reid's weapons. But in England too
Scottish philosophy has been of great significance, in particular
since, in its greatest representative in our time, William
Hamilton, it entered into a peculiar combination with
Kantian philosophy. This leads us to inquire how
German philosophy intervenes in the course of development of English
thought in our time. Kant, who by his own
declaration received the impulse to his epoch-making inquiry
from David Hume, has since himself reacted significantly
upon Hume's fatherland.

Kant, like Reid, steps forth against Hume's skepticism.
But whereas for the anxious Reid it was above
all a matter of securing the reality of the self and of the external world,
on which account his attempt also took on a more superficial
character, Kant went deeper and found in the problem of causality
% --- p. 12 ---
\setcounter{footnote}{0}%
\opage{12}the real point at issue between the two tendencies.
He could not content himself with the Scottish
school's appeal to sound human understanding. That
the common human understanding knew and applied
the concept of cause, Hume did not deny; but the dispute turns
precisely on the significance and validity of this concept. In
any case one ought, thought Kant, rather to use
one's sound human understanding to give proof of what
one assumes than to invoke it as an oracle. In thinking
the matter through Kant became aware that
there is a whole circle of forms of intuition and concepts which
lie at the foundation of all experience, because they are the different
forms and ways in which the material given in sensation
is gathered into unity for intuition and thought. These
forms must therefore, he concluded, be a priori, originally
given in our nature. They cannot derive from experience,
since all experience precisely presupposes them, and since they have universal
and necessary validity, which mere abstractions
from experience cannot have, since experience never
shows us the Unconditioned. Precisely by virtue of the validity of
experience Kant thus asserts the originality of thought.
He has hereby paved the way for an entirely new conception of
a priori or speculative knowledge. This can
no longer, as in dogmatism, be a finished and
closed system, which was supposed to be developable without regard to
experience; but it always presupposes experience, and
takes it up into itself by virtue of its own principle: ``That all our
knowledge begins with experience,'' says Kant,
``of that there is no doubt; but it does not therefore follow that
all our knowledge springs from experience.'' Kant himself distinguished
so sharply between the subjective and the objective that it
looks as if the forms of thought and intuition stand ready-made
on the one side, and the content only comes in
from the other side; but the true consequence of the
critical philosophy Schelling has expressed in the proposition that
our knowledge is just as little a priori as a posteriori,
since this opposition is due only to reflection. Schelling
% --- p. 13 ---
\setcounter{footnote}{0}%
\opage{13}shows further that there are indeed a priori but not innate
concepts. ``Not concepts, but our own nature and
its whole mechanism is innate in us. This nature is of
a determinate constitution and acts in a determinate way,
but perfectly unconsciously, for it is itself nothing other
than precisely this acting.'' The soul is neither a written
nor an unwritten tablet.

It was Samuel Coleridge († 1834) who became above all the
intermediary through which German philosophy made
its influence felt in England. In his youth
he had been an ardent admirer of Hartley and Hume, but
was later repelled by these thinkers' denial of the activity of the mind.
He saw rightly that when the doctrine of the association of ideas
is made the main psychological doctrine, all phenomena
of consciousness are conceived as completely passive.
From this standpoint he has already (in his Biographia
litteraria vol. I. pars II. p. 116 ff) given a striking
critique of the so-called ``Association psychology.'' The study
of Kant and his successors led him quite beyond
empiricism. He became so at home in the principles of the critical
philosophy that on his own way he reached
several of the ideas which Schelling expressed in his system of identity,
so that this even brought upon him an accusation of plagiarism.
As a poet Coleridge worked to introduce Romanticism
into England; as a thinker and publicist he worked
to spread knowledge of the fundamental thoughts of German philosophy.
Even so decided an opponent of these as
Stuart Mill places him beside Bentham and regards
these two men as the two fertilizing minds of their
time. Mill acknowledges that the attacks on Locke's philosophy
from the German as well as from the Scottish school have essentially
spurred on psychological inquiry and called attention
to the difficulties that were still unovercome.
(Dissertations and discussions I.). Buckle mentions
in his history of civilization, besides the influence of the Scottish school,
also ``the sudden and well-deserved admiration for German
literature, to which Coleridge especially gave expression, and which
% --- p. 14 ---
\setcounter{footnote}{0}%
\opage{14}instilled into the English mind a taste for higher and bolder
views than it had hitherto known''; but he unfortunately did not live
long enough to complete a description of
this movement.

The influence which critical and speculative
philosophy exerted on empiricism consisted, first, in
leading it to a sharper consciousness of its own principle
and a more consistent carrying through of it, whereby, on
the other hand, the conviction was bound to develop
that empiricism necessarily points beyond itself. Hume
had rejected the concept of causality, but he acknowledged
the propositions of mathematics as necessary, because they were grounded
on the principle of contradiction and not on external experience.
Kant remarked to this that Hume consistently must also
have denied necessity in the mathematical
domain. This Stuart Mill acknowledges and teaches that all
mathematical axioms are grounded on experience, on firm
and constant association of ideas, and that the principle of contradiction
is merely the expression of the most abstract result of experience.
Only thereby has empiricism reached its completion.
But when all necessity is thus denied, the
question must naturally be raised how empiricism
is to be able to prove its own truth. Herbert Spencer
maintains here against Stuart Mill that, if the whole of empirical
knowledge is not to hang in the air, there must be a criterion
which from the first and throughout lies at the foundation of
it, and by virtue of which every advance in knowledge
is made. ``To say that there are no necessary
truths is the same as to say that the conclusion is not
necessarily contained in its premises but is an
empirical judgment, which is not certainly true. Pure empiricism
starts from an uncertainty and moves through
a series of uncertainties. One can stop it at every
step, as long as it has not accounted for its original
presupposition.'' Before one assumes anything to be true, one must
make clear to oneself what the condition is, in general,
for something to be true. Consciousness must, by
% --- p. 15 ---
\setcounter{footnote}{0}%
\opage{15}examination of its own states, come to a decision
as to which property it really is that is meant when anything
is posited as true.

One of the difficulties in the psychology of the earlier empirical
school was the one emphasized by Coleridge, that the
activity of the mind remained unexplained. Here English
psychology has taken a step which Stuart Mill aptly compares
with the step by which Laromiguiére went beyond
Condillac, in that he emphasized the activity of the mind in
attention. Supported by physiological facts,
Alexander Bain, an English psychologist who stands
close to Mill, has indeed demonstrated an original active element as the point of departure
for mental life, namely the primitive, automatic
movement, which precedes all sense-perception,
which in the individual sense-perceptions is guided in determinate directions,
but without which these sense-perceptions cannot take
place, --- the first germ of the activity of thought and of will.


The eye having thus been opened to sides of reality
to which men had formerly been blind, it was at the same time
led to bolder and more comprehensive points of view. Beside
the philosophers there appear here the historians and the natural scientists.
The idea of the development of the human race and of the human mind
in the different domains makes its way through
the philosophy of history founded by Kant and Herder; the
speculative philosophy of nature and comparative natural science
show that the idea of evolution is not limited in its
validity to the world of mind, but that it embraces nature
and mind as links of the same chain. The older school's
anxious and narrow psychological empiricism vanished before
these magnificent views. The task then necessarily had to
present itself of gathering all individual investigations
into a comprehensive view of the world, governed by
one guiding principle. It is a task for which there are
only scattered contributions in the earlier English thinkers,
but which the great philosophical systems on the Continent
were attempts to solve. The English
% --- p. 16 ---
\setcounter{footnote}{0}%
\opage{16}philosophers of our time too maintain that our knowledge is only relative and
phenomenal and cannot grasp the Absolute; but relative
knowledge itself can nevertheless be organized into a coherent
whole, and since it is impossible for thought to hold the Relative
dualistically separated from the Absolute, which it must posit in one
form or another, we see also the most profound
of these philosophers, Herbert Spencer, acknowledge
that in the measuring out and articulation of the domain of relative
knowledge indirect determinations are won of
the absolute force, to which the idea of evolution, the ruling
idea in the realm of phenomena, points back.

In this way empirical philosophy
approaches speculation, partly by acknowledging the
original presuppositions on which speculation too
builds (the a priori criterion of truth; the original
activity of the self), partly by the endeavor after an all-embracing view of the world,
which in the conditioned finds determinations for
the Unconditioned, in the Relative indirect forms of the
Absolute, --- it approaches in particular that tendency in the German philosophy of our
time which may be called the idealism of reality,
because it proceeds from the consideration of reality
and arrives at the result that this is in the last instance
of an ideal nature.\footnote[1]{See my book ``Philosophien i Tydskland efter Hegel'' pag. 220---319.} A similar course we find in the English philosophy
of the present day. Here among the English there is a
stronger empirical and skeptical tendency, but the practical
idealism that is so characteristic of these apparently
so prosaic inquirers must in the end also bring
speculative consequences with it, in so far as the theoretical
interest does not already lead in this direction. Also
as regards the German ``idealism of reality'' it holds
that the Absolute is essentially conceived as a scientific
and practical ideal. The time seems on the whole to be past
when men believed that by the way of dialectical development of concepts
% --- p. 17 ---
\setcounter{footnote}{0}%
\opage{17}they could grasp the Absolute. Even the highest and most clearly
drawn concept of God is yet only a symbol, a limiting concept,
which we arrive at, but which we cannot carry through.
Yet it must be maintained that we arrive at this concept
of necessity. That it must be a hypothesis
is true; but it is a hypothesis we are compelled in one form or another
to set up, and the question therefore
cannot be put aside in philosophy which hypothesis
comes nearest to the truth. Speculative criticism has
work enough here, a work that will never end, as long as
the human mind does not grow tired of occupying itself with its highest
ideas. This too is one of the points in which English
philosophy differs from French positivism. August
Comte wished indeed to abolish all theology and metaphysics, but
nevertheless posited an Absolute, a ``highest being,'' which coincided
with humanity itself, regarded as a whole
in past, present and future, in whose service and for whose
development all individuals are to work and wholly sacrifice
themselves. Against this holds the objection that Spencer advances,
that the Unconditioned cannot be sought in any single
domain, but must be the universal force that stirs
in everything and lies at the foundation of all existence. ---

From this general characterization we can draw
the result, which in the following exposition will find its
closer justification, that English philosophy in our time
shows us how empiricism, precisely when it is carried through,
leads beyond itself. As we have seen, the English philosophy of
our time takes its point of departure in the principles of the earlier
English school, but by the development it gives
these principles it leads them beyond the limitation of the school.
We must therefore ascribe to it no small significance
for the further development of philosophy and the working out of its
universal problems. ---

As already remarked, by English philosophy we understand
in particular the empirical and realistic tendency.
But in order to give a clear picture of the philosophical
movement in England we must nevertheless also bring forward the
% --- p. 18 ---
\setcounter{footnote}{0}%
\opage{18}representatives of the opposition to this tendency. Having therefore first
occupied ourselves with Stuart Mill, the consummator of empiricism,
and with the psychologist Bain, who in essentials
adheres to him, we shall characterize in brief Mill's
opponents Whewell and Hamilton, who have been influenced above all
by the critical philosophy. The third section of our exposition
will give a survey of the doctrine of Herbert Spencer, which
at once has its presuppositions in the empirical and
the critical school, which most clearly and definitely leads beyond
the boundaries of empiricism, and which marks the grand
conclusion of the philosophical endeavors in England in
our time.

\chapterhead{Chapter Two.}{Pure Empiricism.}


I. \emph{John Stuart Mill} (b. 1806, d. 1873) received,
as one sees from his Autobiography, impressions already early
that became decisive for his whole theoretical
and practical direction. His father, James Mill, led him
early into the philosophy of the 18th century and taught him
to regard its doctrines not merely as theories, but
almost as expressions of his religion, because they alone were able
to combat the prevailing prejudices which hinder
the progress of humanity toward freedom and felicity. It
is through James Mill that the English philosophy of our time
stands in connection with that of the 18th century. He was
essentially formed by the study of Hartley and Hume and, like
these two thinkers, regarded the association of ideas
as the fundamental psychological law, a conception which
has already been prepared by Thomas Hobbes. Although Stuart
Mill has not dealt especially with psychology, this
psychological theory nevertheless underlies his whole system; to
it, as he himself constantly emphasizes, all his logical,
% --- p. 19 ---
\setcounter{footnote}{0}%
\opage{19}metaphysical and ethical investigations go back. We must
therefore, both in order to show the historical connection
between the English philosophy of the 19th and that of the 18th century
and in order to understand Stuart Mill's standpoint in particular,
briefly mention the earlier forms of the said
psychological theory.

In its development from Locke to Hume the philosophy of the 18th
century was led to eliminate the concept of the self and
to conceive mental life merely as a train of ideas and
sensations, a train which was then in turn sought to be
simplified by tracing all ideas back to sensations of sense.
The laws of mental life were then conceived not
as expressions and forms of the energy of the self, but as
formulas for the mechanical interaction between the individual
members of the train of mental phenomena, in the last instance
of sensations of sense. A complete doctrine of
these laws, the laws of the association of ideas, was first sought by David
Hartley (b. 1705, d. 1757) to give in his work ``Observations
on man.'' Hartley gives the theory of association a
physiological foundation, in that, adhering to Newton's
emission theory, which the great investigator of nature employed not merely
to explain the effects of light, but also to
explain the effects of the nerves, he traces all sense-perceptions and ideas
back to vibrations in the substance of the brain, evoked
by impressions from without. When now certain determinate
vibrations frequently return, there must arise in the substance of the brain
a disposition to repeat them under certain conditions.
What these conditions are, we learn from the laws of the association of ideas,
which we find by psychological analysis.
Whereas Hume still distinguishes between three different
laws of association: association by resemblance, coincidence
(contiguity) and relation of causality between the ideas,
Hartley traces these laws back to a
single one, the law of coincidence, of contact (contiguity):
two sensations of sense are connected or associated when the
impressions that evoke them occur simultaneously or immediately
after one another. When afterwards the one impression
% --- p. 20 ---
\setcounter{footnote}{0}%
\opage{20}recurs, there will be a disposition to evoke also
the idea corresponding to the other impression, in so far as
in the brain there will be a disposition to repeat
the vibration corresponding to this other impression. An idea
is only the remainder of an earlier impression, resting on
a weaker vibration, which is as it were the after-sound of the
strong vibration evoked by the impression itself. Through the
fusion of individual ideas by means of their
frequent repetition in union with one another there arises a
compound idea, which in the end seems to form an indissoluble
unity, because the individual elements of ideas of which it
is composed disappear from consciousness, in that each single one
of them is overcome by the sum of all the others. The
compound idea seems in the end to have no kinship at all
with the ideas and sense-perceptions from which
it has come to be, just as a medicine may in its taste be
wholly different from each single one of the substances of which
it is prepared. The compound ideas can now in turn
be compounded with one another, and so on to infinity.
Hence it comes that many ideas seem to arise as it were
of themselves, from within, which one has then, like Locke, derived
from a special mental faculty, reflection; but a continued
analysis will be able to trace all, even the highest and most compound
ideas, back to single ideas, and these
in turn to impressions of sense. On the feeling of pleasure and displeasure
the associations of ideas exert an educative power: the sensuous
sensations gradually give way to the ideal feelings, in that
the immediate objects of pleasure and pain recede before
ideas that do not in themselves bring with them sensuous
enjoyment or suffering. On this rests the ability to cultivate
one's will and to raise oneself above the sensuous. And in
this way the moral sense and the moral
ideas are formed. But on the other hand all prejudice
and all influence of authority also rests on the power of
the association of ideas.

Coleridge has rightly emphasized the strict consistency
with which Hartley carries through his theory. When
% --- p. 21 ---
\setcounter{footnote}{0}%
\opage{21}mental life is conceived as resting on mechanical connections
between individual elements, then, as in the connection
between material atoms, it can only be immediate contact
that brings about the connection. Therefore in
Hartley the psychological hypothesis also stands in close connection
with the physiological theory. For though there is not, as
Lewes emphasizes in his criticism of Hartley, any such
likeness between physiological and psychological phenomena
that we could learn anything about the latter from the former,
and our acquaintance with the sensations is greater than our
acquaintance with their physiological foundation; yet association psychology
rests precisely on this, that the mental
phenomena are conceived in analogy with the bodily ones. When
therefore the physiological hypothesis is given up, association psychology
has already lost its absolute character.

James Mill's (b. 1773, d. 1836) psychological doctrine
differs from Hartley's by the abandonment of the physiological
hypothesis and by a more determinate and careful carrying through
of the psychological analysis. His ``Analysis
of the phenomena of the human mind'' appeared in 1829 and
was in 1869 published anew by his son, furnished with exceedingly
interesting notes by the latter, by the psychologist Bain, by the philologist
Findlater and by the historian Grote. In these notes
corrections and explanations are given and the investigation is carried
further, so that one has here as it were a repertory
of empirical psychology in England. --- In James
Mill, as in Hartley, it is the law of contiguity (of contact)
that is laid down as the foundation. Sensations that are wont to
occur simultaneously or immediately after one another associate
with one another, and the ideas follow after or with one another
like the sensations of sense whose images they are. Besides
the frequency of the coincidence of certain sensations, it is
also the interest that this coincidence has for the individual
that strengthens the association; interest leads us to
take note of the frequent and constant connections,
for it is on them that we base the plans for our action. Soon
there forms itself in consciousness an opposition between
% --- p. 22 ---
\setcounter{footnote}{0}%
\opage{22}constant and variable relations. I have seen the bird fly from west to east as often
as from east to west, but
I have always seen the stone in the end fall to the earth when it
was thrown, never continue to fly in the opposite direction.
Two ideas that in this way always show themselves to
belong together enter into an indissoluble union, a fixed
association, and in this consists what we call belief or assumption.
When we see the one phenomenon occur (e.g.
the throw of the stone), we also expect that the other (the fall
of the stone) will follow. ---

The psychological mechanism that Hartley and James
Mill taught lies now, as already remarked, at the foundation of
Stuart Mill's whole doctrine. It forms the presupposition
to which he returns at every point. But
he does not accept the doctrine of association wholly in the form
in which those older investigators presented it; his own investigations
led him at several decisive points beyond
it, as one sees in particular from his notes to his father's
work. This brings then not a few inconsistencies into
his doctrine, in that he has not become aware that the corrections
he makes in the older theory often directly
lead to the complete abolition of its presuppositions,
which he nevertheless believes he can retain. As clear and
open a mind as this thinker is, as great a pleasure
and joy as he finds in debate with other standpoints,
and as unreservedly as he acknowledges, when he has
become aware of difficulties that his theory cannot
solve, still he does not notice that he has here abandoned his
original presuppositions. He mentions in his Autobiography
the struggle he had to go through in order to free
himself from the one-sided points of view that his father's strong influence
had early instilled in him; and we shall find
the traces of this struggle in the indicated inner contradictions in
his doctrine. Association psychology is the mighty weapon
he turns against his opponents, but it is at the same time a
weight that burdens him himself and makes him unable
to complete what he strives to attain. Quite
% --- p. 23 ---
\setcounter{footnote}{0}%
\opage{23}in general that contradiction may be said to arise
from the fact that the theory of association rests on an external and
mechanical conception of mental life, which is indeed justified as
an empirical foundation, but which overlooks a remainder that here is
the real kernel of the matter, namely the real self, the
subjective energy, which must be present, if only in rudimentary
form, in the obscurest sensation, and which is
the power that alone can bring unity and coherence into
the circle of ideas. The association of ideas is one of the most important
real factors in mental life, but can itself be rightly conceived only
by being seen as an elementary form of the energy of the self.
When one, as we shall see Stuart Mill do, concedes
that this energy constitutes the real essence of the mind, and that
it cannot be explained by the doctrine of association, then it is
inconsistent to continue to regard the principles of the latter as
the sole foundation of psychology.

Had Stuart Mill's relation to the earlier English
philosophy been better known, one would not so
often have seen him designated as a positivist and as a disciple
of Auguste Comte. He has himself expressly rejected this
designation in his work: ``Auguste Comte and Positivism.''
Not only was his whole view established before Comte's
writings appeared, but he also departs in very essential
points from the great French philosopher. As the latter's
chief merit Mill emphasizes his conception of the development
of the human mind through three stages, the theological, the
metaphysical and the positive, and the manner in which he
thereby grounds the classification of the sciences. But to
this acknowledgment Mill adds some peculiar limitations.
Because it is admitted that we know only
phenomena and their laws, while the Absolute lies
beyond our reach, it does not follow that the theological
stage should be wholly overcome in the sense
that there should come a time when the human race wholly
ceased to ask after the original cause, the
absolute principle of all. The ``positive'' philosopher may
form on this the opinion that he considers the right one; it
% --- p. 24 ---
\setcounter{footnote}{0}%
\opage{24}will then become the object of positive criticism; but it is
not necessary that all thinkers agree in these
matters. ``He who regards all events as links in
a fixed order, each of them being the unchangeable
consequence of a preceding condition or of a combination
of preceding conditions, adopts completely the positive
way of thinking, whether he acknowledges or
does not acknowledge a universal ground whose original consequence
the whole system of nature is, and whether this universal
ground is conceived as mind or not.'' Because further the
positive stage is to replace the metaphysical, it does
not follow that it would no longer be necessary to make
investigations and analyses of the presuppositions and
principles on which our knowledge rests. But the significance
of such investigations Comte is inclined to overlook,
just as he has given logic and psychology no
place at all in his system (cf. above pp. 9-10). The general
theory of science has two tasks to solve: it is to
show us the ways by which human thought attains
knowledge, and the means by which the truth of knowledge
can be established. But Comte has treated only the
first task, not the second; he has not seen the significance of
verification, which inductive logic shows.
With this it is connected that Comte regarded hypotheses
as means to satisfy our intellectual
need of clarity and coherence among the phenomena,
instead of as ways to find the real causes
of the phenomena. He even declares that ``the most philosophical
point of view in the end leads to conceiving the study of the laws of nature
as destined to represent the external world to us, with
the satisfaction of the essential inclinations of our understanding'',
among which he counts the need of order and harmony, of
beauty and ideality. The center of gravity thus threatens here
to shift from the objective to the subjective, although
positivism, according to Comte's doctrine, was precisely to exclude everything
subjective. Rightly Mill finds in this the germ of Comte's later
mystical speculations, where the ``subjective
% --- p. 25 ---
\setcounter{footnote}{0}%
\opage{25}synthesis'' was to take the place of the objective. Mill is
here more positivist than Comte, but he leaves positivism
by maintaining logic as a distinct science. ---
Just as Comte overlooked the significance of logical analysis,
so he also denied the reality of psychology as an
independent science, because it rests on self-observation,
which according to his view leads only to subjective illusions.
In place of psychology he would put phrenology. Mill
maintains the independence of psychology as a science that
has its own peculiar principles and laws, namely the laws
of the association of ideas, which are for psychology what the law of gravitation is for
astronomy, and what the elementary properties of the tissues are
for physiology. --- Philosophy has thus for Mill a
deeper ground and a wider scope than for Comte. Mill
himself has as a thinker a wider horizon and a greater mobility
than Comte; but on the other hand he stands behind the French philosopher in strict consistency
and systematic completeness. ---

Stuart Mill's chief work is his Logic. He has here
with mastery portrayed the method of the inductive investigator
in its finest nuances and has, more than any other logician,
come face to face with the real, practical
functions of knowledge. It will be difficult to find
a counterpart to this work, unless one goes back
to Aristotle. What he did for the syllogism,
for deductive logic, Mill has done for induction,
for the logic of experimental inquiry. Already Bacon and
Descartes polemicized against the syllogism and denied its
significance as a method of knowledge; but the new logic
that Bacon promised could not yet be presented at that time. Just as
Aristotle could first set up the theory of the syllogism
after reflection and dialectic had been practiced and developed
by Socrates and Plato, so the theory of induction
could be carried through only after experimental inquiry
had labored long and in different fields.
And then it also proved that formal (syllogistic)
logic was not enough, but that a doctrine of
% --- p. 26 ---
\setcounter{footnote}{0}%
\opage{26}objective truth and certainty is needed, a ``philosophy of evidence'',
which is not merely analytic and secures against self-contradiction,
but really carries knowledge further, enriches it with new content
by the way of synthesis.

By inferring one understands in general to derive
one proposition from another. But this is, according to Mill, not
the original form of inference. Originally one infers
neither from the universal to the particular nor from the
particular to the universal, but from one particular to another
particular. Sense-perception shows us two phenomena (a and b)
connected either by simultaneity or immediate succession,
and according to the law of association the perception of
the one phenomenon (a) afterwards evokes the idea of the other (b).
This is the most elementary inference, which is purely instinctive
and mechanical, and which is made even by animals
and small children. When the child has once burned itself (b)
on the hot stove (a), the sight of the stove
(a) will afterwards evoke the idea of the pain of burning
itself (b), and therefore ``a burnt child dreads the fire.'' --- But when
the experiences from which the inference is to be drawn are manifold
and separated in time, they must be held fast and brought together
under a common name; only thereby does it become possible to
give oneself a sharp and conscious account of the operation one
performs. In place of the elementary inference, which went
from the one particular to the other particular, there now enters
a more complex process. First there is set up by induction
(in the narrower sense) a general proposition as the expression
of the sum of the given experiences (ai + a2 -f- a3 = A), and from
this general proposition (A) we then derive the other
particular (b) by syllogistic deduction. But this complex
process is only another form of that elementary
inference: partly the universal is only the sum of the particulars,
and partly we have just as much right to draw an inference
about a whole class as about a single phenomenon, in so far as
it possesses the attributes of the class. The syllogism, the inference
from the universal to the particular, is only a test
we make, the universal itself is only the sum of the 
% --- p. 27 ---
\setcounter{footnote}{0}%
\opage{27}earlier experience, a kind of register, a brief formula, which cannot
itself be the original starting point.

We see thus association psychology lying at the foundation
of Mill's logical theory. But how do we then get beyond
the accidental in this foundation? How can real insight and knowledge
develop out of the blind,
instinctive succession of ideas? Stuart Mill
says himself that what men can conceive, and what they
cannot conceive, is to a great extent accidental and depends
on their experience and habits. By cultivating in particular
certain associations of ideas men can make themselves wholly
unfit to conceive many things, while on
the other hand they can conceive much that again seems
inconceivable to others. How can we then be sure
that we cultivate the \emph{right} associations of ideas? This logic
must necessarily teach us, since otherwise it would not be
what it is to be according to Mill's definition, ``thi science of
the conditions of correct thinking .'' But it is easily seen that
Mill can set up a concept of the right only by an inconsistency,
whereby this becomes different from what is actually
accepted. The associations of ideas that have actually arisen could
not have formed themselves otherwise, and just as little will one,
according to Mill's presuppositions, really be able to alter and
``cultivate'' them, for whence is this formation to proceed,
when in consciousness there are no other factors than the
individual ideas, and these work only according to the involuntary
laws of association?

According to the common conception, what matters is to hold fast
the connections between phenomena that have
the character of necessity, or in other words, to order
the phenomena according to the principle of cause. But Mill is not
satisfied with this answer. By necessity, he says,
one commonly conceives a mystical bond that
unites things, by cause something that not merely \emph{always}
\emph{precedes} something else, but that directly \emph{produces}
something else. Neither of these concepts can,
according to Mill, be used by inductive logic. It finds only a
% --- p. 28 ---
\setcounter{footnote}{0}%
\opage{28}train of phenomena, of which some constantly occur
in union with one another, --- the only connection
it knows. Of necessity there can here be talk only
in so far as by it one understands \emph{unconditional} connection, in that
the one phenomenon is never present without the other
also occurring. Of cause there can here be talk only
in so far as by it one understands the total sum of the positive
and negative antecedents, upon whose occurrence a certain
phenomenon invariably also occurs. Cause is
therefore always phenomenal, and the question is how we
attain the knowledge of the phenomenal causes. --- One
easily sees that Mill, in his zeal to put aside the inner necessity
grounded in itself, does not notice that he merely
transfers the difficulties to another place, namely \emph{between}
the phenomena instead of in their inner nature. An unconditional
connection the empirical-inductive way of looking at things cannot attain, nor
in any case can our knowledge resting on associations of ideas.
When Mill nevertheless posits such an Unconditioned,
it is not merely a logical usurpation; it is also because
there operates here, as will appear later, unconsciously
a principle that leads beyond his conscious presuppositions.
---

It is above all in specifying the ways by which
the phenomenal causes can be known that Mill has
carried the logic of induction further. Bacon had not got beyond
accidental induction, which consists in collecting and
summing up the cases that speak for, and those that speak against
(the positive and the negative instances); the theory of
the experiment was reserved for Mill to set up. And he has
done it with such clarity and completeness that a
student of nature like Liebig has felt himself aided in his investigations
by the study of his Logic.

There are, according to Mill, four methods for empirical
inquiry. They are partly methods of discovery, partly methods of proof.
We give here a quite brief characterization of
them. 1) The method of agreement: When two or more
cases of the natural phenomenon that is to be investigated have only
% --- p. 29 ---
\setcounter{footnote}{0}%
\opage{29}one single circumstance in common, then only this
circumstance can be the cause sought (or effect).
2) The method of difference: When a case in which
the natural phenomenon occurs, and a case in which it
does not occur, have all circumstances in common save
one, which occurs only in the first case, then
this is the cause sought (or effect). This method
leads us by direct experiment with certainty to the cause.
Where experiment is not possible, we must content ourselves
with the method of agreement, which does not lead us to
real causal laws, but for the time being only to phenomenal
uniformities. 3) The method of residues: We subtract from the phenomenon
the part that by earlier inductions is known as
the effect of certain conditions; the remainder of the phenomenon must
then be the effect of the remainder of the conditions present.
4) The method of concomitant variations:
When a natural phenomenon is changed by a change
in another natural phenomenon, it must stand in causal connection
with the latter.

It is not all discoveries that we can immediately
make by means of these methods. They all presuppose experiments
or at least observations. But natural phenomena
are often too complicated for one to proceed directly
to investigate them. The operative conditions may be
so mingled with one another in the real world
that an inquiry advancing by means of the experimental methods
would never be able to be done with them.
Here the deductive method shows itself as a necessary
supplement to the inductive. Deduction presupposes, however,
always a preceding induction. The task of deduction
is to find the law of an effect from the different
factors of which it is the product. But then each of
these factors must first have been investigated separately by observation
or experiment. From the laws of the effect of the individual factors
are then derived the laws of the effect of their combination.
This deduction is a syllogism, which can also
have mathematical axioms for its premises,
% --- p. 30 ---
\setcounter{footnote}{0}%
\opage{30}so that the effect of the combination is found by direct
calculation. But the result of deduction must, in order to be
recognized as scientific truth, be verified i.e. it
must by careful comparison prove to agree
with facts that are found by direct observations. Thus
Newton's theory of gravitation had to be verified by its leading
by deduction to Kepler's laws.

To explain a natural phenomenon or a connection
between natural phenomena means to trace it back to
some familiar phenomenon. But the explanation
really only puts a new riddle in the place of the earlier one;
beyond the phenomena we can never get.
The essence of nature remains equally mysterious,
however far we advance. The highest task knowledge
can here set itself is, according to Mill, to reach the
fewest assumptions from which the order and coherence of nature
can be explained. Natural explanation finds in the answer
to this question its definite limit. On the one
hand the ultimate laws of nature cannot be less numerous
than the qualitative differences between our sensations
and mental states. If we have got so far that we have
shown a peculiar kind of cause for each of the irresolvable
states of consciousness, then we cannot get
further. But besides this subjective limit there is on
the other hand also an objective limit. There are coordinate
causes in nature that we cannot trace back
to any common law. We know no law for the distribution
of the original natural forces and substances
or for the determinate relation that actually obtains
between them. The elementary forces stand in different
relations to one another, without our being able to discover any
rule therein. We know the laws of the operation of the factors, but
not of their original union. Thus there is in a kaleidoscope
always a regular order among the pieces of glass,
and yet the arrangement each time rests only on an accidental
shaking.

According to the common conception there is at least
% --- p. 31 ---
\setcounter{footnote}{0}%
\opage{31}one law that holds beyond all experience, namely that
everything has its cause, the law of causality. Induction itself,
in the sense of an inference from individual cases
to a general law, rests on this conception in so far
also on the principle of causality, as that which gives
thought the right to infer from the particular to the universal,
can only be the presupposition that everywhere in nature
there is a determinate, lawlike connection. Mill now conceives,
as we have seen, induction essentially as an inference
from one particular to another particular, which originally takes place
as a wholly mechanical and instinctive association of ideas.
Experience is thus nothing other than the sum of these associations of ideas
with exclusion of the accidental and transient. He
cannot therefore recognize the law of causality as an originally
operative principle, but regards it as an abstract
generalization of all the experience we have (so far) had.
An Absolute we could never reach by this way. We find
uniformity and lawlike connection in all the individual
circles of phenomena that we investigate, and draw from this the
conclusion that in \emph{the whole} of nature there must be found such
lawfulness. Properly we have a right to say only that
it is found in the part of nature known to us. ``I am of the
conviction,'' says Mill, ``that anyone who is accustomed to abstraction
and analysis, and who sincerely uses his faculties, after
having sufficiently exercised his imagination, will find
no difficulty in picturing to himself that, e.g. in one
of the many firmaments into which astronomy now divides
the universe, phenomena may follow one another in
an accidental way and without any determinate law; nor is there
anything in our experience or mental organization that
can afford any ground whatever for believing that this is not
the case.'' It is a petitio principii, an assuming of
what one was precisely to prove, when one grounds the possibility
of experience on the law of causality instead of grounding
the law of causality on the reality of experience. To draw
the inference ``Fire burns'' one needs no presupposition
that the whole of nature follows uniform laws; it
% --- p. 32 ---
\setcounter{footnote}{0}%
\opage{32}is enough that there prevails uniformity in the effects of fire
--- but this uniformity we do not presuppose;
we learn it precisely from experience. And as it goes with
each single circle of experiences, so it goes with experience
as a whole.

The induction on which the law of causality is grounded
is not experimental, but consists merely in what Bacon calls
inductio per enumerationen simplicem, a direct sum of
the individual cases in which the rule has held good, without
deduction and verification. If it were a special fact
that was at issue, this would be a highly insufficient
ground to build on; but the uncertainty in the mode of proof is diminished
by the fact that we have here to do with the most abstract
generalization, which embraces the whole domain
of our experience. What holds of the law of causality holds
also of the mathematical axioms; they too are
only inductions from experience. But they concern relations so simple
and so frequently occurring that one does not even
notice that one is performing the induction. Nevertheless
geometry would have been a physical science, if its
object could not equally well have been considered in an
image or ideal construction as in reality itself,
and if it had not had its demonstrative character,
a character which, moreover, other branches of natural science
gradually also assume. There is thus in mathematics
no higher or stricter necessity than elsewhere
within our experience. A=C, if A=B and B=C; ---
two straight lines which have once intersected
will always diverge; --- such propositions are inductive
truths, whose validity, like that of the law of causality, rests only
on this, that they have always been found correct and not a
single time false. An absolute impossibility we do not know;
the ground why we cannot suppose that two straight lines
which intersect can converge is therefore only that we
have never seen it, A more perfect proof of general
truths we do not have.

Concerning these main points, concerning the a priori and
% --- p. 33 ---
\setcounter{footnote}{0}%
\opage{33}a posteriori validity of the principle of causality and of
the mathematical axioms, Mill has carried on a long and interesting dispute with
the man whose work on the history of the inductive sciences
according to Mill's own declaration to a great extent made
it possible for him to complete his Logic, William Whewell.
According to Whewell, who in essentials adheres to Kant's doctrine,
every knowledge presupposes not merely that the phenomena
present themselves in a certain way, but also that the knower
brings along the right idea, the right point of view for
their conception. All experience presupposes therefore certain ideas
or fundamental concepts. But, says Mill, these ideas are themselves
only results acquired by abstraction from experience.
When one wishes to express an observation in words, one must
refer it to previously acquired concepts. If there is now
no concept under which it can be brought, then one must
form a wholly new concept as the expression of the new facts.
It was an accident that Kepler beforehand
had the idea of an ellipse when he studied the orbit of Mars;
therefore he could describe and designate it as an ellipse;
otherwise he would have had to form this idea by comparison
of the phenomena, in order then to apply it to their
interpretation.

The logical axioms must share the fate of the
mathematical ones; they are just as little as these original
principles, but the highest generalizations of our experience.
There can thus belong to them no necessity that
goes beyond what experience teaches. A necessary proposition
would have to be one such that we cannot do otherwise
than think it, whose opposite we cannot grasp. But
what does it mean that something is inconceivable to us? It means only
that in advance indissoluble associations of ideas have formed
in our consciousness, which do not permit us to think it.
Every apparently universal and absolute proposition is grounded
on such a fixed association. Only by being attentive
to the significance of these associations of ideas can one
avoid being entangled in subjective prejudices, of which the religion
and speculation of all times afford examples. Therefore it was
% --- p. 34 ---
\setcounter{footnote}{0}%
\opage{34}a benefaction to humanity that James Mill in
his ``Analysis'' carried through the irresistible association of ideas
as the principle for the explanation of all mental phenomena.
There are facts which in all our experience show themselves
connected, and which therefore in our consciousness must also bring about
a firm assumption unshaken by any doubt.
But on the other hand the history of science contains
examples enough that what at one time was regarded
as impossible has later proved really to be true.
One assumed e.g. in the Middle Ages that it was impossible that
there could be antipodes, and that the stars, in order not to
fall down, must be fastened to fixed spheres. When
we continually see that ``inconceivable'' things in the course of time become conceivable,
then the inconceivability of the opposite of a
proposition cannot be a criterion of the truth of this
proposition.

Whewell remarks in reply that when one age rejects
as unthinkable what another age regards as true, or
conversely, this comes from the fact that consciousness continually develops
from the more imperfect to the more perfect, a
development in which new ideas and
points of view continually arise; but this does not shake the truth that the
last, most abstract criterion we have is the impossibility
of the opposite. Lewes and Spencer also appear here against
Mill, although they are otherwise in essentials in agreement with him.
Spencer emphasizes in particular (see above p. 14) that without
such a criterion our knowledge could not
take a single valid step. --- Mill's last word is
that even if certain relations could be pointed out that in and
for themselves are unthinkable to us by our nature, this
unthinkable could nevertheless exist. That something is unthinkable
can have only a purely subjective significance, unless
one, like Schelling and Hegel, would assume that the possibilities of thought and
of existence coincide, that there is an original
correspondence between the microcosm and the macrocosm;
but that is an assumption which could be grounded only
by a supernatural revelation. So long, says
% --- p. 35 ---
\setcounter{footnote}{0}%
\opage{35}Mill finally, as thinkers still so eminent as Spencer
and Lewes can regard a speculative necessity, a
necessity for thought, as proof of the existence of a
real necessity in the universe, so long must it be said that
the old metaphysics still reigns in philosophy. Have
we the authority to say that there is something God or nature
cannot do? Do we then know all the combinations of nature?
There is no necessity outside the operation of nature
itself, and this operation we learn to know only through
experience.

We do not know what is in itself; therefore all
ontology is impossible. Speculative philosophy seeks in vain
to establish laws for being in itself. Metaphysics is
not philosophia prima but philosophia ultima, the final
result of all our experience. Our general ideas contain
nothing that has not been put into them by our experiences
or our habits of thought. There are no abstract ideas.
In the case of compound ideas, which correspond to existing, individual
objects, we are able to turn
our attention exclusively to single ones of the characteristic marks. It
comes thereby to look as if we could conceive them
separated from the rest, especially since by signs and words we can
hold fast what is similar and oftener recurring. When we e.g.
see a white horse, we are able to fix our attention
merely on the characteristic mark ``white'', which is then found again
in other objects; and since we have formed a
special word for it, we say that it is ``the same quality''
in all cases. But it is nevertheless always the ideas
of concrete phenomena that make up the content of concepts,
and we do not think by means of concepts, but by means
of concrete ideas and of the names by which we hold fast
these ideas. --- Mill thus approaches here
the extreme nominalism, for which concepts become one
with mere names. ---

After one has become acquainted with Mill's Logic,
one cannot without astonishment hear him declare that
the results of this work stand in no necessary
% --- p. 36 ---
\setcounter{footnote}{0}%
\opage{36}connection with any of the peculiar doctrines of the two
chief philosophical tendencies. For quite certainly this holds
for such parts as the exceedingly important one
that presents ``the methods of experimental inquiry'';
but the whole light in which the logical processes are placed, the whole
explanation that is given of their nature and ground, confirms
nevertheless in full measure what we emphasized from the beginning,
that it is a quite determinate doctrine inherited from the earlier empirical
school, association psychology, that lies
at the foundation from first to last. Association psychology points
back in turn to Locke's doctrine of the mind as a blank
tablet; only when the mind is conceived as perfectly passive
can its essence be exhaustively explained by the laws
of the association of ideas. In other places Mill also counts himself in the philosophical
school whose chief representative is Locke, while
the opposing one is represented by Kant. But he denies
that the chief difference between the two schools consists in this, that
the one regards the mind as purely passive in knowledge,
while the other ascribes to it an original activity.
The empirical school too acknowledges, according to him, an original
mental element in our ideas and does not regard
concepts as mere copies of the impressions but as
products of mental labor. But in the same essay
in which this statement is found (Dissertations and discussions
III.), Mill points out precisely as Bain's merit that he has
led association psychology beyond passive determinism
to the acknowledgment of an original activity in the mind.
Herein lies the concession that the older empirical school
did \emph{not} teach the activity of the mind. And it is a great question
whether one of the essential peculiarities of
this school is not thereby given up. It is in any case
impossible both to hold fast that original activity in the essence
of the self and to regard the laws of the association of ideas as the only ones
that hold for the life of the mind. The original energy must
surely stamp every impression in a peculiar way
and must likewise be determinative for the connections in
which the ideas enter with one another. 
% --- p. 37 ---
\setcounter{footnote}{0}%
\opage{37}Recognition of an original activity leads us properly from
Locke to Kant.

Stuart Mill places the chief difference between the two
philosophical schools in this, that the one (the a priori, the
intuitive or the introspective school) starts from some
of the most compound phenomena of mind as original
and believes it has in them fundamental presuppositions about the
nature and essence of things, while the other (the experimental or
the a posteriori) seeks to reduce the apparently original
principles by continued analysis. It is clear that
hereby no difference of principle is stated; for
both schools would surely agree that one should
carry the analysis as far as one can in general carry
it. The analyses we find in Spinoza and Leibnitz, in
Kant and Hegel can also in clarity and consistency
measure up to those we meet with in the thinkers of the other school.
And Mill concedes himself that analysis \emph{can} be carried only
to a certain point. We stop in the end at states of consciousness
that cannot be traced back to simpler and more
elementary ones. In particular the sensation of likeness and
unlikeness lies at the foundation of all other sensations. Complicated
relations of likeness and unlikeness we can trace back to simple
and uncompounded ones; but the very faculty of apprehending relations of likeness and
unlikeness must constantly be presupposed. Here too
we come upon an element that points beyond the empirical,
a contribution to knowledge from the side of the mind. The conception of
likeness rests on a comparing, a mental activity.
We are thus referred back to the activity of the mind.
Kant saw this as consisting in a combining of the manifold given to
sensation into the unity of intuition and of thought.
And the comparing activity points in the
same direction; comparison is the condition of combination.
All association of ideas presupposes, as William
Hamilton has shown\footnote[1]{Mill acknowledges this. See his note to James Mill's ``Analysis'' I p. 112 f.}, this comparing; for in order that an
% --- p. 38 ---
\setcounter{footnote}{0}%
\opage{38}idea may be able to evoke another, with which it
is wont to go, it must first be recognized, be seen as one that we
have had before; and the recognition takes place by comparing.
From this it follows again that the association of ideas itself becomes only an
elementary form of the activity of the self. But then the
principle Mill has laid at the foundation of his logical theory is only
a derived principle, and the final explanation is not given;
that would be the case only if the logical laws were traced
back to that original activity, in which they must have
their germ. It is Mill himself who has not carried through
the analysis. And when the matter stands thus, we have
not merely come beyond empiricism, but also beyond
skepticism. Our knowledge is then seen to rest on presuppositions
that are given in our mental organization, in
the essence of our consciousness, and that it is impossible for us to abolish.
The very act of thinking or imagining something always
carries with it that the content of thought and of idea is ordered
and connected in a certain way; hence the presupposition
that there must everywhere be lawlike connection, a presupposition
that can be given up only when thinking itself
is given up. The necessity we recognize may enter as
a link in a greater connection that we cannot survey;
but thereby it is not abolished.\footnote[1]{When Stuart Mill thinks that one could well imagine a world where two and two were not four but five, then one has rightly answered to this that in that case it would not be 2 + 2 that was = 5, but 2 + 2 + 1, the 1 being due to the circumstance by which the constitution of that world differs from ours. See Mansel: The philosophy of the conditioned p. 131.} This presupposition of
a lawlike connection is therefore found in all men and in
all times, even if only through the more recent science of nature
it has received its full confirmation and secure execution. Also
in the mythological ideas, even in the wildest superstition
there lies at the ground the conviction of order and connection.

By going back to the presuppositions we have thus
\opage{39}found Mill's fundamental error. In this way criticism
acquires greater significance than when one, as is for the
most part the case, argues from the skeptical consequences
that may be drawn from Mill's doctrine, consequences
that Mill's supple dialectic could moreover make
harmless. By directing the criticism against the presuppositions
we gain at the same time the advantage that it becomes possible for us to maintain
the great significance of Mill's Logic (cf. above pp. 25 and
28). It does not do justice as an exhaustive presentation of
the process of knowledge; but when the one-sidedness and relativity
of the standpoint is seen, one must admire the energy and
the acuteness with which the tasks are treated. A
presentation that anxiously strove for correctness
and completeness with regard to first principles
would perhaps hardly be able to penetrate so deeply into the
essence and forms of concrete knowledge as Mill from his
standpoint has been able to. ---

What brings Mill to leave the original presuppositions of
association psychology is that he has got his eye on the
concept of selfhood, of subjectivity. As this
concept presses forward, it brings unrest and contradiction
into the empirical presuppositions. In his specifically psychological
investigations Mill comes still nearer to this concept
than in his Logic. To these we now turn, in that
we first set forth his conception of psychology in general.\footnote[1]{In the Logic the method of psychology is dealt with; the other investigations are found especially in ``Examination of Sir William Hamiltons philosophy'' and in the notes to James Mill's ``Analysis''.}


Mill maintains, as already mentioned, the independence of psychology
as a science. It is not merely a part of
physiology. That every mental state has a certain nervous state
as its immediate condition and nearest cause
cannot be said to be demonstrated in the same strict sense
in which it is demonstrated
% --- p. 40 ---
\setcounter{footnote}{0}%
\opage{40}as regards the sensations of sense. In any case we do not know the essence of these
nervous states, nor their relation to the mental
states; we are therefore restricted to studying these
latter directly. Psychology has for its object the original
or derived laws according to which one mental state
follows upon another. To be sure, there can here in particular cases
be no question of exact insight; we can find the fundamental laws
and the fundamental conditions, but these are modified in the
individual cases by the individual circumstances. Mill compares
psychological knowledge with the doctrine of ebb and flood;
here too the main causes (the attraction of the sun and the moon)
stand fast and can be explained and predicted; but there
is a multitude of secondary factors (the direction of the wind, local
circumstances, the character of the bottom of the ocean etc.),
which each depend on particular laws, and which exert
considerable modifications. --- We know already that the laws of the
association of ideas are, according to Mill, the fundamental laws of psychology.
In order to find these, one must, as in natural science,
apply a double method: partly the facts must be
traced back to laws, these in turn to more general
laws, and so on; partly the compound
phenomena must be resolved into single and elementary ones. The first
method has its analogy in Newton's discovery of the law of gravitation
as explanation both of the motion of the heavenly bodies and
of the motion of falling on the earth; the second has its model
in chemical analysis, by which substances are resolved into
their constituents. But when one has thus traced a
mental phenomenon back to simpler and more elementary
phenomena, by whose combination and association it has
arisen, then it is not, as Hartley and James Mill supposed,
enough to study these elementary phenomena; one cannot
deduce from them the nature of the higher, more compound phenomenon,
any more than one can deduce the properties
of water from those of oxygen and hydrogen. --- Proper psychology
is an empirical science. But from the general
laws to which it leads one can find by way of deduction
how the human character will form itself under
% --- p. 41 ---
\setcounter{footnote}{0}%
\opage{41}certain determinate physical and moral conditions. Psychology
then passes over into ethology, the doctrine of human
character. There can of course be no question of predicting
determinate facts; but from the rightly known psychological
laws one can derive what inclinations and
strivings will appear under certain circumstances. In order to
acquire validity, the results of this deduction must be verified,
shown to agree with actual observation.
--- On psychology and ethology sociology finally rests,
the doctrine of human society. This too
is a deductive science. After an investigation
of each in particular of the factors that intervene in social life
and there exert their influence through the general
psychological and ethological laws, the united
operation of all these factors is to be determined. The combination can here often
become so complicated that the deduction cannot be carried through,
just as physics cannot solve the problem of the mutual
action of three bodies gravitating toward one another.
Sociology must therefore sometimes be inversely
deductive i.e. instead of beginning by deducing from
the psychological and ethological laws, it must begin with
deriving an empirical law from historical and social facts,
in order thereupon to find a deeper grounding of this empirical
law in human nature. The idea of the progress
and development of humanity is an example of such an
empirical law, which in order to become a rational law, a
causal law, must be seen in its origin in human
nature itself.

One might miss ethics in this survey of
the sciences of mind, and it has sometimes been believed that Mill
recognized no other ethics than ethology, the doctrine of
the development of characters. But this is a misunderstanding.
According to Mill ethics is essentially an art, a part of the
art that corresponds to the science of human
nature and human society. The peculiarity
of art in contrast to science is that it speaks
in the form of rules and precepts, not in the form of statements
% --- p. 42 ---
\setcounter{footnote}{0}%
\opage{42}of facts. Art sets itself an aim as desirable,
defines it and hands it over to science, in order that this
may investigate its causes and conditions. The doctrine of
the ends and the general principles that can be
derived from them constitutes teleology or what the German
philosophers call the principles of practical reason,
an art of living, which consists of three parts: ethics, politics
and aesthetics, --- the doctrine of the right, of the expedient
and of the noble and beautiful. By means of this
art a guiding rule is won for the judgment of all aims
and objects of our drives.

The difference Mill thus sets between ethics and
ethology necessarily presupposes a factor in the mind,
whereby it can raise itself above the psychological
mechanism that association psychology teaches. To understand
the transition from psychology to ethics we must consider the concepts of the self
and of the will.

It is a presupposition Mill has in common with positivism,
that we know only phenomena; we know nothing
in and for itself, the mind in itself as little as matter.
The Scottish school started from common sense
and its instinctive assumption of an external
object as the cause of our sensations. But ``common
sense'' is something far too combined
for association psychology to be able to recognize it as something
original that cannot be subjected to analysis. In the assumption
of external objects independent of us there lies really
nothing other than the assumption that, besides our actual,
present sensations of sense, there is given a real possibility
of many other sensations, which could occur
under certain conditions. Things are nothing other than permanent
possibilities of sense-impressions. For us they stand as
more than possibilities, because we involuntarily apply the concept of substance
and of cause, which has validity only within the
series of phenomena, to the very possibility and ground
of this series. And the idea of this real substance
gains in strength when we notice that not only we ourselves, but
% --- p. 43 ---
\setcounter{footnote}{0}%
\opage{43}others too proceed from the presupposition of its existence.
--- Stuart Mill thus adheres (like James Mill) to
Berkeley's idealism and regards it as an illusion when
one thinks that a material substance is necessary to
bear the phenomena, just as the Indians thought that the earth
needed a point of support and therefore rested on an elephant.
As he follows Berkeley with regard to the concept of matter,
so he seems to follow Hume with regard to the concept
of mind.
The mind is only the train of our actual and
possible sensations. It differs from the train
that constitutes matter in that it also includes the
ideas that are not confirmed by experience, and
all our thoughts, feelings and resolutions, but especially
in that this train is restricted to ourselves, each for
himself, and is not, like that, common to all men under the same
circumstances.

This renewal of the phenomenalism of Berkeley and Hume
is a necessary consequence of association psychology.
As this doctrine in the last century called forth
Reid's appeal to common sense and to
common morality, so it has now also aroused fear
of absolute skepticism, and Mill has thought it
necessary to demonstrate that the proofs of the existence of God
and the immortality of the soul are ``just as good'' according to his
doctrine as under the usual presuppositions. For us
it is of the greatest interest* to see how Mill here too quite
definitely goes beyond his own presuppositions and
demonstrates their insufficiency. Immediately after (in the work on
Hamilton) having given the stated explanation of the concept of
the self he declares, namely, that this explanation does not
suffice, that it in particular does not comprehend phenomena such as
memory and expectation. ``When we therefore,'' says he,
``speak of the mind as a train of sensations, we must
add, a train of sensations that is conscious of itself
as past and future. And we stand then at the
dilemma either to assume that the self or the mind itself is
something different from the train of actual and possible sensations
% --- p. 44 ---
\setcounter{footnote}{0}%
\opage{44}or else to assume the paradox that a train
of sensations can be conscious of itself as a train.''
The self and memory, he says in his notes to James
Mill's ``Analysis'' (where one will on the whole find the clearest
and sharpest expression of his psychological view), are
only two sides of the same fact: without memory
no self, and conversely: without self no memory. With
right Mill calls it a paradox that a ``train'' should
be able to be conscious of itself; but when this is so, what
becomes then of the definition of the mind that Mill and
association psychology assume, that it is a train of
sensations? Is it then not necessary to seize the
first alternative and assume that the mind or the self does not
coincide absolutely with the train of sensations? This
Mill also seems to do when he declares that the bond
that holds the train together is just as real as the
individual sensations that form the train, and is not a
mere product of thought, and that we have here an original element
that has nothing in common with any other thing
whatever and therefore must be designated by its own name. When he
also declares that we stand here before an insoluble riddle,
then one must nevertheless remark that he himself in an artificial way
has to a great extent produced this riddlesomeness by fixing
a definition of the essence of the mind, of which he immediately
afterwards must concede that it does not hit the real
kernel. --- It is not easy to understand how Mill has
been able to be accused of materialism; for it is clear that
in so far as he lets himself be driven from his empirical-skeptical
standpoint, he must precisely go in an idealistic direction. Proceeding
from empirical psychology he shows us the necessity
of a psychology that, as it has been since Kant the task of
German philosophy, develops the central concepts of mental life,
especially the concept of self-consciousness. And
an inner connection between the two lines of investigation will
become possible, if it can be shown that the whole
system of mechanical processes that association psychology
% --- p. 45 ---
\setcounter{footnote}{0}%
\opage{45}presents is a serving foundation and means for the development of
self-consciousness. ---

In close connection with the concept of the self stands that of the will.
With regard to the conception of this concept Mill has a
support in his logical theory, which he applies here in an
interesting way. He denies, as we have seen, necessity
and puts in its place fixed and uniform sequence.
It seems to him that by necessity one commonly
understands a mystical power that from without overwhelms
the individual will. If necessity is conceived
in this way, the doctrine of necessity (determinism)
becomes fatalism, as among the Mohammedans, according to whose belief
the predestined will happen whatever one does to prevent
it, and as in Owen, according to whose doctrine man's character
is formed for him and not by him. Mill therefore wishes
preferably to abolish the expression necessity. Every action
presupposes as its condition earlier states of consciousness
(thoughts, feelings and drives), without which it would not
have come to be. Effect follows here upon cause just as
regularly as in the physical domain; but of inner,
mystical necessity there is here just as little talk as
there. We can ourselves develop our will by applying
the laws of the association of ideas to procure for the nobler motives growing
power. Originally the action is determined by pleasure and
pain; but by the power of the association of ideas there becomes gradually
the means an object of our desire without regard to the original
end: the action is then willed in itself, quite
apart from our finding pleasure in what it procures for us. The difference
between a bad and a good man does not consist in this,
that the latter acts against his strongest inclinations,
but in this, that his striving to act rightly and his
abhorrence of the wrong are strong enough to keep
every other drive in check and bring it to silence.
Since now such a mental state is possible, men are
fitted for self-education and moral guidance. Self-education
is everyone's moral duty. A necessary
means to it is that thought forms for itself a clear standard
% --- p. 46 ---
\setcounter{footnote}{0}%
\opage{46}of right and wrong, in order that moral drive and aversion may work
in the right way. --- Mill's doctrine shows clearly enough also
on this point beyond his original presuppositions.
For how can the concept of self-education acquire reality, if
the self is not conceived as an original energy that is able
to penetrate and develop the given conditions (impressions
and ideas) and guide the association of ideas in a certain determinate
direction? ---

The question of the freedom of the will was one of the
points in which Mill differed from the acute
Scottish philosopher, against whom he has directed one of his
most important works, William Hamilton, in whose doctrine the free
will in the sense of absolute freedom of choice plays a chief role,
in that it, and the moral responsibility conditioned by it,
was in his conviction the only foundation
on which the belief in the existence of God could be built. But
there are also other points of Mill's polemic that we here
must bring forward for the characterization of his standpoint. Hamilton
maintains energetically the relativity of knowledge; the Absolute
and the Infinite we cannot think, because our
thought cannot conceive that which is complete in itself, and
which has no limits. Although Mill naturally must be
at one with Hamilton in the maintaining of the relativity of knowledge,
yet he polemicizes against this motivation.
It is, he thinks, pure abstractions Hamilton fights
against. \emph{The} Absolute and \emph{the} Infinite is in itself meaningless.
If the Absolute really existed, it would have to
possess all possible predicates in the fullest extent:
it would have to be absolutely good and absolutely evil, absolutely wise
and absolutely stupid, absolutely powerful and absolutely weak, in short,
it would be self-contradictory. When Hegel nevertheless
has set up such a concept, he has thereby only
earned the merit of having dissolved speculative
metaphysics by a series of reductiones in absurdissimum.
--- It is clear that Mill has not at all understood Hegel's
thought or at all had an eye for the problem that lies
in the concept of the Absolute. His strength lies, on the
% --- p. 47 ---
\setcounter{footnote}{0}%
\opage{47}whole, in the treatment of individual, concrete questions; he
does not gather the individual philosophemes into a total view.\footnote[1]{``It is the full, the whole vision that he lacks; the great capacity of his mind unfolds itself upon objects that after his treatment of them continue to be just as isolated as they were beforehand; they are not brought to unity and harmony with one another; the philosopher's hand is still always needed with regard to them.'' Quarterly Review. July 1872. p. 81.}
He can without feeling any intellectual want let knowledge
end in an accidental plurality, and in so far as he
raises the question of the reality of the idea of God, this
question does not stand in connection with his proper philosophizing.
The intellectual drive that does not rest until it
can get the scattered threads connected, or at least
see them converge, this drive is foreign to him.
His empiricist and prosaic way of thinking, which, as one
has reproached him,\footnote[2]{Ibid. p. 117.} does not feel ``the mystery of life'', regards
it as possible, even if he himself does not seem to follow this
way out, to settle down to rest in a deism, where the supernatural
is relegated to ``the beginning of the world'', and where the Deity
as with Epicurus lives in indifferent remoteness from
the world. On the contrary it has been the chief problem of philosophy
from its first days to seek an unconditioned principle
as ultimate explanation and reconciliation of the oppositions of existence.


Mill thinks that in place of \emph{the} Absolute, \emph{the}
Infinite, one should put something absolute, \emph{something} infinite (an
absolute and infinite being), the concrete in place of the
abstract, and he does not regard it as impossible to idealize
the highest qualities we know, and thereby reach a relative
knowledge of the nature of God. At least he maintains
this possibility against Hamilton in order to cut off the theological
consequences that the latter drew from the impossibility of
knowing the Unconditioned. One notes that Mill here has
come into a domain where he is not at home.
% --- p. 48 ---
\setcounter{footnote}{0}%
\opage{48}From Hamilton's standpoint he has been aptly answered\footnote[1]{\emph{Henry Mansel}: The philosophy of the conditioned. p. 115.}
that ``something Absolute,'' ``an absolute Being'' is a self-contradictory
concept for us, since every ``something,'' every
``being'' we know is relative and finite, so that the added word ``absolute'' leads us beyond the
limits of our knowledge.

It is, by Mill's own declaration in his Autobiography,
not so much theoretical as practical motives that
led him into conflict with Hamilton, in whom he saw the
most significant representative in England of the philosophical
tendency opposed to his own. ``The difference between the
two philosophical schools, the one that rests on intuition
and the one that rests on experience and association, is
not merely a matter of abstract speculation; it is full of
practical consequences and lies at the root of all the
greatest oppositions in practical views in an age of progress.
The practical reformer must constantly demand
changes in things that are supported by powerful and widely diffused
feelings, or doubt the apparent necessity
and immovability of existing facts;
and it is often an indispensable part of his case to demonstrate
what origin these powerful feelings have, and
how these facts came to seem necessary
and immovable. Therefore there is a natural enmity
between him and a philosophy that sets aside the explanation
of feelings and moral facts by circumstances
and associations of ideas and prefers to treat them
as original elements of human nature, a
philosophy that has an inclination to maintain favorite doctrines
as intuitive truths and to regard intuition
as the voice of Nature and of God, speaking with a higher
authority than our reason.'' One hardly errs when
one finds in this utterance the deepest motive of Stuart
Mill's whole standpoint. But as to the matter itself,
% --- p. 49 ---
\setcounter{footnote}{0}%
\opage{49}it must be remarked that Mill here characterizes the ``intuitive''
school as though the Scottish school were a worthy
scientific representative of it, and that one need only
mention names like Spinoza, Kant and Fichte to demonstrate at once
the unwarrantedness of Mill's assumption that only the
empirical school should represent the philosophy of freedom of mind and
of progress.

But in Hamilton Mill combats not only the ``intuitive''
philosophy, but also theology, supernaturalism.
Hamilton's polemic against the philosophy of the Unconditioned (Schelling
and Hegel) had, namely, as its background a moral-religious
faith that was not supposed to be able to find any expression or
any foundation through knowledge, because its object
is the Incomprehensible, that which is absolutely exalted above everything
human. Against this doctrine Mill expresses himself thus:
``If, instead of the glad tidings that there
exists a Being in whom all the excellences which the highest
human thought can conceive exist in a degree inconceivable to us,
I am told that the world is governed by a Being
whose attributes are infinite, without our being able to understand
what they are, any more than what the principles of his
government are, except that `the highest human morality
which we are capable of conceiving' does not sanction
them, --- convince me of that, and I will bear my
fate as well as I can. But when one tells me
that I am to believe this and at the same time to call this Being
by the names which express and affirm the highest
human morality, then I say in plain words that I
will not. Whatever power such a Being may have
over me, one thing there is that he shall not be able to do: he
shall not compel me to worship him. I will not call
any being good that is not what I mean when I
apply this predicate to my fellow creatures; and if
such a Being can sentence me to Hell because I
would not call him so, then I will go to Hell.''\footnote[1]{Examination of Sir William Hamiltons philosophy p. 103.}

% --- p. 50 ---
\setcounter{footnote}{0}%
\opage{50}It is only through the reading of the Autobiography
that one gets a complete explanation of the energy and
the bitterness with which Mill expresses himself. One sees from it
that he was early led by his father to the conviction that
the belief in a Providence was not only theoretically untenable
but also morally reprehensible, because it necessarily veils
all the evil and all the corruption that is found in
the world, indeed even leads to the worship of a Being who, knowingly and willingly,
has created such a world! --- The
only religious faith that is possible under such presuppositions
must be a practical faith; it must have its
object before it and not behind it, in possibility,
not in reality. Stuart Mill seems, though he nowhere
states himself clearly in his own name on these matters,
to side most nearly with those who, ``though they cannot
believe that the world can have a Creator and Ruler who possesses
absolute power and perfect goodness, yet have that which
constitutes the real value of all possible religions, an ideal
conception of a perfect Being, toward which they constantly
fix their gaze as the guide of their conscience;
and this ideal of the Good is, in general, nearer to perfection
than the objective Deity of those who
consider themselves obliged to find absolute goodness in the Creator
of a world so filled with suffering and so disfigured by
injustice as ours.'' --- It is a question whether Mill
escapes the difficulties by putting a subjective
idea of God in the place of an objective one, and whether on the whole an
ideal practical striving, as he demands it, can
be possible, if the reality from which every striving
must after all proceed, and in which every possibility must
have its point of attachment, is so devoid of ideas that the Manichaean
dualism, according to what James Mill taught, should be
the only morally defensible view. But in any
case Mill here points to an ideal that cannot be explained
by the working of the association of ideas but points back to
deeper sources, sources which Mill in so many hints
has opened for us, but only to stop them up again at once.
% --- p. 51 ---
\setcounter{footnote}{0}%
\opage{51}It is interesting to see how in his practical
conception of the Deity he meets with I. G. Fichte, whose
idea of God was the moral world-order, realized by
free action. Fichte is perhaps the one among the German philosophers
with whom Mill sympathizes most; but Fichte is more consistent
than he, in that he has seen that practical
ideality necessarily presupposes theoretical ideality and cannot
be attained by a purely empirical road. ---

As Mill in general in his philosophical view
builds on the earlier English inquiry, while
nevertheless often departing from it more than he himself notices,
so it is also with regard to ethics. Mill's
ethics is Utilitarianism, ``morality grounded on the principle of utility or
happiness.'' This principle appears in English
philosophy already in Bacon, whom Macaulay, as is well known,
has presented as the prophet of the useful. Richard
Cumberland (in the seventeenth century) seems\footnote[1]{Cf. \emph{Hallam}: Introduction to the litterature of Europe. (7th ed.). IV. p. 160.}
to have to be reckoned as the real founder of utilitarianism;
according to him all moral laws can be reduced to the one:
striving after the good common to all rational beings,
whereby our own good is promoted as a part of the whole,
so that the individual's road to his own welfare is everyone's road to the
common welfare. Yet in Cumberland the individual regard
is entirely subordinate to the regard for the universal, for society.
When Locke's philosophy had broken through, and all knowledge
was conceived as resting on sense-perception, ethics too
was founded on the drive toward felicity, on one's own
interest. Whewell (History of Moral philosophy in England)
remarks that this doctrine agreed with the national
way of thinking, and Tocqueville shows that it is the natural one
for a democratic commercial people like the American.
To this it may be added that it can very well be brought into connection
with theological points of view. So in William Paley,
% --- p. 52 ---
\setcounter{footnote}{0}%
\opage{52}whose ethical work (published 1785) was made the basis of instruction
at the University of Cambridge right into this century. Paley,
who in his youth had won the prize for a treatise
in which he placed the Epicurean morality above the Stoic,
taught that one ought to do what is useful, because it is God's
will, and because it is necessary to salvation. The individual
interest he supported by the hope and the fear
that religion calls forth. But the theory of utility was secularized
and thereby at the same time first properly carried through by the famous jurist
Jeremy Bentham († 1832). Bentham's importance
lies no doubt nearest in the field of practical reform,
and he has a special merit in his polemic against the
many subtleties and the traditional formalism that
made English law a perfect chaos, where the
most shameful abuses could thrive. But his very tendency toward reform
had its ground and its support in his energetic
and methodical thinking, which sought to deduce from the principle of utility
norms for the judgment of existing laws. It
is from Bentham that Mill's ethical and social views have
their origin; in his youth he was absorbed to such a degree in
Benthamism that he embraced it with a sectarian's
zeal. True, he later rectified several points in this
doctrine, but he has not arrived at any clarity about the
difference between his own and Bentham's conception.

There lies in utilitarianism the possibility of a double
point of departure and center of gravity: one can lay the emphasis on
the egoistic interest and show how this, for its own
sake, must enter into a harmonious relation to the interests of
others; and one can start from sympathy and the
individual's relation to the larger whole and show how the
egoistic interest is only a subordinate element within
this. As already indicated, Cumberland goes the
latter of these ways, Locke and Paley the former. The relation
between Mill and Bentham rests likewise on the fact that they each
go in their own direction, conceive the same principle from two different
sides. According to Bentham every human being follows
in each of his actions his own pleasure and interest.
% --- p. 53 ---
\setcounter{footnote}{0}%
\opage{53}Self-sacrifice consists only in giving up a lesser pleasure for
a greater. But the happiness of one cannot be separated from that
of another; everyone's well-understood interest must therefore lead
him to work for the happiness of others. Hereby arises
what Bentham calls the ethical. ``Do not dream,'' says
he (Works I. p. 27), ``that people will move their little
finger to serve you, when their advantage in doing so does not stand
clear before them. But they will wish to serve you, when they
perceive that by doing so they can serve themselves. In these
mutual services consists virtue.''\footnote[1]{One is here reminded of the task that Carlyle set the utilitarians: ``Given a world full of scoundrels, to show how virtue results from their united endeavors.''} Bentham assumes
that sympathy is indeed an original factor in human
nature; but he does not consider it strong enough to bear
the ethical. Stuart Mill differs from Bentham precisely
in the weight he lays on the sympathetic feelings
as a vehicle for the development of the ethical. In Bentham, who
essentially held to the egoistic interest, virtue became
one with prudence and outward rectitude. He also does not emphasize
the wish for self-education, for the ethical cultivation
of one's own personality, but regards man only
as a business being, understands only ``the business part of
human affairs.'' In Mill, on the other hand, the egoistic interest is
not merely subordinate to the relation to others but also to the
ideal need of self-development. Already at the end of the
Logic he expresses himself in a way that seems to point
beyond the ``theory of utility'': ``I fully admit that the cultivation
of an ideal nobility in will and action ought to be an
end for which the particular striving after one's own and
others' happiness (insofar as it is not included in this
idea) ought in every case of collision to give way. But
I believe that precisely the question: wherein consists this loftiness
of character? --- must be decided by going back
to happiness as the standard. For the individual the
character itself ought to be the highest aim, because such an ideal
% --- p. 54 ---
\setcounter{footnote}{0}%
\opage{54}loftiness or an approximation to it will more than anything else
contribute to making human life happy, both in the comparatively modest sense of enjoyment
and freedom from pain, and also in the higher sense
that life does not become what it now almost commonly is,
childish and meaningless, but becomes such as
human beings of more highly developed faculties could
wish to lead.''

It is a false idea that one forms of the morality of utility
when one thinks that it must necessarily be
a low and finite way of looking at things. Bentham and Mill
are enthusiasts, passionately inspired by the progress of
humanity. They consider a condition possible in which
universal happiness prevails, and in this they represent the most complete
opposite to Schopenhauer's pessimism. But
their optimism is not, as optimism so often is, insipid
and superficial; it has for its presupposition a deep knowledge and an inward
confidence in human nature; ---
it is also not conservative, does not hold, like the vulgar
optimism, that all is good as it is, but consists
essentially in the belief that all can become good, however far
it now is from being so, since nineteen-twentieths
of the human race, even in those parts of the world now existing
that are least deeply sunk in barbarism, live in
misfortune. The present wretched education and wretched
social order are the only real obstacles to a
happy life for all humanity, i.e. ``an existence
consisting of few and transient sufferings, many and
various joys, with a definite preponderance of
the active joys over the passive, and having as its foundation
for the whole this: not to expect more from life than
life can offer.'' The higher intellectual and moral enjoyments
will become more and more accessible to all. And
in any case there lies in the very struggle against misfortune
and in the interests and feelings that this struggle calls forth, a
noble satisfaction that stands high above any egoistic
enjoyment whatsoever.

% --- p. 55 ---
\setcounter{footnote}{0}%
\opage{55}What is peculiar to the theory of utility, as
Mill conceives it, is that it presupposes the ethical instead
of grounding and developing it. According to Mill the
point of dispute is not whether one ought to act rightly, but what
the right is; and that, he thinks, only the theory of utility can
teach. But the properly ethical is not a ``what,''
but precisely a ``that;'' the relation of obligation itself is the
point at which all particular tasks, all virtues and duties
must have their point of attachment. Only through its connection
with this central point does something become good in the ethical
sense. The theory of utility, on the other hand, holds to the peripheral,
to the content and effect of actions instead
of to their inner, personal character. Mill even says
that in this he entirely agrees with Bentham, that the motive
has nothing to do with the moral worth of an action,
though it has very much to say with regard to
the worth of the agent. But can one really thus
regard the judgment of actions and of the acting
persons as two different things? The action may be
useful without the person therefore being morally upright. A good and
noble character is not something that can be regarded as a supplement
to the ethical in the action; on the contrary, without an understanding
of the character and the personal qualities no
judgment at all can be passed on the morality of the action; --- therefore
moral judgment is so exceedingly difficult. On the
other hand utility and happiness, ideally conceived, point back
to the personal development of character; indeed, there is,
as we have also heard Mill himself declare, but as he in
his own definitions of concepts for the most part forgets, no
true utility or happiness outside of this. If Mill did not
at bottom conceive the matter in this way, --- and so
here, as in the Logic and the Psychology, went beyond his own
presuppositions, --- there would arise a dangerous enemy
for him in the pessimism of Schopenhauer and Hartmann, which is
the true consequence of the distinction between happiness and
task, between ``what'' and ``that.'' But both pessimism
and utilitarianism separate what ought not to be separated.
% --- p. 56 ---
\setcounter{footnote}{0}%
\opage{56}``The justification of the individual and of the race rests only on their
peculiar end and task, through the solution of which
their own nature is developed, and only in this development itself
can happiness and satisfaction be sought\dots{} The eudaemonological
and the ethical stand in an inner connection, and
the significance of pessimism lies in demonstrating indirectly the necessity
of this unity.''\footnote[1]{``Philosophien i Tydskland efter Hegel.'' p. 285-286.}

The different points of view that thus assert themselves
in Mill must naturally entail obscurity and
vacillation in the definitions of concepts. He declares, for instance, that
``the soul is not in a healthy condition, not in the
condition that agrees with the useful, when it does not
love virtue for its own sake,'' --- but adds in the same
breath that actions and plans are virtuous only because
they promote another end than virtue. To the question:
is virtue a means or an end? Mill has no decisive
answer. At one time it is said that it ought to be desired for
its own sake; at another, that the utilitarian moralists ``place virtue
at the head of the things that are good as a means to the
final end;'' at another, that it is or can become one of
``the constituents of happiness,'' --- ``a part of the end,'' indeed even ``a more important part of it than any of the objects
to which it is a means''!\footnote[2]{\emph{Stuart Mill}: Moral grundet paa Nytte- eller Lykkeprincipet. Translated and introduced by G. Brandes. p. 46-48.}

From the Autobiography one sees that these dialectical
difficulties in Mill's own life led to a crisis, which
he ascribes great importance for his mental development.
As an eager Benthamite, by his own admission even a fanatical one,
he was enthusiastic in his youth for the spread
of utility and happiness in humanity. But --- he came
to ask himself --- would you really feel yourself happy,
if all your ideals were realized? --- and he had to
answer this question in the negative. A long period of melancholy
% --- p. 57 ---
\setcounter{footnote}{0}%
\opage{57}and despair followed upon this No. At last he gained new
courage and new strength through the conviction that happiness, though
it is the aim of life, can yet be attained only by not making
it a direct aim; only those are happy who make something
other than happiness their ideal end; happiness they
then find on the way. ``As soon as you ask yourself whether you
are happy, you will cease to be so.''\footnote[1]{Autobiography p. 142.}

It cannot be demonstrated more clearly that the principle of happiness
points beyond itself, is only a relative principle in the ethical
respect, and that it can therefore not be laid as the foundation
of ethics. Mill has confused the psychological point of departure,
the educative means, with the ethical end.
But his strength lies on the other hand precisely
in showing how this educative means works.
The feelings of pleasure and displeasure are the original motives for
all human activity and action, and Mill shows now, by
a pertinent application of association psychology (an application
in which he has Hartley and James Mill, indeed really
already Plato and Aristotle, as predecessors), how
man by the working of these very motives is led step by step
to set himself higher and nobler ends, to abandon
egoistic limitation and to take up the universal,
humane points of view into his motives. The interest of society
gradually displaces the individual interests, and there
will in time even be able to form so strong a feeling
for the universal that one may come to fear that it
will check individual development. The association of ideas
between the feeling of duty and utility, supported by the natural
force of sympathy, will be able to produce a psychological
power and a social efficacy that exceeds even the
greatest influence that has been exercised by a positive
religion.

Mill begins in his ethics as an Epicurean and ends
as a Stoic. A possibility of uniting the conflicting
% --- p. 58 ---
\setcounter{footnote}{0}%
\opage{58}points of view lies in the idea which in so many places emerges
in him, but which he does not lay as a foundation, the idea
of development. Only through it can his philosophy be freed
from its inner contradictions. At bottom it is thus
not happiness that is Mill's ethical principle, but happiness
developed to its ideal, where it is one with the ethical character.
It is this development itself that is the ethical
task, and whose necessity grounds the relation of duty.
The ethical motive is striving after ever higher development,
a development in which the individual and the cultural-historical
tasks stand in inner connection and interaction
with each other.\footnote[1]{A hint in this direction, in the direction of letting the idea of utility merge in the idea of development, is found in Mill in his work ``On liberty'' p. 24: ``I regard utility as the ultimate standard in all ethical questions; but it must be utility in the \emph{most comprehensive sense, grounded on man's permanent interests as a progressive being}.''}

Utilitarianism is empiricism in the ethical
domain. In the theory of knowledge empiricism found its expression
in the conception of the human mind as originally passive;
with the recognition of the original activity in the individual
empiricism and association psychology therefore dissolved. In
ethics it finds a corresponding expression, in that desire,
``the passive state of receptivity,'' as Mill defines
it, is taken as the point of departure, while the will, ``the active
phenomenon,'' is regarded as ``a shoot from the stem
of desire.'' But how can an activity develop
out of a passive state, and must not that recognition of the
original activity also have consequences with regard to the foundation of the ethical, so that the Stoics,
also as regards the original point of departure, are proved right
against the Epicureans? ---

It has earlier been indicated that Mill's whole theoretical
standpoint has its ground in his practical enthusiasm,
his zeal for progress and development. This zeal
% --- p. 59 ---
\setcounter{footnote}{0}%
\opage{59}shows itself clearly also in his ethics; it is perhaps his
impatience to get to the practical tasks
that makes us in his ethical definitions of concepts often
have difficulty in recognizing the great logician. But
above all, of course, the practical enthusiasm comes forward
when we come to the social science itself. In few
or none in our century has there stirred a stronger
interest in humanity and progress, combined with so
superior an intelligence and such varied knowledge.
Mill is not merely a philosopher, but also has a
name as an economist. His ethical ideals do not float
in the air, but are grounded by sober investigation
of the real conditions. The sharpest, prosaic understanding
has led him to consequences that seem to the
accustomed way of thinking overstrained and fantastic. Without being
able here to go more closely into the practical problems
with which Mill has occupied himself, we must nevertheless briefly
show the connection between his treatment of them and
his general philosophical standpoint.

The social science is, according to Mill's Logic, subject to the same
methods as natural science: the results of induction
form the foundation for a deduction whose results are
in turn confirmed by direct experience. Mill therefore
does not start from dogmatic presuppositions, but seeks to
find definite points of departure in experience. He
here often expresses himself quite skeptically. If there is question, e.g., of
woman's position in society, he does not start from
certain ideas of woman's nature, but even declares
that we do not at all know it. Women have
hitherto lived under certain restrictions, indeed even under
direct compulsion; what is called their nature is in the highest
degree an artificial product, the result of a forced pressure in one
direction and an unnatural excitation in another direction.
Certain of their faculties have always been developed by the power of circumstances,
while others have withered and apparently
died out for lack of nourishment and care. As long as
ethology is still in its infancy, as long as one still
% --- p. 60 ---
\setcounter{footnote}{0}%
\opage{60}has so little knowledge of the laws for the formation of character
and of the ways in which external conditions here intervene, so long at least can no positive proof yet be adduced
that the intellectual and moral
differences between men and women are grounded in
nature. Only that which cannot be explained by education
and external circumstances has one a right to call natural;
how far-reaching an influence these circumstances can exert,
one does not yet know at all.\footnote[1]{Kvindernes Underkuelse. Translated by G. Brandes. p. 37-40.} On the other hand the whole
historical development consists in a progressive abolition
of the right of the stronger, which for so long periods was the sole
and unconditional rule in all relations of life. Through this liberation
from external compulsion, which is the greatest peculiarity
of the modern world, it becomes possible for individuals to develop
their individuality without being stopped by artificial
barriers. The exclusion of woman from the rights and the
roads that stand open to men is only a remnant
of that brutal power which is now combated in most fields.


With the abolition of the right of the stronger every
external authority will gradually fall, every power that would set itself up
as guardian over individuals instead of letting
them fight the battle of life with their own hands and on their own
responsibility. According to the old view it was believed that the individual
would go astray when left to himself.
But the modern view teaches that the affairs in
which the individual alone is interested will not go well,
except when they are left to his own direction, and
that the intervention of authority only does harm when it is not
precisely a matter of protecting the rights of others. In the treatise
``On Liberty'' (one of Mill's clearest and most beautiful
writings, of which he himself also thought that it would
last longest of all his works) this
thought is developed further. A sound intellectual and moral development
is only possible on condition of freedom; when
% --- p. 61 ---
\setcounter{footnote}{0}%
\opage{61}society, be it democratic or aristocratic in its
constitution, wants to extend its influence over both the
mental and the bodily life, when it seeks to lead
individuals into certain directions of life, to instill in them certain
definite opinions, then the development of character is checked, and
higher independence and peculiarity become impossible.
Therefore Mill demands absolute freedom for thought and speech and for
the laying out of the plan of life as a whole. Only with regard to
actions that may affect the welfare of other people
can barriers on individual freedom be justified.
If this rule is not followed, the result will
be a pervading mediocrity; men will become
dwarfs, and one will soon learn that with little men
nothing great can be accomplished. The danger is especially great in democratic
states, since public opinion, which in them
becomes the ruling power, reaches far deeper and more imperceptibly
into the mental life than external statutes and
barriers set by force. Mill's endeavors therefore aim
at finding political means to assert the
rights of minorities. But the most powerful means he seeks
in a reform of public opinion, an absolute,
and not merely external and formal recognition of the right of
personality.

One must not conclude from this that Mill regards the role of
the state as played out, so that everything could now be left
to individual striving and competition. He disputes
such views as Buckle's, which derive all evil and
error in history from the influence of governments. He
maintains in his ``Political Economy'' that there is a great mass of
tasks that concern the interests of humanity and of future generations,
or also those members of society
who need external help, --- tasks that cannot be
solved except by the power of the state. The state ought not only to provide
for legal security and exercise control and arbitration in the mutual
affairs of individuals; it ought also to intervene where
either neglect on the part of individuals may occur,
as in education and instruction (only that it must not
% --- p. 62 ---
\setcounter{footnote}{0}%
\opage{62}make this a monopoly), or where the powers of individuals
do not suffice (poor relief, colonization). There can, in
his conception, be fixed no definite limits here;
one must keep to the general rule that the intervention of state power
is permissible only where its utility
is evident. And at the same time state power ought always to work in
such a manner that the individual powers are not pushed
back, but on the contrary are developed and educated through its
measures, so that government assistance becomes ``a
course of instruction for the people in the art of solving great tasks
by individual effort and voluntary cooperation.''

Brought up and trained in the school of Bentham and James Mill,
Stuart Mill was from his first appearance an enthusiastic
democrat. His political faith later underwent
a double modification. In part he soon came to think
that democratic development might entail
danger to individual development, that minorities would easily
be oppressed; this fear he already expresses
in his article in the Edingburgh Review on Tocqueville's book
``La démocratie en Amérique,'' and he came back later, as
already shown, to it in the treatise on liberty.
But in part he also saw that the social question
was more important than the political. Earlier he had
with the older political school regarded the right of property
as something fixed and absolute; but gradually he was led
more and more to see that it, like every right, is conditioned
by social relations. The right of property also could not
according to his whole philosophical view be a dogma,
from which one must necessarily proceed. It is an institution,
which has arisen in consequence of the natural course of
things, and which has its great advantages, but which nevertheless in and
for itself cannot lay claim to absolute validity.
Therefore Mill judges more favorably of the communistic and
socialist systems (especially Fourier's) than the economists
are accustomed to do, just as he finds it natural,
that communistic ideas have spread in a time like
ours, ``where a general examination of all first principles
% --- p. 63 ---
\setcounter{footnote}{0}%
\opage{63}is recognized as unavoidable, and where for the first time in history
the most suffering classes of society have a
voice in the inquiry.'' Most of the difficulties
that criticism can point out in the socialist systems apply,
in his opinion, perhaps in still higher degree to the present
conditions. When one says, e.g., that under a socialist
constitution everyone will shove as much work as
possible onto others, because he is not particularly interested
in it, Mill answers that the same must hold of nine
tenths of present labor: how small a part
of what is produced in England is produced by men
who themselves reap the proceeds of it? A factory worker has
less personal interest in his work than a member
of a socialist association. Besides, Mill maintains that
humanity is capable of a far higher degree of public spirit
than our age has become accustomed to regard as
possible. And in any case, he thinks, one could
hardly get a distribution that was more unequal and more unjust
than the present. --- The decision which of the two
systems, that of communism or that of private property, is to be given the preference,
must depend on which of them can best be
reconciled with human freedom and development. Mill finds
no necessity of abandoning the system of private property,
if only it is freed from all the evils which
are not inevitably bound up with it. If the laws did only
as much to mitigate social inequalities
and favor a greater distribution of fortune as they now
do to heighten those inequalities and concentrate possession,
it would be seen that the principle of private property does not
stand in any necessary connection with the bodily
and mental sufferings which the socialists regard as inseparable
from it. Likewise he disagrees with the socialists in their
polemic against free competition, in that he points out
that one has only the choice between competition
and monopoly, and that ``protection against competition means
the same as protection of sloth and idleness, as
% --- p. 64 ---
\setcounter{footnote}{0}%
\opage{64}an exemption from the necessity of being as intelligent
and active as other people.''

Mill counts himself in the Autobiography among the adherents
of socialism. Yet he is, as one sees, at odds with the socialists
on essential points. What is socialist in his view
rests on his demanding the intervention of state power
in order to solve the social question. The conditions of the workers
cannot be changed without a thoroughgoing alteration in
their whole way of thinking, and this alteration in turn presupposes
an improvement of their external conditions, which can give them
desire for and taste for nobler enjoyments and higher culture. Only
thereby can the dullness and thoughtlessness be eradicated
that now make it impossible for the proletariat to raise itself, in that
by its strong multiplication it constantly keeps itself
down. The measures that are to be taken must thus work
partly toward the development of intelligence, partly toward the abolition of
poverty. The first is attained by a national education of children
and by popular instruction, and as means to the
second Mill proposes large colonization schemes and the parceling out of
landed estates. When a generation of workers was thereby
placed under new and better conditions, there would have been
provided favorable conditions for the development of a new
generation with higher ends and greater energy. --- We
cannot here enter upon any appraisal of these
proposals, but restrict ourselves to pointing out the remarkable fact
that the same man who more strongly than anyone else has
pleaded the cause of free individuality so urgently demands
the intervention of state power in the social field. It
might seem a contradiction to appeal to these
opposite principles. But one sees easily that Mill posits
an inner connection between them. The highest end of the power of society
is to develop individual independence and
personality, to secure the individual's activity at the points
where everything can be left to its free play, and to nourish and
tend individual freedom where it has not yet been able to
unfold. It is of course far from the case that Mill
has overcome all difficulties. But he has the
% --- p. 65 ---
\setcounter{footnote}{0}%
\opage{65}great merit of having shown how wrongly one conceives
the nature of freedom when one holds that all intervention by
state power in the field of social tasks ought to be
over, and that everything could now be left to the ``natural''
course. There is still a long way to go before the truly
natural development can set in. The state, by
whose help man has been led out of the animal state of nature,
cannot let go of him before it has led him to
the human state of nature. Mill has finally also the
merit of having pointed out the significance of the ideal and ethical elements
for the social questions and thereby of having led beyond
the narrow horizon that generally asserts itself
in their treatment.

As the result of social development Mill imagines
a condition where both capital and population
are about stationary. Acquiring cannot be the
final task; humanity must some day come so
far that it sets the end above the means. One then no longer works
to increase capital; all inventions aim
rather at shortening the working time in order to win
time for freer and nobler pursuits; the intellectual, aesthetic
and moral life plays the principal role and claims
the greatest part of energy. The population is stationary,
because one sees that too strong an increase of it would
hinder the higher and nobler development. Even if there were
food and clothing enough for more still, there would nevertheless be
overpopulation, when men lived so closely that it became
impossible to enjoy solitude, and when free nature entirely
was taken into the service of human utility. ``A world
from which solitude was banished would be a poor
ideal. Solitude is an essential condition for depth in
thought and in character; and solitude in conjunction with
the beauty of nature is the cradle of thoughts and impulses that
not only benefit the individual human being, but which
society too can ill dispense with.'' --- Even if
the picture that Mill here sketches should be a Utopia, it has
its significance and interest, not only as graphic
% --- p. 66 ---
\setcounter{footnote}{0}%
\opage{66}expression of the goal of a humane endeavor, but also as a testimony to the noble, fresh and philanthropic mind
of the great thinker whom we --- despite all the criticism we have tried to bring to bear on his doctrine --- cannot
leave without admiration.

II. \emph{Alexander Bain} (b. 1818) was born in Scotland
and may be said to have continued the endeavors of the Scottish school
in the domain of empirical psychology. His doctrine can
be called ``the Scottish psychology as it would look
in the nineteenth century.''\footnote[1]{\emph{Ribot}: La psychologie Anglaise contemporaine. Paris 1870. p. 222.} But in place of the timidity of the older
Scottish Philosophers, their appeal to the pronouncements of ``common
sense'' and their inclination
to halt at certain so-called facts of consciousness
which they regarded as fundamental, there appears in their modern
successor a confident and comprehensive application of the inductive
method, a careful study of physiology and
anatomy as the foundation of psychology, and a sharper
criticism and analysis. Yet one still detects the aftereffect of the older
school in the empirical-descriptive method, in
a want of principled carrying-through of the concepts won by induction
and analysis, and of synthetic development
and presentation, by which a total view
of mental life might be gained. It is this, together with the fact that in his whole
standpoint he stands so close to Mill, that justifies placing
him among the representatives of pure empiricism. In
England itself he is reckoned to Mill's school. His chief works
are ``The senses and the intellect'' and ``The emotions and
the will,'' in which, following the descriptive method of botany and zoology,
he gives a natural history of the mind with
a multiplicity of particular investigations and acute
remarks that make these books highly instructive and interesting,
though the survey is often hindered precisely by the
% --- p. 67 ---
\setcounter{footnote}{0}%
\opage{67}multiplicity of the material. In a smaller work, ``On the study of
character,'' he has given a sketch of an ethology according to
Mill's principles.

We have already had occasion to emphasize that
Bain has corrected the sensualist conception of the mind
as wholly passive. It is a decisive proof
that the English philosophy of our time is not a mere renewal
of the sensualism of the last century (perhaps
rather sensationism). What was peculiar to the latter was
that it reduced thought to sense-perception and conceived
sense-perception as a purely passive process, a mere suffering and
receiving. In opposition to this the opponents of sensualism
generally conceived thought as something purely active and
distinguished it sharply from passive sense-perception.
This distinction rests, properly speaking, on the same presupposition
as sensualism. There must be an inner relation between
sense-perception and thought. When one acknowledges that we indeed
have no innate concepts, but that the basic forms of our
mental organization must be original, that we, as
Leibnitz expressed it, are innate to ourselves, then one must
also acknowledge that the mental energy which manifests itself in
the higher forms of knowledge must be present in the first
manifestations of mental life, in the elementary sensations of sense,
and that in the course of development it arises out of these and
becomes conscious of itself. This is just what Bain teaches: ``Sense-perception
is already a function of the understanding; it is the
purely outward manifestation of the ideal processes'' (Senses and
Intellect p. 339). And these ideal processes have their
germ in the very starting point of mental life.

Even in the lowest manifestations of mental life there are, says Bain,
elements of mind or thought (intelligence).
The primitive fact that expresses the essence of the mind
is the capacity to distinguish, to apprehend differences (discrimination).
This is an original mental activity, which
cannot be explained further, but on which all consciousness rests.
It has as its presupposition the capacity to retain an
impression even when its source is no longer present, so
% --- p. 68 ---
\setcounter{footnote}{0}%
\opage{68}that it can be compared with another impression. Psychological analysis
must therefore halt at it. Already Thomas
Hobbes impressed upon us this basic fact, in that
he taught that without change and motion neither
sensation nor consciousness was possible.\footnote[1]{Per sensionem vulgo intelligimus aliqvam de rebus objectis per phantasmata judicationem; phantasmata scilicet comparando et distingvendo. . . . Adeo sentire semper idem et non sentire, ad idem recidunt. \emph{Hobbes}: De corpore IV, 25, 5.} An action without
change produces no impression on us. We do not notice
the motion of the earth in space, but we do notice
the motion of the ship on which we sail, and still more that of the carriage
in which we ride. The sense of difference is thus
the original mental element. All consciousness rests,
properly speaking, on contrast and surprise and is the more
intense, the greater the contrast and, as a consequence, the surprise.
After a deep darkness the light is felt more strongly
than otherwise.\footnote[2]{In this fact the interesting investigations of Weber and Fechner on the degrees of sensation (of consciousness) take their point of departure. Cf. ``Philosophien i Tydskland'' p. 253---255.} And it is with knowledge as with
consciousness in general; the more differences are apprehended,
the clearer and more determinate is knowledge. No characterization
of an object is possible except through the determination
of its differences from other objects. Even in the
most elementary sensation of sense there is such an apprehension
of the different, and to that extent a judgment, though
in speaking of judging one commonly thinks of the more highly
developed forms of the power of discrimination. Bain thus agrees
with William Hamilton and the Scots\footnote[3]{Lotze, too, teaches that sense-perception is a judgment. See op. cit. p. 295.} in the doctrine that
every act of consciousness contains a judgment. Sense-perception
accordingly presupposes understanding (intellect). The old
proposition: Nihil est in intellectu, qvod non erat in sensu
--- is utterly false. There is, as will soon be shown, a
% --- p. 69 ---
\setcounter{footnote}{0}%
\opage{69}multitude of concepts that do not have their origin in
sense-perception.

The first sensation is not a sense-perception of
anything external, as the earlier sensualist school
assumed. ``Motion precedes sense-perception and is from
the beginning independent of every impulse from without. Activity
belongs more intimately and more inseparably to our nature
than any of our sensations of sense do, and is in fact
a co-operating factor in each of our sense-perceptions, in that it gives
them the character of composite phenomena, whereas it
is itself a simple and elementary property.'' But whence
comes the original motion? A series of facts
shows that the nerve centers can send out impulses to the muscles
merely by virtue of their inner tension, which makes a discharge
necessary. When the organism has received proper nourishment,
the nervous system must have an outlet for its increased energy.
We are thereby brought to see mental force in its connection
with the other cosmic forces. Through nervous force
it stands in connection with the chemical forces,
with magnetism, electricity, heat and mechanical force,
which all point back to the same origin and can all
pass over into one another.

The first mental phenomenon is spontaneous motion,
which we already encounter in the fetus in the mother's womb,
but which also manifests itself in the newborn and in older
individuals in the liveliness that is called forth by good and
abundant nourishment, and in the spontaneous awakening from sleep.\footnote[1]{Bain here relies on the famous physiologist Joh. Müller.}
This spontaneous motion we become conscious of through muscular sensation,
which is thus the first sensation, a
sensation that differs from sense-perception proper
in that it has its source not in an external but in an internal
impulse. In muscular sensation, as in every
sensation, we can distinguish two elements, a feeling element and a
cognitive element. These two elements stand in an inverse
% --- p. 70 ---
\setcounter{footnote}{0}%
\opage{70}relation to one another. The more I feel the pleasure or
pain which the muscular motion calls forth, the less am
I able to perform the discriminating and distinguishing on which knowledge
rests.

The original motion is without aim and means, is
an involuntary discharge of nervous force. But the direction of the motion
is gradually determined by the impressions it calls forth;
by them it is led into definite channels. One can
observe this in newborn lambs. They involuntarily try
to move their limbs, but this motion goes now
sideways, now forward; only gradually do they learn that motion
forward is the most natural. When the shepherd leads them
to the mother ewe and lets them touch her, the impression of warmth
which they thereby receive makes them seek to retain or
to resume this position. When their mouth is brought
to the mother's udder, and the involuntary sucking movements
have begun, the enjoyment they feel in this will spur them on to
continuation or repetition. Soon they recognize the mother at
a distance and can perform the movements that are required
to get to her. The original, spontaneous motion
is thus the presupposition for the sensations of sense proper
being able to teach the animal anything.

From the muscular sensations arise concepts such as resistance,
inertia, weight and the mechanical properties of nature.
They are measured by the degree to which we must exert ourselves.
The difference between 4 and 6 pounds' weight we apprehend by
the feeling of increased exertion which we must apply when
we have to carry 6 pounds instead of 4. The concepts of
duration and extension have their origin in muscular activity
in connection with the sense of touch.  Only by
the aid of this sense do we gain definite points with which to
measure motion, both as regards space and time.
The combination of the muscular sensations with sight,
hearing and the other senses gives us the more composite
sense-perceptions. By means of these knowledge develops,
but its germ is present from the first.

Already at the first step beyond, the first
% --- p. 71 ---
\setcounter{footnote}{0}%
\opage{71}closer determination of the original spontaneous motion,
the fundamental laws of mental life reveal themselves, the laws of
association of ideas, by virtue of which mental phenomena enter into so
close a connection that the one calls forth the other, ---
laws which hold for the whole domain of mental life from the lowest
to the highest stages, and which are, in Stuart Mill's phrase,
for the mental world what the law of gravitation is for the
external world. These laws can be reduced to two: the law of contiguity
(contact) and the law of similarity (likeness).
Mental phenomena that occur simultaneously with
or immediately after one another tend to call each other forth,
when only one of them occurs later. And likewise
mental phenomena that resemble one another tend
to call each other forth. Older psychologists, in particular James
Mill, tried to explain all mental phenomena by
the law of contiguity. But this law presupposes the law of
similarity, as William Hamilton has shown. For in order
that a phenomenon may be able to call forth the idea of
another phenomenon with which it has previously occurred
together, the former must itself be recognized, known as
the same that previously occurred together with the other. The law of
similarity is thus a necessary condition for the law of contact
to be able to operate. But where the similarity is
complete, where it can be said that it is the same phenomenon
that recurs, there one need not specially
emphasize it. Conversely contact forms the tacit
presupposition for the law of similarity to be able to operate; for in order
that the similarity between two phenomena may be apprehended, they
must first be brought together before consciousness. That there really
are two different laws, although they may thus seem to
pass over into one another, is seen, among other things, from the fact that decisive mental
differences depend on which of the two laws is
the predominant one. Association by contact in space or
time (contiguity) plays a part especially in the accumulating of
knowledge and in the execution of practical affairs.
We then remember things in the same order in which they
are wont to present themselves; habit and routine rule.
% --- p. 72 ---
\setcounter{footnote}{0}%
\opage{72}Association by similarity, on the other hand, is predominant in the discoverer,
the thinker and the poet, and conditions on the whole
productivity and originality, whereas the law of contact
especially conditions receptivity.

Let us consider first the law of contact or contiguity.
That states of consciousness which have once occurred
together later call each other forth cannot be explained further.
It is a fundamental mental fact, which in turn points
back to a fundamental peculiarity of nervous activity.
Like the ideas, so too the corresponding
nervous states must be connected. The idea of a sense-perception
is merely a fainter sense-perception, resting on the same nervous current
as the original sense-perception, only that this
current too is fainter than in the original
sense-perception. Repetition forms, as it were, channels for these
currents, so that they can move more easily in
certain directions than in others. A vivid idea of
a certain motion can involuntarily call forth this motion;
the idea of an abyss can arouse dizziness.
Indeed, thought itself is properly only a recalled, an arrested
motion. The brain is therefore not a storehouse in which thoughts
and ideas are supposed to lie heaped up. Nor can there
be innate ideas; all ideas rest on function
and motion and do not exist without them, and so not
before them either (cf. Schelling; see above p. 12 f.).

Association by contact can take place either between
states of consciousness of the same kind (e.g. a sound and another
sound, a motion and another motion) or of
different kinds (e.g. color with resistance, motion with
distance, etc.). Of the first kind the association
between different single motions is especially important. Almost
every significant manifestation of life in animals and men rests
on very composite motions, which come into being only gradually.
One of the most important examples of association
between states of consciousness of different kind is the
process by which the idea of an external world arises.
In general one regards this idea
% --- p. 73 ---
\setcounter{footnote}{0}%
\opage{73}as natural and instinctive and has no inkling that it has
a very involved and composite origin. Many different
motions and sensations must, by frequent repetition,
enter into a firm connection in order for that idea to be able to
arise. It is the sensations of sight and touch in union
with the muscular sensations that produce the ideas
of solidity and extension, which are the essential
elements in the idea of external material things. In particular
Bain here lays stress on the sensation of our own
motions. When we move our limbs and as a consequence
receive certain sensations of sense, we derive these from
external things. If our sense-perception were perfectly passive,
our conception of the universe would be like a dream, in which
one image replaces another. Our sense-perception is, however,
as shown in what precedes, rather active than
passive. And we believe in an external reality because we know
that when we perform a certain motion, we shall thereafter get
a certain sensation of sense. ``Our external world is the total sum
of all opportunities for developing energy.'' The assumption
of an external world is thus a generalization or abstraction
founded on particular experiences, by which
we sum up the past and predict the future;
but by an illusion our consciousness then posits
this abstraction as an original and independent source
and cause of those particular experiences, much in the same
way as Plato posits the ideas of the Good and the
Beautiful as independent realities outside the particular
good and beautiful things from which they are abstracted. But
such an absolutely independent external world we have no
ground to assume. ``The world can be known only in its
relation to the mind. Knowledge is a mental state;
the concept of a material thing is a mental thing. We
are unable to investigate the existence of an independent
material world; an attempt at it would be self-contradictory.
We can speak only of the world which
presents itself to our mind. By an illusion which language
carries with it, we believe ourselves able to contemplate a world
% --- p. 74 ---
\setcounter{footnote}{0}%
\opage{74}that does not enter into our own mental existence. But
this very contemplation is itself a mental activity.''

There is thus in Bain the same approximation to
Berkeley as in James and Stuart Mill. Yet one can draw from
him, no more than from these, the consequence
that everything must dissolve into semblance and delusion; for he holds fast to
a pervasive duality in the view of the world,
resting on the opposition between subject and
object, idea and reality. The basic activity of our knowledge
is, according to Bain, discriminating and comparing.
It therefore always presupposes plurality and change.
We know only relations, difference and similarity.
A thing can be apprehended only when there is a background, a
foil, in relation to which it can be judged and characterized.
With the discovery of new foils, of new oppositions
and differences, knowledge advances. In the last
instance all oppositions can be reduced to the opposition
between consciousness and existence (actuality), subject and
object. A removal of this opposition is impossible; an
absolute subject and an absolute object are equally unthinkable.
Berkeley's error consisted in letting the actual be absorbed in the
ideal, the objective in the subjective. This opposition
one must hold fast, even if one rejects the assumption of
a world absolutely separate from the mind. Actuality and
ideality are akin in their essence; idea and real
motion rest on the same nervous function. A whole series of transitions
can also be pointed out from the objective
pole to the subjective; but the opposition itself between the
two poles cannot be removed. ``Extension is the most
pronounced objective fact, pleasure and pain the most pronounced
subjective states. Between the consciousness of
extension and the consciousness of pleasure there is the greatest
difference which human experience can posit in
the whole existing world. It is the extreme object
and the extreme subject; and the carried-through analysis shows
us that the extreme object rests on the sensation of
applied muscular energy.''

% --- p. 75 ---
\setcounter{footnote}{0}%
\opage{75}Association of ideas on the ground of similarity plays a part especially
in the higher intellectual activity, whereby order
and coherence are brought into the material accumulated by the working of
contiguity, and whereby this finds its explanation and
is conceived in new ways. The identifying (seeking similarity and
unity) force operates in all thought, but can be
present in very different degrees. Its strength shows itself
in the magnitude and multitude of the differences through
whose cover the agreement must be sought. On it rest
classification, abstraction and definition, whereby one identifies
several objects by emphasizing a common feature
or a common group of properties as a special object
of consideration. The more the law of similarity operates, the more perfect
become abstraction and classification. What an
advance there is, e.g., between Aristotle's division
of animals into water animals, land animals and air animals and Cuvier's
system, which is grounded on the internal structure of animals and according to
this divides them into vertebrates, articulates, mollusks and zoophytes.
In place of the external and accidental similarities the internal,
essential agreements come. By classification, abstraction
and definition we form general concepts. General
judgments we form by way of induction. Here we do
not, as in abstraction, turn our gaze to a single
property or a single group of properties, but we
emphasize a certain conjunction, a certain connection of two
or more properties and facts. When we state the constant coexistence
of two facts, we express a
law of nature.

The laws of contiguity and similarity are the fundamental laws
of association of ideas. But there are given not only simple but also
composite associations of ideas. When I, e.g., contemplate
a mountain, this may lead my thought along many different
paths. According to the law of contiguity I may be led
to think of the neighboring mountains, of valleys and villages in the vicinity,
of the minerals of which the mountain consists, of historical
memories that attach to it. According to the law of
similarity I may come to think of similar mountains,
% --- p. 76 ---
\setcounter{footnote}{0}%
\opage{76}that I have seen or read about, of paintings and poetic works,
etc. Which of all these directions thought chooses
depends on which of them is strengthened by connection and
combination with other ideas. When I, e.g.,
am a historian by profession, the historical memories operate
most strongly; if I am instead a geologist, the physical character of the mountain
will captivate me most, etc. Associations that
singly are too weak can gain the power when they operate
together with several others. Most often feeling and
will play an important part here, either hindering or furthering.
When the composite associations are of a purely accidental
character, they serve precisely to lead thought away from the
essential. From this all self-contradiction and inconsistency
are to be explained. Only when the force of similarity is sufficiently developed
will it free the mind from contradictions by placing
the conflicting side by side and drawing attention to
the discrepancies. When one, e.g., has two opinions
that stand in contradiction with one another, and between which one
must therefore choose, it is the law of similarity that grounds
the necessity of choice. But this working is hindered by the influence of inclinations
and moods which have attached themselves
to certain ideas. Hence it comes that absolute
consistency is so seldom found in human views
and characters.

The laws that have so far been pointed out for the associations of ideas
concern the reawakening of earlier ideas. But
the mind is not bound to what it has previously received; it
can form new connections. Then we have what Bain calls
constructive association. The creating of imagination rests, as
Hobbes already indicated, upon this; we have the idea
of gold and the idea of a mountain, and the fairy-tale poet
now forms from them the idea of a mountain of gold. Not only
art and poetry rest on constructive association, but also
all conception of nature presupposes it. For the real nature
we understand only when we unite in a total coherence what
analysis and abstraction have separated. The same
holds when we wish to enter into the moods and
% --- p. 77 ---
\setcounter{footnote}{0}%
\opage{77}passions of others; we understand them only when we construct their
mental state by combining the particular elements
(fear, anger, longing, etc.). In all such
constructions, even the most original and creative, it is
no wholly new principle that operates; what is needed is only
a vigorous working and co-working of the laws of contiguity and
similarity. ---

Bain follows the old threefold division of mental phenomena
into knowledge, feeling and will, not in the sense
of special faculties or powers, but as groups of related
phenomena, as natural orders, each of which
in turn comprises several subdivisions. The life of knowledge
is comparatively the easiest to understand; analysis can more readily
be carried through in the domain of clear consciousness, where the
particular elements are as a rule known. More difficult
becomes the conception of the life of feeling and of will, because the
obscure natural ground and the unconscious life play so great a
part here. Here Bain's descriptive and analytical
method also does not suffice; it brings forward only the phenomena
that move on the surface, but does not demonstrate
the deeper coherence. Spencer has, in his criticism of
Bain's work ``Emotions and Will,'' emphasized that not only must
the analysis be carried through to more primitive forms than Bain
has done, but that there must be laid as the basis of the whole study of
the life of feeling a comparative investigation of its
appearance in the different species of animals and races of men
and in the same individual at different ages, so that one
may come to know the conditions and course of development of every feeling.
Yet Spencer himself has treated this difficult psychological problem only
in passing.

As the basis of his classification of the feelings Bain lays
three criteria, namely, besides the character of the states of
consciousness themselves, their physiological effects and their aesthetic
or moral forms of expression. In this way he comes
to set up the following classes. 1) When the impressions
agree, harmonize, so that the feeling on the whole
can have its free course, pleasure arises; in the opposite
% --- p. 78 ---
\setcounter{footnote}{0}%
\opage{78}case a feeling of constraint and pain. 2) Perhaps the
most elementary feeling is wonder, the feeling of surprise;
it is, as we have already heard, tied to the very
fundamental law of consciousness, according to which relation, opposition
and change are necessary for all apprehension. 3) Fear
and all that belongs to it, superstition, anxiety, etc., together with
its opposite, courage in its different kinds. 4) The feeling
of tenderness and benevolence, which is of a predominantly
passive character, and sets the glands rather than the muscles
in motion; the feeling of tenderness (the tender affection) must
not be confounded with sympathy, which is a very composite
feeling that can develop out of the feeling of tenderness, but
does not necessarily presuppose it. Sympathy presupposes
that one has entirely entered into another's state and
regards it as one's own; but the feeling of tenderness may well
restrict itself to one's own person, and we then get 5) self-love
and self-satisfaction, which have their ground in
the passive enjoyment and delight in one's own self. 6) The feeling
of our power and of the superior activity we exert,
in whatever domain it may be (politics, wealth, artistic
and scientific proficiency, etc.). This feeling can
fuse with self-love, but the two are nevertheless
different in their essence. There are men, especially hard and
stern natures, in whom the consciousness of power plays a great
part, without their being able to be called egoists or
self-admirers. 7) The feeling of anger, where it is not a
pleasure that calls forth another pleasure, but a pain that
brings about another pain. 8) The feeling of pleasure that is
brought about by eagerly pursuing an aim (in earnest or in
play), by exertions in which one is wholly absorbed in one
object (emotion of pursuit; plot-interest). 9) The intellectual
enjoyment, the feeling of pleasure that the activity of thought
calls forth, in particular through the discovery of similarity,
agreement and consistency; on this rest the joy
and interest in scientific discoveries, poetic
images and industrial ingenuity. On the other hand inconsistency
and disharmony call forth a painful sensation especially
% --- p. 79 ---
\setcounter{footnote}{0}%
\opage{79}in thinking natures.\footnote[1]{``There is a kind of pain which is characteristic of intellectual natures, and which is very little felt by others; it is the pain of inconsistent ideas, opinions and principles. Strong feeling united with a less developed understanding easily carries one past the offense which contradiction gives; but on a mind endowed in the opposite way inconsistency acts like a grating sound and becomes a spur to intellectual exertion, which heightens the natural inclination to let the ruling faculty of the mind operate. To this pain we must ascribe an important function, namely the incitement to strive after truth and not merely after manifold intellectual images and ideas. . . . Herein lies a great motive to carry through scientific investigation and to root out popular errors.'' \emph{Bain}: On the study of character. p. 272 f.} These feelings become of importance
also in the moral respect by leading to truth and justice.
--- To these elementary feelings are finally added
the most composite of all, the aesthetic and moral
feelings.

Already older English investigators, Alison and after him
James Mill, applied the laws of association of ideas to explain
the feeling for the beautiful and the sublime. Coleridge,
following Kant, drew attention to the difficulty of
distinguishing between the beautiful and the agreeable, when one
confined oneself to the mere associations of ideas. The English
aesthetician John Ruskin lets the associations of ideas serve as
means of calling forth higher ideas; beautiful is, according to him,
what recalls divine ideas and attributes
(Infinity, Unity, Peace, Harmony, Purity, etc.) which
find their expression in the world.\footnote[2]{\emph{Milsand}: L’esthétique Anglaise. Etude sur Mr. John Ruskin. Paris 1864. p. 149. Cf. \emph{Stuart Mill's} Note 48 to James Mill's ``Analysis'' II.} Still more than Ruskin
Bain here points back to Kant. In order that an impression
on the whole be of aesthetic nature, he says, it must
not be purely sensuous, since it then takes on an interested
and egoistic character, which makes it purely individual
and renders a pure and common enjoyment impossible. Therefore sight
and hearing are the most aesthetic senses; with these it is
easiest possible to disregard the egoistic drive. The
% --- p. 80 ---
\setcounter{footnote}{0}%
\opage{80}beautiful proper rests on harmony. The sublime rests
on sympathy. Our thought goes involuntarily in the same direction
as the impressions: the great in nature awakens the thought
of the great in us, especially in the highest forms
of the human mind. The true sublime is human power:
the sublime that we find in nature we apprehend only by
clothing it in human forms and interpreting it according to
human ideals.

The moral feelings are, according to Bain, still more
composite than the aesthetic. He polemicizes against the
one-sided attempts at explanation which derive all actions from
a single drive or side of man's essence. He
polemicizes likewise against the assumption of a universal moral
law that should be able to serve as rule and standard in all
particular cases. Where do we find such a law? We can
set our watches by the Observatory at Greenwich; but where
do we get a common signal to regulate our actions? There
is as little a general moral law as there
is a general reason; each individual has his reason
and passes his judgments according to his individual standard.
There is as little anything good in itself as there is
anything beautiful or true in itself; these ideas do not exist
outside the individual consciousnesses: even if individuals
agree in their judgments, this agreement
can as little be derived from a truth existing in itself
as one can ascribe to a choir a general
voice, different from the voices of the individual singers. ---
On the contrary, we need only cast our eyes around us and
read history to find the powers that really rule in
the ethical world. It is not an invisible moral
ideal, but real authorities, that operate. We find no original
feeling of duty and justice; the original
springs of action are fear of punishment, respect for authority,
the feeling of tenderness and sympathy, self-love, etc. Of
these egoism will always gain the upper hand, unless
the egoistic drive itself is brought into the service of duty,
which happens through the relation to a punishing authority.
% --- p. 81 ---
\setcounter{footnote}{0}%
\opage{81}Obligation is the fundamental moral phenomenon. But obligation
is in question only where transgression of a law
entails punishment. Responsibility is thus punishability. And
punishment in turn presupposes an external power, a real authority,
from whom the command proceeds, and who can chastise its
transgression. But by the power of association of ideas the
original motive is gradually lost sight of; the virtues that
at first were practiced only under constraint gradually become ends in
themselves, just as money becomes an end in itself, whereas
originally it was desired only as a means to the satisfaction
of certain needs. The associations of ideas have a tendency
to produce new centers of gravity, where often more powerful
feelings can be concentrated than the original motives
were able to call forth. Thus conscience emancipates
itself gradually from the fear of punishment, especially by
the help of the working of knowledge and sympathy. The submission
and feeling of obligation which the individual has learned
in the relation to the external power are now transferred to the relation
to the inner, ideal power. Thereby man becomes mentally
free, raises himself above the slavish way of thinking.

As a contribution to the ethical history of man's development
this description is excellent. Far from, as Bain thinks,
disproving the self-validity of the ethical law, it confirms
precisely this. For it demonstrates that there is a law of evolution,
by virtue of which man raises himself from lower
to higher and nobler forms of mind. This law is a law of nature,
but a law of nature which the self-conscious mind can make
its rule of life, according to which it can form its ideal; here
is then the transition from law of nature (in a narrower sense) to
moral law. It is by virtue of this law that man can
turn with criticism and protest against authorities, over against
which he perhaps outwardly must succumb; it is by virtue of
this law that the authorities struggle among themselves, and the one
supplants the other. Under such conflicts it shows
itself clearly that the real power is the inner, the
ideal power, which uses those external forces as means for
its own unfolding. According to Bain's doctrine one would have to conclude
% --- p. 82 ---
\setcounter{footnote}{0}%
\opage{82}that morality ceases when the authorities lose their power.
This consequence is also drawn by Kirchmann, a
shrewd German author,\footnote[1]{I. H. v. \emph{Kirchmann}: Die Grundbegriffe des Rechs und der Moral. Berlin 1869.} who proceeds from similar
presuppositions to Bain. The ethical is then itself only
a means, which falls away when the end is reached. But
if the ethical has its ground in man's essence
as disposed to development, then there will never come the
time when it could be replaced by other ideas, --- not
even in ``the stationary state.''

According to Bain a definite distinction must be drawn between duty and
merit. What one is obliged to do is only what one
would be blamed or punished for omitting; what one
does beyond this is merit. This is not merely
consistent, but at the same time in a high degree characteristic of the
empirical morality, which grounds the ethical on factual
conditions (authority and drive) and therefore cannot embrace
the ideal striving which goes beyond these, --- and this
striving is nevertheless the highest form of the ethical. On this
point empiricism is in agreement with Catholicism, which has a
rule for ordinary Christians and another for saints. Mill
admires in this respect too ``the wise and clear-sighted
men who created the Catholic Church.'' Protestant
ethics has lodged an energetic objection against this
opposition between duty and merit. We are obliged
to acquire for ourselves all the merit we can; --- even
when we do so, we have still done only our duty. There
is in these concepts a dialectic which demonstrates the insufficiency of the empirical
morality. This is one of the points
where French positivism departs from English empiricism.\footnote[2]{Cf. \emph{Littré's} polemic against utilitarianism: La science au point de vue philosophique. 2éd. p. 340 f.}
The energetic Auguste Comte, glowing with enthusiasm, who otherwise admires medieval
Catholicism, agrees on this point with Protestantism
% --- p. 83 ---
\setcounter{footnote}{0}%
\opage{83}and demands of everyone an absolute devotion in the service of humanity,
a complete surrender of egoism. The empirical
school, which with Bentham and Mill builds on ``the
well-understood interest,'' cannot acknowledge such an obligation,
indeed probably even regards self-sacrifice and disinterested
action with a suspicious eye.\footnote[2]{``Disinterestedness is, strictly speaking, a kind of unreason or insanity with regard to the individual's own self, a contradiction of the innermost nature of the feeling being, which moves toward pleasure and away from pain. When we abandon this fundamental principle of our constitution, we stand constantly on the border of absurdity and lawlessness. For when a man has ceased to care for his own welfare or to act in the direction of his own pleasure, why should he then care for the welfare of anyone else? He has denied felicity as an end and must, if he is consistent, extend the same standard to all around him.'' \emph{Bain}: On the study of character. p. 292.}

Bain's ethical principles demonstrate the deficiency of
his psychological investigations. He has not carried through the
analysis of the feelings in such a way that it led to discovering
the fundamental drive in human nature, but remains
at a plurality of primitive feelings. A scientific
ethics must have its real foundation in that fundamental drive.
The vacillation of the concepts of utilitarianism has precisely
its ground in an obscurity with regard to the original
psychological starting point. Next, Bain has also not
emphasized the central significance that knowledge has for
the purification of the ethical consciousness; it is by the labor of thought
that there dawns upon consciousness an idea
of an ideal which points beyond the finite, the factual and
the egoistic. ---

The original activity which, according to Bain, is the starting point of mental life
receded into the background in the investigation
of the life of feeling, but it must by the nature of the case
again become of importance in the investigation of the nature of the will.
The earlier association psychology had properly speaking
wholly to deny the existence of the will. According to Bain the development of the will
has two presuppositions: the spontaneous motion, which
% --- p. 84 ---
\setcounter{footnote}{0}%
\opage{84}precedes all sensation, and the bond which makes
a certain motion follow upon a certain sensation. The birth-hour
of the will is the moment when a sensation for the first
time determines an action. The involuntary motion
becomes, by association with the sensation of pleasure or
displeasure, a voluntary action. In this way the individual gains
the capacity for more composite motions and is trained to
mastery over his limbs. Gradually the transition
from sensation to action no longer takes place instantaneously.
Between the two there interposes itself the idea first of the action's
direct, then also of its indirect consequences.
A struggle arises between the different ideas
and the impulses they give; against the present or near
pleasure and pain there is set the idea of possibly occurring
stronger sensations. This debate, during which
the action is suspended, is deliberation. It was originally
called forth by the experience that hasty acting brought pain
with it. By it the domain of the will is extended, and a higher
development, an ennobling of the motives, becomes possible. The
voluntary actions are distinguished from reflex actions in
that consciousness steps in between the impression and the action.
Thereby a mastery becomes possible not merely over
our external actions, but also over our thoughts and
feelings, which indeed are themselves properly motions. But this
mastery can be exercised only indirectly, by a roundabout way; we
check the feeling by holding back its outward signs or
by turning our attention toward certain ideas
that can call forth other feelings.

The older Scottish school (Reid and his disciples, indeed
even W. Hamilton) assumed the will, in the sense of free
mastery over the limbs, to be an original faculty.
Bain declares this an illusion. The will has its germ
in the original activity, but it must go through a
long process in order to become what is commonly
called free will. Of freedom there can here be no question
in the sense that a power different from the particular mental phenomena,
an inner self, should intervene; the series of
% --- p. 85 ---
\setcounter{footnote}{0}%
\opage{85}associations and motives cannot be interrupted. Bain
treats this question like Stuart Mill and polemicizes
like him both against the doctrine of freedom and against the
doctrine of necessity. In particular he rejects in a striking
way the appeal to the testimony of consciousness, to the feeling
of freedom. Immediate consciousness, he says, can
teach us no more about the psychological laws than
immediate contemplation of the heavens can teach us about the
astronomical laws. Consciousness acquaints us with the particular
mental phenomena, but does not teach us their
nature and coherence, which can be found only by investigation
and analysis. In any case we have no infallible
authority in it. Infallibility and knowledge exclude
one another. The will is, as has been shown, a very
composite phenomenon; in the combination of the many
details that belong to apprehending it, too many errors
can therefore be committed for us to be able to build
upon an immediate assurance. ---

Closely connected with Bain's theory of the will is his
conception of the concept of belief or assumption (cf. already
above p. 73). What is required for us to believe
anything? James Mill, with whom Herbert Spencer here sides,
explained belief as a firm and indissoluble association
between two ideas. But according to Bain, with whom Stuart
Mill has declared himself in agreement, belief cannot be explained by
being regarded as a purely intellectual phenomenon. It is
the connection with a practical end, with an acting,
that distinguishes belief from a mere fancy. The most primitive
form of belief is the expectation that after a
certain action a certain effect will follow. The original
activity has an irresistible urge to discharge itself, and
the discharge takes place naturally in the direction in which it
has previously shown itself to bring satisfaction. The hungry
beast of prey keeps to the path on which it has before
come upon its prey. The sight of the brook lures the thirsty,
because he expects there satisfaction of his craving. The
higher forms of belief too have their criterion and
% --- p. 86 ---
\setcounter{footnote}{0}%
\opage{86}standard in the acting to which they lead; those who do not hope, fear
and act in accordance with their religious belief
imagine only that they believe, but do not in reality believe.

From the original activity there therefore springs a
strong tendency to believe. Here lies the germ of all generalizations,
all general propositions: the first experiences
are raised by the primitive drive to general laws. Thus
arise the intuitive or a priori assumptions; and thus
in particular the idea of the Infinite and the
Ideal has arisen. But advancing experience corrects
step by step the original intuitions. The beginning
is the spontaneous assumption, the intuitive view; but
soon doubt rises at the conflicting cases which experience
points out. Experience is therefore the only certain source of knowledge,
the real authority by which we must measure
all our assumptions. Our views thus do not grow with a
quiet and even growth, but form themselves under a constant struggle
between the spontaneous, generalizing and idealizing intuition
(Bacon's antecipatio naturæ) and the weighing and
testing experience. We have a strong tendency first to
believe and afterwards to prove, and it will probably last
long before the disappointments which this procedure
brings with it come to an end.\footnote[1]{``Man has certainly an unbounded longing, and the metaphysical doctrine of the Infinite corresponds in a way to it. But in real life we find very few of our wishes fully satisfied, not even those which are most honorable to the human mind, such as the drive for knowledge, longing for self-development, the drive to do good. How can we then assure ourselves which longing carries its own satisfaction with it? Especially since the intuitive aspirations go in very different directions in different men, and cannot all be equally true.'' \emph{Bain}: The emotions and the will. London 1859. p. 587.}

This last development characterizes Bain's whole standpoint.
The original activity, which is his psychological
discovery, and by whose acknowledgment he decisively leaves
the older sensualist school, he uses only as a
starting point, from which he in turn strives to free
% --- p. 87 ---
\setcounter{footnote}{0}%
\opage{87}himself, in that he lays the chief weight on experiences and associations of ideas.
Besides that original activity, which
still belongs to the pre-conscious life, he assumes yet another
original mental element, which underlies all
manifestations of consciousness, namely the original sense of difference;
without it there can be no question of experience
or association of ideas. But we thus halt at an opposition
between the active and the passive mental phenomena,
which nevertheless in every real mental function are to go
together, without any attempt being made to show
how the relation between them stands. Sense-perception cannot
take place purely passively, when it has the spontaneous activity
and the elementary function of the understanding as its presupposition.
And just as little can association of ideas be a purely passive
and mechanical act. In particular the law of similarity decisively points
back to the original energy with which the mind
strives after unity in its view and its thought. It
shows itself at all points that receptiveness presupposes
activity, devotion freedom, just as on the
other hand what is received heightens activity, devotion
purifies and ennobles freedom. We thus come also
here to the result that the empirical leads beyond itself,
both by the original ideal presuppositions on which it rests,
and by the ideal end in whose service it labors.
The ideas of the Infinite and the Absolute have, precisely
according to Bain's psychology, their source in the innermost
essence of man. Then it can only be apparent that the
\emph{human} experience disproves them. Experience
can refute the conditioned and historical forms in which
the ideas appear; but this refutation can only become
the introduction to the labor toward a new and more perfect expression,
corresponding to the new presuppositions. Bain's striking remark
that infallibility and knowledge exclude one another holds
nowhere more than in the knowledge of the Ideal and
the Infinite; --- and yet he himself seems to proceed from the assumption that we
cannot attain such knowledge, because we can attain nothing infallible!

% --- p. 88 ---
\setcounter{footnote}{0}%
\chapterhead{\opage{88}Chapter Three.}{The Critical School in England.}


Empiricism leads, as has been shown, beyond itself, both in
a negative and in a positive way. Negatively ---
for by setting experience and association of ideas in opposition
to inner necessity it cuts itself off from the possibility
of any grounding, of getting beyond the
contingent. Positively --- for its own investigations lead to
the recognition of the original activity of the human being
and of human thought as a primitive and irreducible
element in all knowledge. There is a close connection
between the two results, the positive and the negative.
For the ultimate grounding that human
knowledge is able to attain can be achieved only by going back
to the ultimate presupposition of all knowledge, which cannot be sought
elsewhere than in the very nature of the knowing being itself
and its original organization. And here we stand precisely at
Kant's fundamental thought. That we cannot, as Mill supposed,
conceive an absolute disorder, an unconditioned contingency, but
that even imagination refuses its service to this, comes from the fact
that that original element of thought brings about a connection
in all sense-perception and knowledge. And on the other
hand it now becomes clear how we can pass judgments
that are valid for all possible experience and in all cases
(like the logical and mathematical propositions); it is by virtue
of the same element, lying at the ground of all experience and
therefore relatively independent of experience. We are determined
only by experience in that we ourselves determine experience.


By Kant's critical philosophy, whose great significance
rests on having brought out this way of regarding
knowledge, those men are strongly influenced who in
% --- p. 89 ---
\setcounter{footnote}{0}%
\opage{89}England stand as the opponents of empiricism. It is on the
other hand interesting to see that the general tendency of English philosophy
toward induction and the real also clearly prevails
among them. They too recognize that experience
and induction are our essential paths of knowledge; they turn
with sharp polemic against the German speculative systems
and assert the relative and the empirical as the proper domain
of our knowledge. But they assert that experience cannot
be its own foundation, but presupposes certain principles
on which its own possibility and validity rest.\footnote[1]{However, the ideas of German speculation too have gained entrance in England. Not only are there Hegelians in Oxford, but already \emph{Coleridge} independently developed several Schellingian ideas, particularly in the aesthetic domain (the concept of imagination). And here a more recent author must be especially mentioned, the well-known historian \emph{Thomas Carlyle} (b. 1795), who in his ``Sartor resartus'' (written 1831, published some years later, first in North America, where Carlyle is much read) has presented, in humorous form, a humane view of life that has sprung from serious and profound thought. His view of the world is an idealism, for which ``all objects are only windows, through which the philosophic eye looks into Infinity itself.'' ``The whole outer universe and all that it includes is only a clothing; and the essence of all science lies in the philosophy of clothes.'' Everything is symbols; but the highest symbol is in the heroic personalities, the models of humanity. --- Stuart Mill states in his Autobiography how much he owes to Carlyle, who often by his intuitive path developed ideas to which Mill came only a long time afterward through induction. Thus one will find, among other things, in the cited work of Carlyle (Book III. chap. IV.) Mill's sociological view in nuce. Nor does Carlyle conceive philosophy merely as speculation; besides the definition cited above he also gives another, which shows the practical significance of theory: ``What is philosophy but a constant struggle against habit, a ceaselessly renewed effort to go beyond (to transcend) the sphere of blind habit and thereby become transcendental?'' This definition is apt against empiricism, according to which all knowledge rests on habit, so that the attempt to get beyond habit is self-contradictory, although Mill has believed that in empiricism he has the only usable means of combating handed-down habits and prejudices. --- Yet for Carlyle it is not thought but imagination that is the highest; ``The understanding is the window, but imagination is the eye that sees through it.'' The philosophical ideas are for him a serving foundation for mysticism and morality.}
A decisive opposition exists between the investigators
% --- p. 90 ---
\setcounter{footnote}{0}%
\opage{90}whom we are to occupy ourselves with in this chapter, and their
opponents. The representatives of the empirical and realistic school
have their view of life in their philosophy. For
Whewell and Hamilton, on the other hand, philosophy is only an outwork;
the real sanctuary, the view of life itself, rests
in both these men on positive religion. This
must naturally also have consequences in a philosophical
respect, especially with regard to consistency and carrying through
of the principles. The philosophical problems are not held fast
and debated in their purity, when another
interest plays in on the part of both assailants and defenders.
Although the struggle has only exceptionally extended into
the religious domain, the feeling must nevertheless often have arisen
that this was what it was really about. On the other
hand this circumstance has given the struggle a heightened
interest, in that the questions of principle, which the English otherwise
prefer to hide under the mass of details, have thereby come out
more clearly. ---

I. \emph{William Whewell} (b. 1795, d. 1866). ---
That investigators other than those who belong to the empirical
school can also ascribe to experience a significance of principle
and have earned merit in respect of it, of this we have an example
in Whewell, who not only himself began as a specialist
in natural science (mineralogist), but whose ``History of the
inductive sciences,'' according to Stuart Mill's own declaration, was
for him an indispensable foundation in the working out of his
Logic, which without this work could not have been carried through.
The study of the natural sciences led Whewell to
investigate their method and the principles and conditions
of their development. These investigations led him more
and more to a conception of the nature of knowledge and of experience
which on the whole agrees with the Kantian, although
he too, not always to the advantage of definiteness and consistency,
is influenced by the Scottish school. We find
also in him this school's inclination to appeal
% --- p. 91 ---
\setcounter{footnote}{0}%
\opage{91}to the common consciousness and to content itself with the
traditional ideas instead of entering upon
more rigorous philosophical analyses.

Having completed ``the History of the Inductive Sciences''
(1837) Whewell published his second chief work ``Philosophy
of the inductive sciences, founded upon their history''
(1840, revised 1847), in which by means of the
historical development of the sciences he seeks to trace the relation
between the elements of knowledge. To this work we shall
in what follows chiefly keep. ---

Not only the natural sciences, but also the sciences of mind
(psychology and ethics) are inductive sciences. But
there is an element in induction that has often been overlooked.
The great discoveries and advances in the history of science
take place not only through the pointing out of new facts; an
essential condition is the ideas and forms by whose help
these facts are conceived and determined; indeed, the very pointing out
of facts often becomes possible only because the investigator is led by
the right way of regarding things. The single facts are
not merely to be brought together, but to the very
mental activity by which they are combined a new element
is added; a general proposition is formed which did not exist
in any of the observed facts in and for themselves. There must
be found an intermediate determination in which the known and the
unknown facts can meet. Thus Kepler conceived
the orbits of the planets as ellipses, and thereby it became possible for him
to set up his laws.\footnote[1]{``Only those,'' says Whewell (Gesch. der induct. Wissensch. übers. von Littrow I. p. 426) on the occasion of Kepler, ``attain the possession of new truths who can make a lively leap from a point within their knowledge to others, which are often far removed from it, and bring about a connection between them, --- not those who cautiously stand still at every place and wait until they are driven from without to go further.'' Cf. \emph{Reuschle}: Kepler und die Astronomie. Frankfurt 1871. p. 141: ``To the question how Kepler found these extraordinary propositions (the three famous laws), the short answer is probably that he calculated them out of Tycho Brahe's observations, since these were so good that they were satisfied only by the \opage{92}elliptical theory, and he had mathematical power and endurance enough not to rest until he had in all respects done justice to those observations. But we must now descend deeper into the workshop of his genius and follow the course of his research from his fundamental thought, in order to see how the history of his objective, physical work is most intimately bound up with the history of his subjective, metaphysical system ...... At the ground of all Kepler's mathematical investigations and metaphysical researches lies the idea, the conception of which alone would be enough to found his greatness, that the solar system, whose rough outer outline Copernicus had given, is a real system only when law and order prevail in it, so that the laws that govern the motions have their cause in definite inner forces which are essential to the nature of the system and its bodies.''}
 In the same way Newton
% --- p. 92 ---
\setcounter{footnote}{0}%
\opage{92}conceived gravity as attraction and this as a
universal cosmic force. Only on the presupposition of the
introduction of such an idea does a discovery become possible. But in
general one conceives induction rather as a means of proof
than as a method of discovery.

Concept and fact stand in a reciprocal relation to
one another. The concept can never be independent of the fact,
but on the other hand the fact must
constantly be drawn toward the concept. Knowledge presupposes
a combination of thought and object. Without
thought there would be no connection between the objects,
without an object there would be no reality. Every
single act of knowledge includes both elements. The difference
between them can only be relative. What from
one point of view is theory is from another point of view fact.
What in relation to a preceding investigation is
theory is in relation to a following fact. That
the planets move in ellipses was with Kepler a theory,
but becomes with Newton a fact. A true theory is
a fact; a fact is a theory with which one is
familiar. Both are only inferences that we draw from the phenomena;
what we call theory is a conscious inference;
what we call fact is an unconscious inference.

Let us in this respect consider sense-perception itself. A coherent image of a sensible
object
% --- p. 93 ---
\setcounter{footnote}{0}%
\opage{93}is formed only when the individual sensations
are combined and comprehended into unity. An
unconscious inference lies at the ground of all our sense-perceptions.
Concepts such as space, distance, solidity, etc. develop
only through the sensation of our muscular movements, and we
then transfer them involuntarily to the objects. We see a
movement take place and infer force as its ground.
We see different plants or animals, but only because we apply
to them our ideas of likeness and unlikeness. Likeness
or unlikeness we see as little as we see force
or space; or rather, we see them only in so far as
we put them into the phenomena. From this it can be seen how senseless
it is to designate the ideas or concepts as
mere transformed sensations of sense. Even the most elementary
sensation cannot come about without a
contribution from the sentient subject. Our conception of nature
is from first to last an interpretation of nature, a translation
of its content into our own language.

The subject's contribution to the knowledge or
interpretation of phenomena consists in a series of ideas or fundamental concepts,
which bring connection and unity into the content of experience.
Not as though these ideas were to lie ready-made
in the mind from the start. They must not be separated from the
conceiving activity itself, but are precisely this activity's
forms and laws; only by reflection on that activity
do we come to know these forms and laws, and only through experience is that activity set in motion. Consciousness
would thus have no ideas and fundamental concepts, if it did not
stand in interaction with the world of space, time and matter; but then we also know no consciousness at all
that does not stand in such a reciprocal relation. We cannot
imagine any stage of human existence where
that activity and its fundamental laws do not already come into
force and are developed through experience. The forms and ideas
operate long before we become conscious of ourselves, and are developed and
modified in many ways by the influence of events and phenomena. How far this modification extends
% --- p. 94 ---
\setcounter{footnote}{0}%
\opage{94}we cannot judge; but it is certainly of essential
significance for the process that goes on in our
consciousness.

The fundamental concepts are acquired in the course of experience, but
they are not learned from experience in the same sense as specific
facts are learned. The single facts are
for us more or less contingent; for a time they are unknown,
at another time known; some know them,
others not; they are known with greater or lesser certainty.
But the fundamental concepts are the laws according to which all conception,
all experience proceeds. It is this their universal
and necessary validity on which their difference
from the concepts of experience rests. We can learn universal
truths, e.g. the law of gravity, from experience. But if
these truths really have universal validity, then
this is due to the ideas which the understanding applies to the
given.

The two elements that interpenetrate in our
knowledge thus condition a fundamental opposition, which
reveals itself under a series of forms: idea and fact,
concept and object, theory and fact, inference
and observation, but which has its sharpest expression in the opposition
between necessary (a priori) and experimental
truth. The former springs from thought, the latter from observation.
We cannot conceive the fundamental concepts in themselves; they find
their expression in certain fundamental propositions and determinations of concepts,
to which thought comes when it reflects on its own
working. Each single science has its original fundamental propositions
or axioms, which spring from the fundamental concept,
the idea, on which it rests. Thus the axioms of mathematics spring from the ideas of space and time, which lie
at the ground of all conception of the given, the axioms of mechanics
spring from the idea of cause, which likewise is an
original law for the working of thought. From the general
axiom that everything has its cause follow the mechanical
axioms that causes are measured by their effects, and that
reaction is equal to action. The laws of
% --- p. 95 ---
\setcounter{footnote}{0}%
\opage{95}motion derive their form from the idea of cause, their
matter from experience. The axioms are naturally abstract
and must be interpreted by study of actual nature; but
they constitute the firm kernel around which the knowledge of
the latter can develop. That the ideas cited are necessary
means that it is impossible for us to think the opposite of
them, and that every objection to their validity precisely
presupposes it. We may err in the explanation of
the particular; but that there must be an explanation, a
cause, is a presupposition we cannot do other than
proceed from. We cannot, e.g., conceive a stone floating
in the air instead of falling to the ground, without having to
presuppose a cause of this deviation from the
general law.

Just as the ideas of space and time lie at the ground
of mathematics and the idea of cause of mechanics, so
is the idea of purposiveness or of purpose necessary
for conceiving the organic phenomena. Nor can this
idea be derived from experience; it develops of necessity
as the only form under which we can conceive
the phenomena of life. Harvey's and Cuvier's discoveries became
possible only through it. It has indeed been held that as
scientific knowledge advances, the domain of final causes
is restricted and the mechanically operating
causes step in. But this is an illusion.
The idea of purpose must now be conceived in another way than
before, but cannot be given up. We do not see purposiveness
in particular single phenomena, but in the whole
connection. We do not suppose that a cloud is sent, so that
it may rain on our field, but we find the idea of purpose
expressed in the connection between the laws of heat, air and
liquids and the development of vegetable life. The idea of purpose
thereby acquires precisely a more grand character:
instead of a few threads that connect separate
objects, we obtain a total connection embracing the whole of reality.

% --- p. 96 ---
\setcounter{footnote}{0}%
\opage{96}Also in the domain of ethics\footnote[1]{\emph{Whewell}: History of Moral philosophy in England. 1852.} we have an originally
operating idea, which comes to light in the axiom that we
ought to do what is right. Even if there is disagreement about
\emph{what} is right, there is nevertheless no dispute \emph{that} the
right is to be done. This germ of the moral life is developed
further through experience in a manner similar to that in which
in the theoretical domain the fundamental concept is developed
through the observations. Men's drives bring them
together and make fixed laws for their life together necessary;
thereby the moral consciousness is educated gradually
also to apply the concept of law to its own life. The
highest law for a man's actions must be one
that comprehends his whole being, thought and drive
as well as action. From the outer, legal fulfillment of duty
man is therefore led to the inner, ethical one, and gradually
gains a higher standard, by which he can
judge the laws of society. --- In opposition to this
ethical conception, which proceeds from an inner principle as
standard and presupposition for the judgment of everything outer,
the independent ethics, the ethics of principles (the morality of
principles), stands the ethical doctrine that proceeds from
the consequences of the action (the morality of consequences), from
utility (utilitarianism), whose most important representatives are
Epicurus in antiquity, Hobbes, Paley and Bentham in more recent times.
This ethical empiricism overlooks, like the theoretical
empiricism, the original presuppositions. ---

One will easily find again in the view presented here
Kant's fundamental thoughts, but one will find them again in a somewhat
vague and indefinite form. Whewell's merit consists in
his having carried through the history of the sciences with special
regard to method and the problem of knowledge, and in his
having asserted the originality and activity of thought in the process of knowledge
against the empirical school. But he
does not give what was found by the historical and critical examination
% --- p. 97 ---
\setcounter{footnote}{0}%
\opage{97}a more penetrating philosophical treatment. He
is for the most part content, for each particular scientific
domain, to infer back to an
a priori idea, but does not investigate whether these ideas
might not be reducible to one or at any rate a few, nor
how they are related to the nature of consciousness,
or how they have developed. This investigation
is all the more necessary, since not only the ideas
or fundamental concepts stand in an opposition to the empirical,
which cannot be deduced from them, but also in a
mutual opposition to one another, the concept of cause (i.e. of mechanical
cause), e.g., to that of purpose. Several of
the ``ideas'' he takes up even almost in quite popular
form, e.g. the idea of purpose. He cites here Cuvier's words:
``Zoology has a principle peculiar to it, which it
applies with great advantage on several occasions, and that is
the principle of the conditions of existence, which is usually also
called the principle of final causes.'' But if one looks more closely,
these words lead precisely beyond the popular idea
of a purpose. For what Cuvier means is really nothing
other than that a living being's structure and organs must
be so constituted that its existence is possible; therefore
the analysis of the conditions of existence (conditions de Texistence)
can often lead to laws that agree with experience.
But there need therefore be no question of a preceding adjusting of the individual to the conditions;
the agreement may quite simply have its ground in the
interaction of living beings with the outer conditions.

Whewell's view is, as is seen from the presentation,
not that those fundamental concepts lie ready from the start in
consciousness. He suggests, as likewise cited, in a single
place, that the forms under which we conceive things already
form before we become conscious of ourselves, so that we
now know them only in the shape to which the events
of the unconscious life have brought them. This suggestion points
in a right direction, as Whewell indeed on the whole maintains
that the ideas spring from the activity of the mind, have no
% --- p. 98 ---
\setcounter{footnote}{0}%
\opage{98}independent being. But all the more striking is it that
he can content himself with a mere and bare enumeration of
them without closer criticism and analysis. In this respect
his opponent Stuart Mill has a great advantage
over him.

While it cannot be denied that Whewell with regard
to the general conception of the nature of knowledge has
seen more correctly than Mill, he cannot nevertheless compare with
the latter in the sharper specification of the course of inductive research.
Mill holds fast to the exact standpoint, demands the
strict logical proof. Thereby he is led to formulate
the special inductive methods. Whewell emphasizes the
mental productivity that comes to light in discoveries.
Purely logically considered Mill is right when he says
that by Whewell's determination of induction we do not get
beyond hypothesis, because exact verification is lacking.
Metaphysically and historically Whewell, logically Mill, has
the right on his side. ---

Whewell does not observe the limits that Kant fixed
for the application of the fundamental concepts (the categories). The ideas
that we must necessarily apply in considering nature
and human life are, according to his doctrine, the only natural
means by which we can form a concept of the
Creator and Governor of the world. Just as we infer that
every single thing must have its cause, just so we must
infer that the whole world must have a cause. That there
must be a first cause, he regards even as an axiom.
We cannot stop at an infinite series of finite
causes; we must posit a beginning; in order that the phenomena
may become intelligible to thought, it must
rise above the endless succession. And the moral
ideas lead us to think that first cause as a spiritual
providential power. --- Whewell has written one of the so-called
Bridgewater Treatises,\footnote[1]{So called after an Earl of Bridgewater, deceased 1829, who set aside a sum for the publication of such writings.} which aimed ``by way of natural
% --- p. 99 ---
\setcounter{footnote}{0}%
\opage{99}science'' to prove the Creator's wisdom and goodness
as revealed in the world. His uncritical confounding
of thought and experience and the resulting obscurity with
regard to the rational character of knowledge comes to light
especially in his theological ideas and their
connection with the philosophical concepts. As easily
as it comes to him to go to the idea of a first
cause, so easily does it come to him to connect the doctrine of the Bible
and of science. The moral ideas are not only to
lead us to belief in a providence; they are also to ground
the consequence that, when the moral ideal is not fulfilled,
supernatural arrangements must have been made
to realize it. To these arrangements belong
especially the Bible itself, whose narrative of the creation of the world
science, which does not reach up to the beginning of time,
can neither confirm nor reject. --- Here too we find
an indefiniteness in the concepts, a lack of more penetrating
criticism and of consciousness of the decisive differences, which diminishes Whewell's merits in the problem of knowledge.


II. \emph{William Hamilton} (b. 1788, d. 1856). ---
This name leads us to Scotland, from which the first reaction
within Britain itself against empiricism proceeded at the end
of the last century. Thomas Reid's doctrine has since
that time been dominant at the Scottish universities.
It seems to agree with the national character. For its background
it has a firm and earnest moral and religious
conviction, as one may expect in the fatherland of John
Knox; it keeps sober and calm within
the limits of ordinary consciousness and seeks merely to
present this consciousness's ``facts;'' but it
nevertheless lays throughout the weight on experience and wants to
apply the principles of Bacon and Newton, only that it does not with
Locke and Condillac restrict experience to the sensible.
It seeks to find the primitive and
% --- p. 100 ---
\setcounter{footnote}{0}%
\opage{100}necessary statements of consciousness, which lie at the ground of all other experience.
Scottish philosophy is essentially empirical psychology,
only that it seeks to point out facts of consciousness
which ordinary experience, in its opinion, overlooks.
It stands on the ground of empiricism, although it wants to dispute
its consequences. It believes that it can find immediately
primitive facts and stops at a series of
principles which it regards as the last that can be reached.
Thereby it becomes a kind of ``systematic credulity'' and
sins against the true empirical method, which first recognizes
a fact when scientific investigation establishes
it, just as much as against the principle of philosophical analysis,
according to which one fact cannot directly explain
another, but both must be explained by a higher law.

The Scottish school thus points in two directions.
The chief weight may be laid on the empirical tendency, and
we then come, by a sharpening of the scientific method,
to the standpoint of modern English psychology. On
this path Thomas Brown (d. 1820) already struck out, just
as Alexander Bain is connected through this with his
Scottish predecessors (see above p. 66). But one could
also bring forward the critical-speculative element, and thereby
approach Kant. On this path Hamilton has
brought the Scottish school beyond its vague and unscientific
attitude. He brought to this task a deep
and many-sided knowledge of Greek and German philosophy,
altogether an astonishing learning in the history of philosophy,
rare in England, which contributes not a little
to the great interest and instruction his writings afford. But
he was besides a keen logician and dialectician, an investigator
who found the aim and yield of philosophical thinking in
the very activity of thought and in the movement into which the highest
mental powers are thereby set. He has not written
much, but worked for many years as professor of philosophy
in Edinburgh. He himself published only some dissertations, which
he added to his edition of Reid, and a collection of his
articles in the Edingburgh Review under the title ``Discussions
% --- p. 101 ---
\setcounter{footnote}{0}%
\opage{101}on philosophy'' (1852); after his death his lectures
on metaphysics and logic were published by Mansel and
Veitch in four volumes.

The first work that made Hamilton famous was
an essay in the ``Edingburgh Review'' 1829, directed against
Schelling's and Cousin's doctrine of the Absolute. It is an
attempt to show that there can be no philosophy
of the Unconditioned. We can know only the conditioned
and the limited. The infinite, i.e. that which is unconditionally
without limit, and the Absolute, i.e. that which is unconditionally
closed in itself, we cannot grasp, cannot construct in
our thought. Our thought is from first to last conditioned, indeed,
to think is to make conditioned (to think is to condition).
Thought cannot transcend consciousness, any more than
the eagle can fly beyond the limits of the atmosphere; and consciousness
has its definite conditions, to which every content of consciousness
is subject. We conceive existence in time
and space, not in and for itself. But time and space presuppose
relations between different existing elements; what exists in time
and space is therefore relative, and time and space
cannot be applied to the Unconditioned. Consciousness
presupposes furthermore an opposition between a knower
and a known, between subject and object; but this opposition
cannot be applied to the Unconditioned, which thereby
would become relative. The only way out here would be, with
Schelling, to assume that the Unconditioned is raised above the opposition
between subject and object, is the indifference of both,
and that in order to know this indifference we must raise ourselves
above consciousness and reflection, which always presuppose
plurality and relativity. But how can we raise ourselves
above our own consciousness? And with what right can we
posit a faculty that conflicts with the very laws
of our consciousness?

One sees how Kant's doctrine is taken up by Hamilton
and plays a leading part in his view. Yet he parts
from Kant at a decisive point. Kant had indeed restricted
the applicability of our forms of intuition and of
% --- p. 102 ---
\setcounter{footnote}{0}%
\opage{102}thought to the phenomenal and relative, but he nevertheless
recognized the Unconditioned as a principle of reason, which indeed could not
be laid at the ground as a constructive principle, but which
has regulative significance, serves to lead beyond the
relatively given and to set ideal tasks. In this Hamilton finds
the germ of the systems of German speculation built upon the idea of the
Unconditioned. The fundamental error is, in his
opinion, that the Unconditioned is assumed to be a positive
concept. Only the conditioned can we positively know and
comprehend, embrace and represent with our thought; the Unconditioned,
the negation of the conditioned, is unknowable and
inconceivable, is a purely negative concept, like the Unthinkable.
If one now would say with Kant that we nevertheless possess
the Unconditioned as a positive, regulative principle, then
a self-contradiction would be admitted in our highest mental
faculties, which would necessarily lead us into skepticism.
The solution is on the contrary quite simple, if one
holds fast that we can conceive only the conditioned, and that
the Unconditioned is for us a pure negation. The conditioned
lies always between two extremes, which have the character of
the Unconditioned; we cannot grasp either of them, and yet
one of them must be the true. The world must e.g. at some time
have a beginning or be eternal; but we can
no more grasp an absolute beginning than an infinite
eternity. We learn from this that we cannot make thought
the standard of existence. We know only conditioned
existence, not existence in and for itself. What lies beyond
the conditioned is an object of belief. There must be
an existence in and for itself; otherwise there could be no conditioned
existence. ``And by a wonderful revelation there is
instilled in us, by the very consciousness of our
lack of power to grasp anything that lies beyond the
relative and finite, a belief in something Unconditioned beyond
all the reality that we can grasp.''

If philosophy were to have the Unconditioned as
its object, it would be impossible. An absolute truth
is inaccessible to us. But on such a truth the value
% --- p. 103 ---
\setcounter{footnote}{0}%
\opage{103}and significance of philosophy do not rest either. It is not finished results
that it is after. A complete solution of
the problems, which would stop all further discussion, would be
a paralysis of philosophical striving and lead us back
to an intellectual barbarism. Our being and our happiness
consist in activity, and knowledge is of value only
in so far as it is activity and striving, not dead possession.
The significance and necessity of speculation must be sought mainly
in the exercise of our highest mental powers that it
brings about. Lessing is right: for us the eternal striving
after truth is better than the possession of truth. A
waking error is better than a sleeping truth. The
highest to which philosophy comes is the recognition of our
ignorance with regard to existence in and for itself.
But there are two kinds of ignorance. We philosophize in order to
escape the blind, thoughtless ignorance, which does not know
the conditions of consciousness and of truth, and when we have
come to know these conditions, we end in the higher, the conscious
ignorance (docta ignorantia).

Hamilton thus posits an opposition between knowledge
and belief. Only the latter reaches the Unconditioned. But
what motives has it then to seek beyond the finite
and to assert the positive reality of the Unconditioned instead of
leaving it standing as a mere possibility, as a ``Ding
an sich?'' These motives lie in the nature of man as
a moral personality. Man is conscious of an unconditioned
law of duty, and moral responsibility presupposes
in turn freedom of the will. The nature of man has its peculiar
expression in the struggle between moral freedom and
physical necessity. And only as a free moral
being can he believe in a wise and holy ruler of the world.
Admittedly an unconditioned freedom is inconceivable, but we know indeed
that the limits of thought are not the limits of possibility and reality.
By the recognition of this philosophy paves the way for the recognition of responsibility and freedom,
which Christian revelation also presupposes. The end of
philosophy is the beginning of theology.

% --- p. 104 ---
\setcounter{footnote}{0}%
\opage{104}We have here reached the core of Hamilton's view of life.
The relativity of knowledge gives place here to
the moral and religious belief in the Unconditioned. The
mental nature of man, as it is established by
moral responsibility and freedom, was for him the
highest, just as for Kant. His motto was: ``On
earth there is nothing great but man, and in man
there is nothing great but mind.'' Stuart Mill opposes
this motto, which has a decided spiritualistic
tendency. In this core the opposition appears clearly
between the two philosophers, the representatives of the two warring
schools. Mill holds fast, like Hamilton,
that we know only the relative and conditioned, but holds
with right that we must then also conceive the nature of man
as conditioned, that the method that leads us to insight into
the laws of nature must also teach us to know the laws of human life,
since human life itself is a link in the life of nature.
Mill's strength lies precisely at this point; his
conception of the concept of cause helps him beyond the abstractions of the doctrine of freedom and
necessity. He has also clearly seen that Hamilton's philosophy is determined in its whole
peculiarity by the idea of unconditioned freedom.
For Hamilton Mill's doctrine had to stand as fatalism, and
fatalism led, in his conviction, necessarily to
materialism and atheism. The sciences that especially
rest on the concept of necessity, mathematics and physics,
he was therefore not free to regard with other than suspicious
eyes. This remoteness of his from the natural sciences places
him in opposition to the striving of the whole most recent philosophy,
which goes precisely in the direction of a closer connection with
them.

Mill was, as we have seen, more fortunate with the projectiles
he cast into the fortress itself than with his attacks
on the outworks. Hamilton's merit rests on the
clarity and originality with which he has conceived Kant's
doctrine of the relativity of knowledge, and Mansel has shown
that Mill, as far as the main points in his attack on this
% --- p. 105 ---
\setcounter{footnote}{0}%
\opage{105}doctrine are concerned, has not understood Hamilton. Moreover
Mill, in accordance with his empiricist tendency, could not attack Hamilton
at the really decisive point. This point we find in
the relation between the Unconditioned and the conditioned,
from which Hamilton proceeds. The Unconditioned he determines
as the negation of the conditioned. If one proceeds
from this determination, one need not first show
the impossibility of conceiving the Unconditioned; for then
this concept is from the first self-contradictory. For
if the Unconditioned has the conditioned outside itself as
its opposite, then it thereby becomes itself conditioned and relative,
and is thus canceled at the very moment it is posited. But
what right has one to proceed from this presupposition
of the opposition between the Unconditioned and the conditioned?
It is a presupposition which certainly lies in the common
consciousness and is contained in the traditional religious
ideas. But the first step of philosophical thought
must be to rise above this dualism.\footnote[1]{It appears in the history of thought in a twofold form. For the Greeks the finite and limited was the perfect, the infinite the lower and in the last instance the evil. Even Origen therefore conceived God as limited. In the Middle Ages (cf. Augustine and Anselm) the relation was reversed; God became the Infinite, eternally raised above the world of finitude. Immediately at the beginning of modern times, on the other hand, profound thinkers (Cusa, Bruno, Spinoza) sought the inner unity of the infinite and the finite. Cf. \emph{Weisse}: Das philosophische Problem der Gegenwart p. 195 f. R. \emph{Nielsen}: Mathematik og Dialektik. p. 117 f. --- Hamilton has with his knowledge of the history of philosophy clearly seen this; he finds precisely already in Cusa the idea of the Absolute of the more recent speculative philosophy, against which he polemicizes.}

The true Unconditioned must be that which receives its
closer determination from the conditioned, --- that in whose
nature all conditions have their ultimate ground; the true
Absolute must stand in an inner relation to the relative. Admittedly
an immediate and absolute knowledge of the Unconditioned is therefore
impossible; so far Hamilton is right. But we can
through the knowledge of the conditioned gain decisive
% --- p. 106 ---
\setcounter{footnote}{0}%
\opage{106}determinations of the essence of the Unconditioned, by virtue of the inner
relation between them. Here lies a series of problems
with which English philosophy has occupied itself only little,
but which we find epoch-making treatments of in the great systems on the Continent, especially
in Spinoza and Hegel,
problems that inevitably press forward again, however often
they are pushed aside.

Hamilton himself cannot carry through the opposition
between existence in itself and the phenomena, the Unconditioned
and the conditioned. He states (Lectures I p. 125)
that from the peculiar nature of the phenomena we can
draw many valid conclusions concerning the unknown
principle whose revelations they are. Although therefore existence
is revealed to us only in the phenomena, and we can only
have a relative knowledge of mind and matter, we
can nevertheless by inference and analogy rise above the
mere appearances which experience and observation show
us. But how could this be possible, if there
were not an inner relation between existence in itself and
the phenomena?

According to Hamilton the end of philosophy is the beginning of theology.
But the question now becomes, with what
means theology is to build its doctrine. If we look back to
the history of theology, we find that it has constantly made use of
philosophical concepts or at least forms of concepts.
But when now these concepts, when all our thought and knowledge
concern only the conditioned, how then does
a theological doctrine of the Unconditioned become possible? Hamilton
says himself that no difficulty arises in theology which
does not already show itself in philosophy. But then one must
also conversely say that the difficulties that arise in
philosophy must also hold for theology. It is then
inconsistent when Hamilton in one place derives the imperfection and insignificance
of English theology from its unphilosophical
character. Consistent it is on the other hand when he
insists that the statements of the Bible are to hold without reinterpretation,
that the theologian's chief task is interpretation of the Bible, and
% --- p. 107 ---
\setcounter{footnote}{0}%
\opage{107}philological learning therefore the only necessary condition
for theology. Consistent it is also when his
gifted disciple \emph{Henry Mansel} has derived from his principles
the impossibility of a scientific theology.

On the so-called Bampton Bequest lectures are held at Oxford
annually in defense of the Christian doctrines.
Such ``Bampton-lectures'' Mansel too has held, and
he has here ``used his short appearance in a theological
chair to prove that there is no theology.''\footnote[1]{\emph{Rémusat}: La théologie naturelle en France et en Angleterre. p. 75.}
He first declares himself against dogmatism as well as against
rationalism. Dogmatism wants to deduce scientifically
formed dogmas from religious belief; rationalism wants to
reduce the religious ideas to symbolic expressions of
abstract truths. Both aim at attaining a higher
knowledge of the divine things. One thinks that when
one abolishes the figurative and anthropomorphic, one will
attain a truer knowledge of God. But the abstract
forms that are put in place of the so-called figurative are
conditioned by the nature of our consciousness and our knowledge,
which is finite. We are not able to conceive the Unconditioned
as a personal being, which we nevertheless must do for moral and religious
reasons. Our personality is finite and conditioned
and can therefore not be the form for the conception of the
Unconditioned. Every attempt to develop a rational
concept of God leads therefore, as the history of philosophy shows,
to pantheism, which is incompatible with moral and
religious belief. Therefore no scientific theology can
be given. But since thought does not constitute the whole nature of man,
since he as a moral and religious personality finds himself in
a relation of dependence and obligation to a highest
personal being, the object of his worship and the source
of the law of his action,\footnote[2]{``To use the name of God in abstract speculation merely as a conventional designation for the highest link in the chain of thought, and \opage{108}to believe in him for practical reasons in reverence and obedience, are two very different things; and only for the latter, not for the former, is the concept of God as a person necessary. If man were a pure being of the understanding, the problem of the Unconditioned would lose its chief difficulty; but he is also a being with religious and moral faculties, and these too demand to be satisfied by a valid solution of the problem.'' \emph{Mansel}: The philos. of the cond. p. 13.} the religious
% --- p. 108 ---
\setcounter{footnote}{0}%
\opage{108}doctrines must turn to all sides of man, appeal
to his feeling and will, indeed apply analogies and inferences of probability
from the finite, but not ground their
validity exclusively on the recognition of thought. We
must thus sharply separate knowledge and belief. We cannot
comprehend how the Unconditioned and the personal are
united in the divine being; but we must believe that
they are somehow united therein. Our thought
can therefore not be taken as the standard for religious belief;
we cannot form any scientific concept of God,
according to which we should be able to judge the religious doctrines.
--- This is Mansel's doctrine of ``the limits of religious
thought,'' which has aroused so much attention in England
in recent years.

It is a sign of the times that attacks are made on theology
from so many different sides in our days.
What has now been described, which proceeded from the very old bulwark of English orthodoxy,
the University of Oxford, has this in common
with R. Nielsen's polemic at home, that it does not proceed from
a hostile relation to positive religion, but at the same
time wants to assert the peculiarity of religious belief and
the laws of scientific knowledge. Yet there is one
essential difference between Mansel's and Nielsen's standpoint,
and in one point, where one must admit that the English
philosopher is more consistent than the Danish. Nielsen
maintains, what Mansel denies, a philosophy of the Unconditioned.
Having proceeded from Hegel's absolute idealism he holds fast
also in his later doctrine to the possibility of a philosophical
determination of the nature of the absolute spirit, a possibility
% --- p. 109 ---
\setcounter{footnote}{0}%
\opage{109}which for him becomes one with the possibility of knowledge itself
(Grundideernes Logik I p. 53). Here then arises the
difficulty of determining the relation between this philosophical
idea of God and the conception of the Deity which the
demands of the religious life shall lead to; there cannot after all be two
absolute ideas, and the strife between philosophy and religion
begins then anew with the question which of them
is to decide what the \emph{true} Unconditioned is. This
Mansel, and already Hamilton, has clearly seen. They deny that
philosophy can know the Absolute and Unconditioned, and
thereby the place stands open for the religious ideas.
Mansel remarks that the Greek philosophers could have a
philosophical idea of the Absolute, of the principle lying at the ground of
everything, without having to make it one
with the deity of the popular religion; they could regard the popular religion's
gods as derived personal forms of that
Absolute, just as in the Theogony the higher gods have developed
out of Chaos and Night. But within Christianity the absolute principle that one assumes must necessarily be
made one with the God in whom one believes.

As much as one must grant Mansel right in this, it
is nevertheless a question whether the opposition between belief
and knowledge can be maintained, even when one conceives it
as he does. Every positive religion will after all
precisely teach what God and the divine is, and in particular
the Christian church sets up a series of definite doctrines,
which it demands be recognized; in the development of these doctrines
it is, as the history of dogma also shows, impossible
to keep thought and knowledge out. Belief can thus
not be without theology, and it will therefore prove to be
an illusion to think that one can overthrow all theology and
nevertheless maintain positive belief.

It is only when one makes a dualistic distinction
between the Unconditioned and the conditioned that one arrives
at such an opposition between knowledge and belief. If one wants to
maintain the supernatural, then it is quite consistent
to recognize the impossibility of a grounding by way of
% --- p. 110 ---
\setcounter{footnote}{0}%
\opage{110}thought. But the supernatural principle, which includes in itself
a denial of the validity of thought, is as such
foreign to philosophy. This can only conceive
the Absolute as immanent, in its inner unity with the
relative, and it also sees in the moral and
religious nature of man no necessity of leaving this
rational, human standpoint, since we have no other
means to \emph{know} true morality and true
religiosity than precisely thought, knowledge itself.
If there were, as Mansel declares, ``legitimate
sources of conviction, different from those that constitute
the knowledge specially so called,'' then this
conviction would have to develop according to a logic that is different
from the one we know and apply, so that we got two kinds of
logic --- and the old strife would then begin again from the start. ---

We now return to Hamilton. Stuart Mill's criticism
of Hamilton's philosophy aimed partly at combating the attempt
to assert a belief independent of knowledge; --- and on
this point at least Mill was not exactly fortunate, although
he fought for the cause of philosophy; it concerned
problems which he and his school have always pushed
aside, and it was therefore a foreign domain that he here
ventured into. Partly it was directed against Hamilton's
logical and psychological doctrine, which we shall now characterize.


Logic is, according to Hamilton, a purely formal and a priori
science. It is to take no account whatever of actual
existence and its relations, but occupies itself only
with the laws and conditions of thought itself. Concerning the
truth or falsity of propositions in themselves logic decides nothing;
for it everything can be true that does not contain a
self-contradiction. Logic has thus to do only with the
necessity that rests on the mutual relations of concepts,
but can decide nothing concerning the real necessity
that springs from the actual nature of things. Hamilton
even reproaches Aristotle with having violated the
% --- p. 111 ---
\setcounter{footnote}{0}%
\opage{111}formal and abstract character of logic by mixing in considerations of
the real connection.

Within the domain of formal logic Hamilton
has displayed much zeal and ability. He has attempted
to set up a new theory of inferences, which we cannot however
go into more closely here.\footnote[1]{In \emph{Trendelenburg}: Logische Untersuchungen. 2te Ausg. II p. 304---307 is found a brief characterization and apt criticism of this ``New analytic of logical forms.'' In particular it is shown how Hamilton's own theory leads beyond formal logic.} But with regard
to formal logic as a whole we must remark that
Mill is right when he maintains that it may well serve
to test whether our inferences are not self-contradictory,
but not to show how we gain new knowledge.
The extension that Mill has given to logic is
therefore a real scientific progress. Formal
logic tests the agreement of thought with itself, but
does not solve the most important task, the agreement of thought
with reality, the proper problem of knowledge.

Hamilton shares his general philosophical starting point
with the earlier Scottish philosophy, in that
he regards philosophy as the scientific presentation
and development of the facts of consciousness.\footnote[2]{Philosophy in the proper sense thus has to do only with the mental, not at all with the material; cf. Hamilton's motto. Philosophy becomes only ``philosophy of mind;'' ``natural philosophy'' is not what we understand by philosophy.}
The original statements of consciousness, our instinctive assumptions,
every philosophy must proceed from that does not wish to
make itself nothing. What the Bible is to the theologian,
that consciousness is to the philosopher. That those original
statements are true, one must assume, ``if the very root
of our nature is not a lie.'' And the most important of these statements and original facts is the opposition between consciousness
and its object, between self and not-self.
Idealism, which cancels the not-self and makes it a
% --- p. 112 ---
\setcounter{footnote}{0}%
\opage{112}merely subjective phenomenon, dissolves everything into dream and shadows.


This is the same argumentation that Reid applied
against Hume. But in two points Hamilton departs
from Reid. He does not regard consciousness, like the latter,
as a special faculty, different from observation and attention
, but as one with immediate
knowledge itself. We know an object immediately, i.e. we
are conscious of its existence. We know ourselves
immediately, i.e. we are conscious of our own existence. Hamilton
next sets as a necessary presupposition that
the primitive statements of consciousness are not self-contradictory.
Only under this presupposition do these primitive statements,
the principles of common human understanding, have validity
above all derived phenomena of consciousness, whose validity
rests on those.

Consciousness itself cannot be defined more closely, since it
lies at the root of all knowledge. All that we can
imagine and comprehend, we imagine by the help of
consciousness. This is constantly presupposed. If it now itself
were to be capable of being illuminated and comprehended, then we should indeed have
to have another consciousness, which could conceive the first, and
so on to infinity. But although consciousness cannot
be defined, it can nevertheless be analyzed. We can hold fast
and compare the phenomena of consciousness and thereby
infer the general conditions of the acts of consciousness.

Philosophy (as philosophy of mind) falls into three
main parts. In the first, phenomenology or empirical
psychology, the phenomena of consciousness are presented as
they actually are found. They arrange themselves in three groups:
knowledge, feeling and will. The second part, nomology
or rational psychology, presents the laws that
hold for these phenomena, and divides therefore into logic,
aesthetics and ethics. The third part, ontology or ``inferential
psychology,'' finally draws from this results
with regard to the nature of mind in itself; here are treated in particular
the problems of the existence of God and the soul's
% --- p. 113 ---
\setcounter{footnote}{0}%
\opage{113}immortality. --- Hamilton has managed to treat only phenomenology
and logic. ---

Properly this view must founder on the same
reefs as empiricism. It proceeds from empirical
psychology and can only by an inconsistency get beyond
this empirical foundation and set up laws for the true,
the beautiful and the good. In any case the whole view would
be purely subjective, since it is built only on
the subjective phenomena of consciousness, --- if it did not
begin with the postulated presupposition that these
phenomena have objective validity, since ``the root of our
nature cannot be a lie.''

Yet on closer consideration it appears that Hamilton
really has another starting point than that empirical-psychological one,
that like Kant he proceeds from the
original activity of consciousness. On this rests the
critical-speculative element in him, whereby he offers
interesting correctives to the doctrines of the empirical school.

All consciousness presupposes, he says, a duality,
an opposition. Therefore all our knowledge is conditioned.
But opposition and duality cannot be conceived except
by comparison. Therefore there lies in every act of
knowledge a distinguishing and judging. The comparing activity
plays a part in the simplest observation. The
first act of consciousness is the affirmation that something exists;
but this presupposes a separation between its existence and
its non-existence; thus a comparison of these two
states. Next, a distinction is made between the self and the external
object, between the different inner and outer phenomena
etc. From the simplest to the highest and most composite
states of consciousness we trace the comparing
activity as a pervasive condition. Under this
working consciousness follows certain laws, which indicate the conditions
without which it cannot function. These laws
are therefore the fundamental laws of our knowledge (the logical and mathematical
fundamental propositions); they are a priori, in their principle
independent of experience, which follows from their
% --- p. 114 ---
\setcounter{footnote}{0}%
\opage{114}necessity; a real necessity experience cannot teach
us. Only by that activity of consciousness and the laws it
follows does knowledge come to be. Though knowledge therefore has
its beginning in experience, it nevertheless has its proper
origin in the mind itself.

Since Hamilton has such a conception of the nature of knowledge, he cannot with the English psychologists
regard the laws of association of ideas as the only ones that hold
in the domain of mental life. Nor can he conceive them
as expressions of purely passive occurrences within consciousness,
but must seek to trace them back to the energy of consciousness as their ultimate ground. As such an energy manifests itself
already in the perception and observation of
objects, so it is also due to the energy of the self that
the idea of a certain impression is held fast and preserved
so that it can later be awakened again. And this same
energy works also in the association of ideas, whose general law
may be designated as the law of completion (the law of
redintegration): Ideas that belong together, i.e. constitute
a whole for thought, either because they have earlier touched
one another as successive or simultaneous, or because they
resemble one another, have a tendency to call up
one another. This law, which includes in itself the more special
laws of the association of ideas, seems to have its ground in the
unity of the active self, which also strives to bring about unity and
connection among the ideas themselves. --- In
this way, as already shown above (p. 37 f.),
association psychology, and with it empiricism, unravels itself and
comes under a higher principle.

Association psychology believed it could especially explain
the principle of causality from experience by regarding it
as an example of an association of ideas which, on account of the
complete confirmation it has won in experience, is
fixed and indissoluble. But association of ideas can, Hamilton
objects, only explain how we come to refer
a certain phenomenon to a certain cause, but not how
we come at all to ask for a cause.
% --- p. 115 ---
\setcounter{footnote}{0}%
\opage{115}Most philosophers who do not adhere to pure empiricism
are accustomed to set up causality as a special principle, a
peculiar law for our knowledge. According to Hamilton this is
not necessary. The law of causality is according to his conception
only a special form of the law of the conditioned, i.e.
the law according to which our knowledge comprehends only the
conditioned and relative. Since we cannot comprehend an Absolute
or construct it in our thought, we can also
not think an absolute beginning or an absolute end.
When a thing arises, this is for us only apparent;
we involuntarily presuppose that it has existed beforehand
in another form, i.e. we derive it from a certain cause,
or from causes, for there are always several. We conceive not
existence absolutely or in itself, but existence under
the form of time. From existence as a whole not
an atom can be taken away, any more than an atom can be added
to it; we can comprehend neither how the sum of
what exists could be increased or diminished. Ex nihilo
nihil. Omnia mutantur, nihil interit. When anything begins
or ends, it can only be as phenomenon. Even when
we attempt to think the creation of the world, we always think the
world as contained potentially in God's nature, so
that creation is only a transition from possibility to actuality.
--- It is thus in accordance with the limitation and
finitude of our mind, in accordance with the relativity of our knowledge, that we are
compelled to think causes for all things.

This conception of the principle of causality has the merit
that it strives to trace this principle back to
the very nature of our consciousness instead of setting it as
a separate isolated principle. It stresses especially the side of
the causal relation that this indicates a connection between,
a unity of the phenomena. Now Hamilton regards it
indeed as a consequence of our limitation that we cannot
do other than refer every event to an earlier one,
from which it has developed; but when this referring
shall mean that we conceive the phenomena as one in their
ground and only different in their appearance, then lies
% --- p. 116 ---
\setcounter{footnote}{0}%
\opage{116}it just as near to regard it as a perfection
that our consciousness seeks to return to unity and does not
remain at the phenomenal and accidental plurality.  In its
consistency this doctrine leads, as has rightly been remarked, to the idea of an eternal substance, whose manifestations in
time the phenomenal causes and effects are, ---
that is, to the pantheism that Hamilton feared so greatly.
Hamilton's doctrine here forms a complete contrast to
empiricism, which sees in the relation of cause only a relation of succession,
a connection between different phenomena
without inner, necessary coherence.  Kant has already
shown that our judgments about things have objective validity only,
concern real relations only and not dream-images,
when the relation of cause is more than accidental temporal sequence.
The law of causality is an expression of one of the forms under which
our thought comprehends the given into unity.  Such a
unity we must, in virtue of our organization, presuppose in existence;
we must in the last instance assume an eternal substance
as the ground of all necessary coherence among the phenomena.
It is then clear that Hamilton here, like Kant,
himself points beyond the philosophy of the conditioned.  The unity
and coherence that we find in existence rests on this, that the
Unconditioned is the inner ground of the conditioned, is that which
conditions all that is conditioned.  --- And here we have a natural
transition to the thinker who has developed Hamilton's doctrine as well as empiricism
further in just the manner now indicated, who seeks in an original way to go beyond the
imperfections of both, and who may therefore be said to mark
the high point of the philosophical endeavors in England in
our time.

\chapterhead{Chapter Four.}{The Philosophy of Evolution.}


It follows from the empirical character of English philosophy
that it has seldom striven to gather its
% --- p. 117 ---
\setcounter{footnote}{0}%
\opage{117}teachings into a closed system.  It has always had the
keenest eye for the finite and conditioned in human
knowledge, and yet every system seems to presuppose
an unconditioned principle that can bring conclusion and unity
to the conception of the world.  The tendencies with which we in the foregoing
have been occupied make no exception to this.
Stuart Mill has his strength in the treatment of the single,
isolated problems; these stand side by side in his work and
are not brought into inner connection, although it is naturally
the same fundamental view that we meet in his treatment
of each of them.  Mill has no inner principle; his principle
is indeed that there is no principle, but that all outer and inner
phenomena group themselves according to laws that possess no
necessity valid in itself.  William Hamilton could likewise not
strive for a system, since his chief teaching indeed
was that the Unconditioned is wholly inconceivable to us; and although
he regarded necessity as more than a mere outer
connection in space and time, it was yet in his view
only a consequence of our limitation that we apply the concept of cause
to things; with regard to the Unconditioned
this concept had no validity.

What gives \emph{Herbert Spencer} his unique significance
in English philosophy is therefore not least
his systematic endeavor, his endeavor everywhere to
gather the scattered single features into a comprehensive view.
He is not content to treat the single problems.
Although he possesses an astonishing knowledge in the domain of
almost all the special sciences, he regards it
yet merely as the material from which a new, all-embracing
structure is to be raised.  He is a rare example of a
union of analytic and synthetic gifts.  With the other
English philosophers analysis is the predominant element; what is
given and at hand is tested and resolved into its elements,
traced back to its simplest forms.  They do not, on the other hand,
strive for such a reconstruction of the given that its
development might come to reveal its nature in its
whole coherence.  From the simplest forms, the invisible
% --- p. 118 ---
\setcounter{footnote}{0}%
\opage{118}germs, the synthetic thinker seeks step by step
to derive the highest and most significant functions, in that
he grasps the law valid for the whole of evolution.

It is this at once comprehensive and profound
view that makes Spencer, more than the other thinkers who have proceeded from the older
English school, have an eye for the turning point
to which empiricism in our day has come.  He
is not content to return to the presuppositions
and principles by which empiricism abolishes
itself, but consciously draws their consequences,
sees them in their significance as the foundation of the whole view of the world.
It is the influence of Hamilton and the
critical school that lets him look more sharply at the metaphysical
domain than the empiricists are able to.  The more profound
philosophical grounding, which can be won only by
tracing the laws of thought and of reality back to a
principle that underlies both, this grounding,
which empiricism could neither seek nor find, Spencer endeavors
to give.

But despite this metaphysical tendency Spencer
will not leave the domain of experience.  He still counts himself
among the philosophers of experience, although he polemicizes against pure
empiricism, which in the end must end in the want of principle.
He differs from the empiricists not only in the more profound
grounding that he seeks, but also in the new light
that on the foundation thus won he casts on the factual
determinations in which empiricism rests content.  We
have seen earlier that this halting at the positive and
factual is precisely what brings contradiction into the
empirical view, in that the single determinations stand
fixed and immovable over against one another, --- that at
some points an attempt is indeed made to bring movement and development
in, but that the atomistic character of empiricism
hinders this endeavor, and that movement and development in
any case are not laid down as the governing standpoint.
Otherwise with Spencer.  The fixed forms melt,
life and movement arise in the dead masses.  Evolution
% --- p. 119 ---
\setcounter{footnote}{0}%
\opage{119}becomes not an isolated fact that one comes upon at one place or another, but the expression of a world-law that rules
in all phenomena.  And hereby the way is paved for a kind of
dialectic of positivity, which overcomes and reconciles the oppositions
by ordering them as members in the infinite
movement.

. Spencer's endeavors have in recent years won
deserved recognition.  This is seen not only from the wide diffusion of his
writings in England and North America and from
the translations into foreign languages (French and Russian),
but also from the pronouncements of eminent men of science.
George Lewes raises in his History of Philosophy the
question whether there has ever arisen a more eminent
thinker in England.  And although Stuart Mill, on
account of Spencer's criticism of empiricism, always mentions him
with a certain reserve, he declares him to be Auguste
Comte's equal, no small praise in his mouth.  The
famous naturalist Huxley says of Spencer's chief work:
``The only complete and methodical presentation
of the theory of evolution that I know is found in Herbert
Spencer's `philosophical System,' a work that ought to be carefully studied by all those who wish to make themselves acquainted
with the tendencies of the scientific movement.''\footnote[1]{These pronouncements are found cited in \emph{Ribot}: ``La psychologie Anglaise contemporaine,'' who himself also gives a striking and sympathetic characterization of Spencer (p. 145---152).} --- But
it is also only now, when his great work is approaching its
completion, that Spencer has come to be recognized.  His endeavor
was foreign to the English mental tendency; he is open
and resolute in his line of thought and in his pronouncements, and shows
not the empiricist circumspection that lets even men
like Stuart Mill and Darwin accommodate themselves to the prevailing
ideas; --- therefore he could not count on
sympathy, but had to feel himself out of place in English
society.  ``Amid the general indifference
Spencer has remained unshakably true to philosophy.  He
% --- p. 120 ---
\setcounter{footnote}{0}%
\opage{120}must surely have possessed a heroic courage and a rare
feeling of independence to be able to devote himself to rigorous
studies that could procure him only isolated and unnoticed
expressions of approval.  With the talent, the mental
fruitfulness and the almost all-embracing manifoldness
of knowledge to which his writings bear witness, he would, had
he consented to follow the beaten paths, have
easily obtained all the gifts that English
society bestows in rich measure on those who serve it in the way
that it itself wishes to be served.  He has preferred to dedicate
himself to poverty and to --- what is still harder
to bear --- obscurity.  But Spencer deserves more than
empty declarations of sympathy'' etc.\footnote[1]{Aug. \emph{Laugel}. Revue de deux mondes. 15. févr. 1864.} In
character and view of the world Spencer seems to have much
in common with Spinoza. ---

To give at once an idea of Spencer's teaching
we shall bring forward the main thought that dominates it from
beginning to end, the idea of evolution.  This idea is for Spencer
the result of a careful analysis of the facts of experience;
it states a law that rules everywhere in
the world of nature and of mind.  And on the other hand
only by laying this idea at the foundation can one synthetically
present the whole content of our knowledge.  Spencer never
tires of presenting his idea from both these standpoints
and of showing its truth in all domains.  His
starting point was originally in the ethical and
social domain; here he first saw the law of the evolution brought about by
natural causes hold.  But a
firm form, a greater coherence came into his
views only when he came to know the truth ``which
was first obscurely hinted at by Harvey's embryological investigations,
which was afterwards grasped more clearly by Wolff, and
which was brought into a definite form by von Baer --- the
truth that all organic evolution is a transition from a
% --- p. 121 ---
\setcounter{footnote}{0}%
\opage{121}homogeneous to a heterogeneous condition.''  In one of his
first writings, ``Social Statics'' (1850), he presented not
only the evolution of man and of society, but also declared
that it was in principle the same law that held
for organic and for social evolution.  Afterwards he was
led to subject this law to a closer investigation,
which led both to a rational grounding of it and to
setting up more forms of the general law of evolution
than he had done at the outset.\footnote[1]{\emph{Spencer}: The classification of the sciences. 3d. ed. London 1873. p. 46.} This investigation
is contained in his ``First Principles'' (1862).  But before
we pass over to it, we shall make some general remarks
on the extent and the light in which
Spencer conceives the law of evolution.

Astronomy shows us, according to Laplace's hypothesis, how
the original mass of nebula gradually articulated itself
and was transformed, so that the present solar system proceeded
from it.  Geology shows how the earth from its original
glowing condition, which was a consequence of the heat developed by the concentration
of the nebular mass, has attained its
present condition through a series of changes:
the cooling and solidification of the earth's crust, the distribution of land
of water, the bedding of the different strata, etc.  Organic
life shows us a similar process both within the history of the individual
and within the history of the species.  It is always the
higher, richer and manifoldly developed forms that arise
latest.  As regards man, we can point out a whole
series of stages from the half-animal savage to the
cultivated European.  With regard to every single
human endeavor, too, the same holds, whether we
follow the history of language, of art, of science, of commerce or of
industry.  All phenomena point to the same
law.  But from this it follows on the other hand also that this
law is only a hypothesis to which we are of necessity
led; we cannot ascribe to it absolute validity, since
% --- p. 122 ---
\setcounter{footnote}{0}%
\opage{122}we do not know the Absolute and thus cannot know whether
this law has application to it.  The great problems
concerning the essence and ground of things are not solved
hereby; the last mystery is only moved one step back,
but remains in itself just as obscure as it has always
been.

In general one conceives the idea of evolution and of progress
in a one-sided teleological manner, as if all
stages in evolution were only to have as their end to further
human felicity.  But this is a finite
and purely subjective standpoint.  Philosophical inquiry
asks, on the contrary, which laws hold for the phenomena,
for the constant changes and progress,
quite apart from our interests, --- asks what evolution
is in itself, and what its causes are.  And
as evolution itself is only the expression of a phenomenal
law, so also its causes can only be phenomenal
and relative; the absolute cause is inaccessible to us.

As the law of evolution concerns only the relative and
the phenomenal, so it is in itself an expression of
the principle of relativity.  When all is in evolution, then
all is also only relative.  Therefore Spencer often emphasizes
how impossible it is to give sharply delimited definitions.  The
time of absolute boundaries is past.  One cannot point to
any point at which mental life begins; there takes place
a successive transition from the purely bodily functions
to the mental ones.  Just as little can one fix any
boundary between the organic and the inorganic phenomena,
between the living and the lifeless.  The newer physics
shows us what inner coherence there is between the
different forces, electricity, magnetism, mechanical
force, heat, etc.; the one of these can be converted into
the other.  Science itself stands in no absolute
opposition to popular consciousness; it commands no
means and forces that are not already present in
this consciousness, but marks a higher stage of development
of these same forces, an extension of the field of vision, a
% --- p. 123 ---
\setcounter{footnote}{0}%
\opage{123}more consistent and more exact carrying out of the points of view.


The law of evolution is the expression of the fact that the
view of the world based on relativity has become conscious of itself.  It
no longer falls out into the want of principle and incoherence,
but gathers itself in one principle.  Just as the critique of the
concept of matter in the most recent philosophy of nature has led to putting
materiality in the place of matter, so one may say
that Spencer has put relativity in the place of the relative.
But where there is a principle, there is to that extent also an
Absolute, an Unconditioned.  In setting up a principle and
developing on the basis of it a system, he therefore goes beyond
empiricism, and he is conscious of this.  By consistently
carrying through and thinking through the views of the empirical
realistic school he has got beyond
it and has grasped thoughts that underlie the speculative
(idealistic) tendency.  He also expressly
sets it as his task to reconcile Locke and Kant.

But we can already here indicate where the most difficult
point in Spencer's philosophy lies.  For the principle
to which the relative is traced back, and in which the law of the phenomena
has its ground, cannot be sought elsewhere than
in the Absolute and Unconditioned; --- otherwise it would not
be the true principle, but another would have to be sought.
The relative can be gathered under one principle
only if one presupposes that the relative and phenomenal
stands in an inner relation to the Absolute and Unconditioned;
only from this can all unity and coherence be derived.
Every real total view must mark a progress
in the knowledge of reality itself; the nominalistic
thought that only the isolated fact has reality dissolves
itself in the recognition that precisely unity and totality
is the true reality.  And toward such a recognition
Spencer's own philosophical view also
points.

The principle of relativity abolishes, as we have seen, the
fixed boundaries.  But when all boundaries fall, how
% --- p. 124 ---
\setcounter{footnote}{0}%
\opage{124}should the boundary between the relative and the Absolute
be able to remain standing?  Must not this boundary, which for
Spencer is one with the boundary between the knowable and
the Unknowable, itself become relative and fluid, i.e., be subjected to
the law of evolution?  --- There seems to be no other way out,
if Spencer will not fall back into the ``pure
empiricism,'' which he combats.  The dualism between the Absolute
and the relative would make Spencer's own endeavors toward
a ``synthetic philosophy'' fruitless.

Spencer maintains, as we shall see, expressly
the necessity of the idea of the Absolute as the last principle,
the last presupposition of all thought and experience.  Indeed, he
even intimates that the essential forms and laws of existence
that we can know must also be an expression of
the nature of the Unconditioned; he shows how the idea
of this in the course of the ages, under the development of religion and science,
has assumed ever higher and more perfect forms;
--- but herein it is implied that it cannot be absolutely unknowable.
If one admitted an absolute principle, but
held the relative outside all contact with it, then
this principle would be not only an unknowable, but a
meaningless concept.

The uncertainty that we find in Spencer on this decisive
point is naturally brought about by the twofold
influences he has received, --- from the empirical and
from the critical school, and it could hardly be avoided in
a thinker who has set himself so grand an aim as that
of uniting the two opposed currents in the newer philosophy.
The question might now be raised which of the
two currents gains the upper hand in Spencer's view.
If I conceive this rightly, the result in the last
instance seems to be that Spencer, in a still higher degree than he
himself is aware, has overstepped the boundaries of empiricism
and approached the ideal conceptions that have worked their way
forward through the philosophical systems of the Continent.  The grounding
of this the following presentation shall give. ---

Spencer has sketched and in part already carried out a
% --- p. 125 ---
\setcounter{footnote}{0}%
\opage{125}grand plan.  In the year 1860 he issued a subscription prospectus
for a philosophical system, a ``synthetic
philosophy,'' which was estimated at ten volumes.  Of these the
first five are available.  The first volume contains, under the name of
``First principles,'' the foundation of the whole view, in that it
investigates the relation between religion and science, sets up
the presuppositions on which all thought and experience
rest, and the general laws that emerge as
the last results of the investigations of the special sciences,
and also shows how these laws can be derived from those
presuppositions.  It is thus a metaphysics.  The two
next volumes give, under the name of ``Principles of Biology,'' a
presentation of the conditions and development of organic life, and
show how the forms and functions of life are formed
under interaction with the surrounding world.  ``Principles of
Psychology,'' likewise in two volumes, presents in a corresponding
manner the development of mental life in the struggle with its surroundings.
Of the five volumes that remain, the three
(``Principles of sociology'') are to contain the science of society,
the last two (``Principles of morality'') ethics. --- Besides
the writings mentioned here Spencer has published several smaller
works.  The ``Social Statics'' already mentioned develops
some of his ethical ideas; ``The classification of sciences''
gives, besides the subject stated in its title, a determination
of his relation to positivism; besides this he has collected
a number of scattered, very interesting treatises in two volumes
``Essays, scientific, politic and speculative.'' --- In our presentation
we follow his own systematic order.  Only thus
can we get an idea of the unity that runs
through his view, and of the energy with which
he seeks to make the different streams of knowledge
unite into one onward-hastening river. ---

The sharpest opposition in which different views
can stand to one another is that which we find between religion
and science.  It dates from the very first observation
and discovery of a law, of a regular coherence in
nature, and it has repeated itself at every later stage
% --- p. 126 ---
\setcounter{footnote}{0}%
\opage{126}of the development of religion and science.  Yet it is
not an irreconcilable strife.  Here as everywhere there must be
something on which both the contending parties are or can become
agreed.  Even in the false there is ``a soul of truth;''
even the wildest superstition must enclose a justified element.
Otherwise one could not understand how
religion has been able to hold its own in spite of the attacks of science,
repeated again and again, on its dogmas.  Even if one
will trace the dogmas back to the religious feeling as their
source, this feeling itself still remains as
so essential a fact that every philosophical theory that
does not explain it is incomplete.  If one assumes that
this feeling is immediately created, then one must also
assume that it corresponds to its author; if one assumes that it,
like all mental faculties, has developed under the interaction of
individuals with their surroundings, then one must
also assume that there are definite real conditions to which
it corresponds.  In both cases we are thus led to regard
the religious feeling with respect. --- What all religions
have in common can only be the truth that lies at the
ground of them all; and likewise what religion and science
have in common can only be an abstract truth in
which they both meet.  If it succeeds in demonstrating this
truth, in which the two poles of human consciousness
meet, then it will appear that the strife between religion and
science has been of the same nature as the fight
between the two knights over the color of a shield, of
which each had seen only his own side.

The highest religious ideas concern the nature and
origin of the world.  This question has been answered in three
different ways.  According to atheism the world exists
by itself and is not produced; according to pantheism it has
produced itself; according to theism it is produced by an
external cause.  The atheistic conception is untenable, since we
cannot form an idea of anything that has no
origin; existence in and by itself is for our
thought something inconceivable.  In any case this
% --- p. 127 ---
\setcounter{footnote}{0}%
\opage{127}conception gives us no answer to the question, no explanation
of the existence of the world.  Pantheism too is untenable:
what produces itself must beforehand have been
in existence under another form, --- and the question then returns
into the infinite.  In the same way it fares with
theism: it leads back from the world to the Creator and
assumes the Creator to exist by himself; --- but then we end
once more here at the concept, inaccessible to us, of
something subsisting in itself.

The critical consideration thus demonstrates the untenability
of all the main forms that the religious idea has assumed.
We cannot, as Hamilton and after him Mansel have
shown, form any concept of the Absolute and
Unconditioned, therefore every concept of God is full of contradictions.


From this we learn that the truth in the religions cannot
consist in the way in which they solve the problem of the nature
and origin of the world.  We must rather seek their truth
in the very acknowledgment that here is a problem, that there is given
a mystery that forms the background of existence and lies
beyond the reach of our knowledge.  All religions acknowledge
a mystery, but seek nevertheless, each in its own way,
to give a solution *of the riddle, a revelation, whereby
the mystery is then again in fact abolished.  Yet it appears that
the higher a religion stands, the higher the Unconditioned is raised
above the domain of knowledge.  Lowest stands the fetish-worshiper,
who imagines his god under the most everyday
and finite forms; polytheism shows us already
the Godhead under more ideal shapes; monotheism
erects altars to the unknown and unknowable God, and recognizes
at last that a God whom we could comprehend would not
be any God, and that it would be blasphemous to think
that God is as we can think him.  The mystery
is not, as the religions suppose, a relative, but an absolute
mystery.  When the religious consciousness has raised
itself to this point, it will be reconciled with science.

For it goes with the highest scientific ideas
% --- p. 128 ---
\setcounter{footnote}{0}%
\opage{128}as with the highest religious ideas; they are self-contradictory
when they are regarded as absolute.  What force and matter,
space and time are in themselves, we can know nothing about.
At no point do we reach a conclusion.  We can find
general laws, but these point ever back to still
more general laws, and so on.  Our thought
is, as Hamilton has shown, conditioned by opposition and
difference; what has no opposite and cannot be designated
by its difference from another, we cannot think.
On the other hand all knowledge also presupposes a
relation of likeness; we know a thing by referring it to
another like it; but the Absolute, that which subsists in itself,
we cannot compare with anything other than
itself.  Finally, all our knowledge is subjective, consists
in the agreement of our states of consciousness with the outer
phenomena; therefore it is necessarily also relative, does not reach
the nature of things in themselves.

The question then presents itself: is our consciousness
limited to the mere phenomena, to the finite and
relative?  Seen purely logically the matter seems to stand thus,
that the Absolute, as Hamilton says, is a purely negative
concept.  But on closer consideration
the matter stands otherwise.  There must after all be something else in consciousness
than the very limits and conditions of consciousness; there
must be something that is limited and conditioned.  Every
definitely formed state of consciousness comprises only a relative;
but it presupposes an indefinite, unformed content,
an Unconditioned and Infinite, which is the necessary material
for the work of thought.  The very concept of the relative presupposes
as its necessary correlate the reality of the Absolute;
otherwise the relative would indeed be absolute.  Every real
thought is finite, but points beyond itself; every
existent that we think points beyond itself.  Every new
form that thought wants to give the Absolute falls before criticism,
but the Absolute then creates other forms for itself, and
so on to infinity.  If the Absolute were not
in consciousness, thought would lack content, just as it
% --- p. 129 ---
\setcounter{footnote}{0}%
\opage{129}would lack form if only the Absolute were.  There is
something that is conditioned, something that stands in the relative
relations, and which we can express only by successively
abolishing all conditions, all limits.  We then come to
the positive though indefinite consciousness of that which lies
beyond all definite consciousness, of existence in and for itself,
of a substance behind all phenomena, a being behind all
appearances.  This consciousness is the necessary counterpart
of our self-consciousness.

``Besides the \emph{definite} consciousness, whose laws logic
formulates, there is also an \emph{indefinite} consciousness,
which cannot be formulated.  Besides complete thoughts, and
besides the incomplete thoughts that can be completed,
there are thoughts that it is impossible to complete
(to carry through), and which are nevertheless also real in the sense
that they are normal states of reason'' (normal affections
of the intellect). --- ``Since our consciousness of the
Unconditioned is literally the same as the unconditioned
consciousness or the raw material of thought, to which we by thinking
give definite forms, it follows that an ever
present sensation (sense) of real existence is the very
foundation of our mind.''

Religion and science therefore belong necessarily together,
as two forms of consciousness supplementing one another.  Our consciousness
of nature is, seen from one side, religion, seen from
another side, science.  Science attacks every
single supernatural fact that religion lays down, but
at the same time it concedes that its own results are only
relative, do not reach the absolute mystery.

But are we then not compelled to represent the Absolute
to ourselves under one form or another? --- Auguste Comte
supposed that the theological concept of a providence and
the metaphysical concept of nature had played out their role,
and that the ideal that now remained was that of positive science:
to trace all single facts back to one single
universal fact.  Against this Spencer objects that the
highest religious, metaphysical and scientific idea is in itself
% --- p. 130 ---
\setcounter{footnote}{0}%
\opage{130}one and the same.  When religion has freed
itself from all anthropomorphism, it becomes the belief in an absolute
being that indwells all phenomena; and science
leads us at last to the recognition of a universal force, one
existence, which as universal and absolute is inconceivable to us.
And when Comte held that the religious feeling will hereafter
have humanity for its object, this conflicts
with the law of the development of religion, which shows a steady
advance from the finite toward the infinite; the consciousness
that has raised itself to the idea of the Unconditioned cannot
again find its highest object in a finite being.\footnote[1]{The classification of the sciences. p. 35-36. 41.}
--- According to Hamilton and Mansel it is a duty, which with
necessity springs from our nature, to think God
as a personal being.  But it can be a duty neither
to affirm nor to deny God's personality.  Our duty
is only to recognize our limitation and not go beyond
the barriers of our thought.  There can be no conflict
between our intellectual faculties and our moral obligations.
When we refuse to believe in God's personality,
it is not because it is too high a determination,
but precisely because it is too low, finite and human a
determination of the Unconditioned.  Might there not
be given a form of existence that stood as high above consciousness
and will as these stand above purely mechanical motion?
We can certainly form no concept of
such a form of existence; but that is not a reason
to doubt its existence; --- we have indeed seen that the conditioned
cannot grasp the Unconditioned.  And yet thought
will perhaps ever continue its attempts to find expression
and definite forms for the absolute being that is its
own background and basis.  Nor will there be anything
wrong in this, so long as such definite forms are expressly
recognized as mere symbols without inner correspondence
with what they signify.  Under the constant
% --- p. 131 ---
\setcounter{footnote}{0}%
\opage{131}formation of such symbols and the ensuing rejection
of them as inadequate, consciousness is kept
alive to the opposition between the Unconditioned and the
conditioned.  The adherents of the positive religions cling
fast to the finite forms of the Unconditioned, as if they
could be absolute, and are therefore less pious than
the thinker, who knows that all such forms are only inadequate
symbols.  And even if he stands alone with
this conviction, he yet knows that it develops with
necessity and will have the future on its side.  ``He must
remember that just as he (with all his faculties, endeavors
and convictions) is a child of the past, so he is
father to the future, and that his thoughts are like children born
to him, which he must not let die for want
of care.  He may, like every other man, with
right regard himself as one of the countless organs through
which the unknown cause works, and when the unknown
cause produces a certain conviction in him, then he is authorized
thereby to express it in speech and action.
For --- to apply the poet's words in their highest
sense --- `Nature cannot be improved by any
means unless Nature herself has produced this means;
above the art which you say adds something to Nature stands
an art that Nature herself produces.'  Therefore
the wise man will not regard the conviction that stirs in
him as accidental.  He will without fear utter the highest
truth he can see, because he knows that whatever the consequence
may be, he thereby plays his proper part in the world; --- because
he knows that \emph{if} he can bring about the change he strives
for, good! --- if he \emph{cannot,} also good, though not
\emph{so} good!''\footnote[1]{\emph{Spencer}: First Principles. 3d ed. p. 123.}

One sees that Spencer makes himself no illusions with
regard to the reconciliation between religion and science
that he believes he has found.  It is also clear that the
% --- p. 132 ---
\setcounter{footnote}{0}%
\opage{132}positive religions cannot enter upon this reconciliation
without giving up precisely that which makes them positive, namely
the belief in a supernatural power revealed under certain definite forms.  According to Spencer the Absolute is the presupposition
to which every inquiry into the nature of consciousness and
of existence of necessity leads back.  It
does not stand in dualistic (supernatural) severance from the phenomena,
but in immanent relation to them, is ``the unknown
cause,'' ``the universal force,'' which underlies
them.  Therefore it is not clear how Spencer
can go on designating the Absolute as the Unknowable;
it would be so only if the phenomena did not stand in
inner relation to it, so that they could not furnish even
an indirect contribution to the conception of it.  Consistently
Spencer would have to conceive the phenomena as relative forms
and expressions of the Absolute, as determinations of content, in
which it expresses its nature, so that through
them we could win a real, though conditioned and relative, knowledge
of the Absolute.  This appears clearly when we
bring forward Spencer's determination of the task of philosophy
and its first presuppositions.

By philosophy he understands knowledge developed to unity and
totality (completely-unified knowledge).  It
seeks the general laws that comprise the particular laws,
which are the objects of the single sciences.  In that philosophy
everywhere seeks unity and agreement, it presupposes
the validity of the function of consciousness by which
we conceive things as like or unlike one another.  Even
the most skeptical thought must acknowledge this validity;
for doubt itself can be maintained only by means of that function of consciousness
that underlies every inference,
both denial and affirmation.  But besides this
function itself a result is also presupposed, to which it
has led.  Our ideas divide into two great groups,
which cannot be made one: a series of lively ideas,
which we comprehend under the name of object and
refer to outer things, and a series of weaker
% --- p. 133 ---
\setcounter{footnote}{0}%
\opage{133}ideas, which we comprehend under the name of subject or
self.  This distinction is the result of a series of acts of consciousness,
which establish likenesses and unlikenesses between the ideas.
The subjective and the objective, self and not-self,
stand thus as fundamental oppositions within the
relative, as the two main forms of the manifestation of the
Unknowable.  The relative reality that our knowledge comprises
we must of necessity think in inner connection with
the absolute reality, inaccessible to us.  The fundamental conditions
under which reality appears to us, ---
force and matter, time and space, motion --- forms that
underlie all conception and experience, we must therefore
ascribe real validity.  True, we have no right to assume
that they have absolute reality, are forms of the Absolute
itself; but since these forms have for us an unconditioned
and unchangeable necessity, we must also assume
that the unknown, absolute force acts unchangeably on
us: without such an unchangeableness those forms would be
unthinkable.  In the last instance all those forms can be traced
back to force: matter and motion are only variously
conditioned manifestations of force, and space and time are the conditions
under which these manifestations appear to us.  Force
is also the proper and original object of all consciousness.
The fundamental character of consciousness consists indeed in the conception
of likeness and unlikeness, thus in a comparison,
but this presupposes difference between impressions; we conceive
thus that which brings about change, --- but we can only
conceive it by means of the sensation we have of our
own force to bring about changes.  This consciousness
of force becomes thus the last point to which analysis
leads us.  It is naturally not the absolute force that we
here conceive, but only a conditioned effect of the unconditioned
cause.

The existence of force is now the necessary presupposition
for all science and philosophy.  So long as there is thinking,
so long as consciousness functions, so long must there also
take place transitions and changes: but thereby it is
% --- p. 134 ---
\setcounter{footnote}{0}%
\opage{134}presupposed, then, that force exists, that matter is not annihilated,
that motion is steadily continued.  If force ceased
to act, if matter became nothing, if motion stopped absolutely,
then we should get an act of consciousness that consisted in a
relation between something and nothing, i.e., a relation in which the one
term was lacking, i.e., a relation that is no relation,
a self-contradiction.  An absolute end we can just as little
think as an absolute beginning, indeed, we cannot think
of any consciousness for which this could be otherwise.
Every single form of force we can think of
as converted into another, but force itself we cannot abolish.
It is not any single, no relative force that
exists, but force as a whole.  We win here thus a
new standpoint for the truth that was maintained above, that
we must necessarily presuppose an Unconditioned as background
and bearer of everything relative and phenomenal.  The force
whose existence we learn is the absolute force, of which
we have an indefinite consciousness as the necessary correlate
of the force we know.  To maintain the existence of force
is thus only another form of maintaining an unconditioned
reality without beginning or end.

What thus follows from the nature of our consciousness,
a priori, is confirmed a posteriori by the results of modern physics and
chemistry, which establish the persistence of force and
the indestructibility of matter as facts of experience.
Indeed, every single experiment properly presupposes this
fact; for as soon as there is to be weighing or measuring,
a unit of force must be laid at the foundation, which must be assumed
constant; but if force did not exist, one would not
have any certainty of the constancy of this unit, and
all measurement, all exact experimentation, would thus be
impossible.

The only truth that goes beyond experience,
in that it lies at the foundation of experience, is the existence of force.
As the foundation of experience it must also be the
foundation of every scientific organization of experience.
% --- p. 135 ---
\setcounter{footnote}{0}%
\opage{135}To it a carried-through analysis leads us back, and upon it
a rational synthesis must be built.

From this fundamental presupposition of the existence of force it is
then again a necessary consequence that there must be a fixed
lawfulness in nature.  When in two given cases
complete likeness obtains between the antecedent
conditions, then the consequences must also be the same; ---
otherwise we should come to the consequence that either an entirely new force
must have arisen, or a force had ceased to exist; but in both cases we should
come into conflict with the proposition of the existence of force, which
cuts off an absolute beginning and an absolute end.\footnote[1]{One easily sees that this presupposition underlies Mill's ``four methods for experimental inquiry'' (above p. 28 f.).}

Likewise it follows from the existence of force that the different
single forces can pass over into one another,
so that a certain degree of the one force is equivalent
to a certain degree of the other force.  The apparent cessation
of a single force means only that it is converted
into another form.  Thus mechanical motion,
heat, light, chemical attraction, electricity, etc. can be converted
into one another, indeed they can be converted into sensation, feeling
and thought, as these can be converted into those.  \emph{How}
this conversion takes place, we do not know.  How
vibrations in the air can call forth sensations of sound,
chemical changes in the brain mental states, these
are mysteries that we cannot grasp, but not deeper mysteries
than the conversion of the physical forces into one another.
These conversions are not merely empirical facts; it follows
from the fundamental law of consciousness itself that every manifestation of force must
be a consequence of a preceding manifestation of force, and that at
its cessation it must pass over into another kind of manifestation of force;
if not, the proposition of the existence of force is abolished and
with it consciousness itself. --- In a similar way it can
be proved that motion must always go in the direction
where the least resistance is offered, and that motion, on
% --- p. 136 ---
\setcounter{footnote}{0}%
\opage{136}account of different forces acting simultaneously, must become
rhythmical ---

What we have found so far are only the abstract laws
that hold for the single factors.  But the knowledge of unity
that philosophy seeks is completed only when we turn
to the concrete phenomena that correspond to and can be derived
from those abstract laws.  With the concrete phenomena in themselves
philosophy has nothing to do; they are investigated in
the special sciences.  But when the geologist has explained
to us the history of the earth, when the chemist has taught us to know the
combinations and changes of substances, when the physiologist
has found the law of the growth of the organism, then there arises the
question whether there is not something common to all
these different processes, whether elements cannot be pointed out
that recur in the history of all these concrete phenomena.

The concrete phenomena have matter and motion
as their elements.  The law of their history must therefore
be one that states more precisely the form of the constantly renewed
distribution of matter and motion.  Now every existence
that we know undergoes two different processes.
It comes to be for us in that elements that formerly
were in a scattered and confused condition gather into a
whole, in that matter is concentrated and the single parts lose
their motion.  This is what we call evolution, which
thus consists in the gathering of matter and dispersion of
motion.  The other process that the phenomenon undergoes is
a dissolution, which conversely consists in the absorption of motion
and dispersion of matter.  Evolution and dissolution
designate the two main periods in the history
of every phenomenon.

Evolution, whose general character has been stated as
a gathering of matter and dispersion of motion, must
now be determined more precisely by way of induction and experience.
We must in astronomy and geology, in biology, psychology
and sociology seek the more special peculiarities
of evolution, as it reigns throughout
% --- p. 137 ---
\setcounter{footnote}{0}%
\opage{137}the whole universe, through all manifestations of the unknowable
fundamental force.

The first of these more precise determinations consists in this,
that evolution goes from a less coherent to a
more coherent condition.  According to Laplace's hypothesis
the solar system has condensed from an immense mass of nebula into
the single globes of which it now consists; and each of
these globes has itself undergone a process of condensation.
Every plant grows by concentrating in itself elements that
before existed in a scattered condition; every animal grows
by again concentrating in itself elements that before existed
scattered in plants or other animals.  Human
history shows us an ever closer union between
families and tribes, until there now seems to be forming a
civilized society that embraces the whole earth; and if we compare
a civilized state with a band of savages, we shall at once
see what a difference there is in concentration.

But the transition from a scattered to a concentrated condition
is at the same time a transition from a homogeneous to a heterogeneous condition.
The original mass of nebula became, during its concentration,
gradually articulated: rings formed,
which were detached from the main mass, and which again divided into globes,
until our present solar system with all its inner differences
developed.  The earth was at first a glowing
mass; then the crust solidified, and there came to be an opposition
between it and the inner fluid and glowing core; on
the surface of the earth there developed differences between land and
sea, mountains and plains, climatic differences, etc.  From
its first homogeneous germ the organism develops into a
rich manifoldness of parts and organs.  Human
society shows us the same phenomenon in what is
called the division of labor; originally everyone had to produce
for himself all that he needed; gradually the
different labors have become divided, so that everyone gets his
occupation.

The transition from the homogeneous to the heterogeneous brings with it
a transition from the indefinite to the definite.
% --- p. 138 ---
\setcounter{footnote}{0}%
\opage{138}Everywhere order takes the place of confusion.  This is seen most clearly
in the evolution of the organism.  Its structure and its organs
become constantly more formed and articulated.  From a
gelatinous mass it develops into a definite, bounded
whole with sharply marked parts.

If we now add to this that it is not matter alone, but
motion also, that undergoes these changes, ---
the functions of the organism developing thus
parallel to its structure, --- we get the following complete definition:
Evolution is a gathering of matter and a simultaneous
dispersion of motion, in the course of which matter passes
over from an indefinite and incoherent homogeneity
to a definite and coherent heterogeneity, and in
the course of which motion undergoes a parallel transformation.


Experience demonstrates that such an evolution takes place
everywhere.  Celestial bodies, organisms and human
societies have in essentials the same history.
But what experience thus teaches us, we must give a
deeper grounding; we must point out the causes of evolution
(naturally relative and phenomenal causes; the absolute,
noumenal cause is inaccessible to us), and these causes
we must derive from the highest truth, which lies at the foundation
of all experience and therefore goes beyond all experience, the proposition
of the existence of force. --- It is clear that in the
homogeneous mass, which always forms the starting point of
evolution, the single parts must receive different influences
from the surroundings.  There must first form an
opposition between the inner and the outer, and next
between the different parts of the outer and of the inner.  The
homogeneous is always unstable and therefore passes of necessity
over into difference.  Therefore, e.g., a species of plant
or animal must vary on account of the different conditions
under which the different individuals live.  Otherwise
different forces would have the same effects, which
conflicts with the proposition of the existence of force. --- But when
different parts are thus formed, then also one
% --- p. 139 ---
\setcounter{footnote}{0}%
\opage{139}and the same force must have different effects on the different
parts.  The force becomes heterogeneous, like the aggregate
it acts on.  And this means again that one and the same
cause can have manifold effects; the effect is
always more composite than the cause.  Light is modified by
the effects it calls forth by falling on different
objects, and thus the different colors arise.
This law too follows from the existence of force: different
things mean indeed nothing other than things that react in different
ways; if therefore one and the same force were not modified
in different ways by the different things, then
a (reacting) force would have been annihilated, which is impossible.
--- A third cause of evolution is to be sought in the segregation
of the like.  As the wind gathers the withered
leaves, so in general one and the same force acts
uniformly on that which is alike in all respects.
Things and persons who are subjected to the same conditions
will also come to resemble one another.  Hereby is brought about
what Darwin calls natural selection, which comprises
those beings that are able to fight the struggle of life
without succumbing.  Kindred beings are thereby brought
together, and greater definiteness of form and expression becomes
possible. ---

Now that the nature and causes of evolution have been demonstrated,
the question arises: does evolution continue
without interruption to infinity, or does it find its conclusion in
a definite lasting condition?  Every motion continues
until equilibrium is attained, and the greater the resistance is
from without, the more quickly is equilibrium attained.  When the external
resistance is not considerable, the inner motion
between the parts can continue, so that the whole finds
itself in a moving equilibrium (eqvilibrium mobile).  But
every such moving equilibrium has yet a tendency
to go over into absolute equilibrium.  The external forces never
cease entirely to act, even if their action is comparatively
imperceptible. --- The solar system may be said to be in moving
equilibrium, in that every motion or perturbation is counterbalanced
% --- p. 140 ---
\setcounter{footnote}{0}%
\opage{140}by a motion in the opposite direction.  The motion of the planets can be
regarded as an oscillation between aphelion and perihelion.
But from the diminution in the period of revolution of Encke's comet
it has, as is well known, been concluded that the celestial space is filled
with an ether that hinders the motions of the globes and
in the end must be able to stop them; then all the smaller
globes will again be gathered into one immense mass.  Still more important
is it, however, that the sun will one day have radiated all its heat
into cosmic space; and since the sun's heat is the source of all motion
and all life, this will bring with it a complete rest,
an absolute equilibrium. --- As regards human life, there
will one day come the time when the human
mind has worked its way forward to complete harmony with nature.
Be this goal ever so infinitely far away, it follows
yet from the nature of motion and of life that it is that toward
which motion goes; in other words, evolution
must end with the bringing about of the highest perfection
and felicity. ---

This too follows from the fundamental law of the existence of force.
Since we cannot think of force in its absolute
unity, but only under the form of interaction between different
forces, all motion takes place under resistance.
That which is moved therefore constantly gives off a part of
its motion to that which offers resistance, the medium
through which it moves; in the end it must then have
lost all its motion.

When equilibrium is attained, evolution passes over
into dissolution.  Whereas evolution consists in the gathering
of matter and dispersion of motion, dissolution consists
conversely in the gathering of motion and dispersion of matter.
That dissolution sets in follows from the continued
action of the external forces.  When equilibrium is attained, there is indeed no
force left over to meet these external forces with.  Therefore
we see everywhere, both in the history of human societies and in that of the
single organism, both in chemistry, geology and
astronomy, that an evolution at its highest point passes
over into dissolution.

% --- p. 141 ---
\setcounter{footnote}{0}%
\opage{141}Have we, then, reached the absolute end here, and will the universe at last congeal into an eternal death, become an immense heap
of extinguished solar systems? It follows of itself that we cannot find any absolute answer to such a question.
We know only the evolution within our part of the universe
and do not know whether there are systems of stars everywhere
in infinite space, so that new motions might continually
arise through their influence. The Infinite we cannot
construct. If, on the other hand, we hold to the fundamental law of
the persistence of force, then it is a necessary assumption that
motion can never cease. Equilibrium and dissolution
can form only the introduction to a new period of evolution.
That motion is rhythmical because of the multiplicity of forces
is a law that holds also for the universe in its
totality. In one period the attractive force has the upper hand,
the masses concentrate, and motion is dispersed; in
another period the repulsive forces have the upper hand,
motion is concentrated, and the masses are dispersed. It means,
then, that evolution and dissolution alternate in the infinite.
The result of evolution is the conditions that bring about dissolution,
and the result of dissolution is the conditions that
bring about evolution. Before our world there lies,
then, an infinite series of periods of evolution and dissolution,
and after the end of our world there will follow a new
infinity of such alternating periods. The principle in
these is the same, although the particular results will never
be the same.

This is the conclusion to which our limited knowledge
leads us. Just as little as we can conceive
any limit to time and space, just as little can we
conceive any limit to the unknowable force that
manifests itself in the universe. Itself infinite it must work in
the infinite. --- Since we do not know the Absolute, we know
also no absolute necessity, but only relative necessity,
necessity for our thought. Just as little
do we know the inner relation between the phenomena and the
absolute force that underlies them (the connection
% --- p. 142 ---
\setcounter{footnote}{0}%
\opage{142}between the phenomenal order and the ontological order).
All our fundamental concepts (matter, motion, force) are only
symbols of an unknown reality. When, therefore, the charge
has been raised that the philosophy of evolution is materialistic,
because it traces everything, even the phenomena of life, of the mind
and of human society, back to the general laws of matter,
motion and force, one overlooks
precisely that these fundamental concepts themselves are symbols of the Absolute,
which manifests itself in all phenomena, in the so-called
material as well as in the mental, and cannot find its
complete expression in any single class of phenomena.
If one holds this fast, then the strife between spiritualism
and materialism loses its significance, and one will see
that the philosophy of evolution does not drag the higher down by
its explanations, but precisely the reverse, raises the so-called
lower. ---

Spencer's ``First Principles'' is one of the most epoch-making
works of our century. Nowhere does one find, as here, all ideas, all scientific results worked together
into a firm, organic unity. The richness of the
inductive foundation, the sharpness and certainty of the deductive
grounding and of the psychological analysis, which forms
the presupposition of induction as well as of deduction,
are equally admirable. Spencer stands, so far, on the
standpoint of empiricism, in that he maintains that the highest
point we can reach is the most general phenomenal
law that experience teaches us to know. But he does not stop,
like Mill and Comte, at empirical universal propositions; he
traces the empirical laws back to a single law, the law
of the persistence of force, and then seeks to derive this law
from the very nature of thought and consciousness. Thereby there comes
an inner coherence between the different members of
his view. Philosophy as a unitary conception of existence
must proceed from an original presupposition, which
cannot be sought elsewhere than in the very nature of thought, of knowledge.
This presupposition, and the necessity of tracing
all other knowledge back to it, no thinker in
% --- p. 143 ---
\setcounter{footnote}{0}%
\opage{143}our time has maintained and carried through with such energy
as Spencer.

In particular matters there were many preparatory works. In the philosophical
respect Spencer is especially influenced by Hamilton.
He takes over the latter's doctrine of the relativity of knowledge,
but shows that relativity necessarily presupposes an
absolute principle, whose manifestation and symbol it is. Likewise
he has taken over and carried through Hamilton's doctrine of
the principle of causality as grounded in the very nature of our consciousness.
The proposition that everything has its cause was, according to
empiricism, derived from experience; but Spencer shows that
it is only another form of the law of the persistence of force
and, like this, has its ground in the innermost nature of our consciousness.
He differs from Hamilton in the close connection
that he thus establishes from the first between the results of physics
and of psychology. Hamilton was a strict spiritualist
and had no sense for natural science; for
him such a synthesis, whereby a
common fundamental law is set up for the world of the mind and of nature, would stand as
materialism and atheism, and Mansel has also accused
Spencer of having pressed Hamilton's doctrine in this direction
by accepting only the negative (critical) part of his
philosophy and omitting the positive (moral-religious, or
rather: theological). According to Mansel Spencer's doctrine is a
combination of pantheism and positivism (Philos. of the
cond. p. 39---40). But the truth is that Spencer, by
a consistent carrying through of the Hamiltonian principles, has
arrived at his peculiar doctrine, which, in spite of the vacillation
we have already pointed out in it as regards a decisive point, is yet not an eclectic combination, but a
well-cast development of a deeply grounded fundamental thought. One
might perhaps compare his relation to Hamilton with
Spinoza's to Cartisius.

Besides philosophical preparatory works, the results of natural science also form a necessary foundation for his
views. As moments of special importance we emphasize:
Laplace's theory of the evolution of the solar system; the doctrine of
% --- p. 144 ---
\setcounter{footnote}{0}%
\opage{144}the mutual transition of the forces of nature into one another, as
it is most clearly presented by Joule and Helmholtz; v. Baer's
definition of life as a development from the homogeneous to
the heterogeneous; Darwin's hypothesis of the origin of species.
Yet Spencer had in part himself anticipated several of these
results, and in any case his originality consists
in the working out and the joining together of all these particulars
into a self-contained whole.

In certain respects one could compare
Spencer's philosophy of evolution with Hegel's system, whose
fundamental thought is the dialectical movement from the immediate
through the opposition (reflection) to the Idea. Just as
in Spencer evolution begins with a homogeneous state,
which on account of its homogeneity is in an unstable equilibrium,
whose upsetting leads to articulation and individualization,
so according to Hegel immediacy lies open to
contradiction, which of necessity dissolves it into reflection,
until the rift and the distinction are removed in the Idea (Spencer's ``moving
equilibrium''). But what in Hegel is an abstract
logical process, in which the content of reality is pressed
into the forms of the movement of thought, is in Spencer a real
natural process, which embraces the real universe in all its
particular parts. True, Spencer too traces the idea of evolution
back to the nature of consciousness, but this metaphysical
grounding forms only the presupposition and does not furnish a
compelling form for the real content. Whereas Hegel's method
is absolutely deductive, in Spencer induction and
deduction support one another, in that the results of induction are deduced
from the ultimate presuppositions of all knowledge and consciousness.
--- Hegel had found the motive for his doctrine in an
analysis of the movement of thought, which showed three different stages,
of which the one led with logical necessity to the
other. In Spencer, on the other hand, it is, as we have seen, an
ethical-historical idea that has furnished the motive; this idea
was then brought into connection with the idea of life in general,
and thereby the way was paved to a wider conception and deeper
% --- p. 145 ---
\setcounter{footnote}{0}%
\opage{145}grounding, whereby the ethical-historical standpoint became only
a member within a larger coherence.

The main opposition between Spencer and Hegel rests
nevertheless on this, that the latter meant to give a theory, a development in concepts
of the nature of the Absolute, whereas Spencer only
gives his view as a hypothesis, which indeed in his
conviction is the only tenable one, because it is the
only one that agrees at once with the fundamental law of consciousness
and with the results of experience, --- but which yet cannot
teach us anything about the Absolute in itself, since all our
knowledge concerns only the relative and the phenomenal.
According to Hegel the Absolute develops in everything and comes to
itself through the evolution of the world; according to Spencer we dare
not apply the idea of evolution to the Absolute, since every
determination in concepts of what lies beyond the phenomena
involves us in self-contradiction.

Here we stand once more at the main difficulty in Spencer.
When the relative and the phenomenal are recognized as forms
of the Unconditioned,\footnote[1]{First Principles. p. 278: ``The sphere of knowledge is co-extensive with the phenomenal --- co-extensive with all modes of the Unknowable that can affect consciousness.''} then the knowledge of these forms
must also lead to determinations of the essence of the Unconditioned itself.
It would be precisely self-contradictory if the determinations
that hold for the forms did not have at least a relative
validity for the essence itself, which finds its expression in these
forms. ``One sets the Infinite over against, above, beyond,
outside the finite, indeed even puts up an immense partition,
so that they should not come too near each other. But
they do not lie outside each other at all. The finite is the
content of the Infinite. Therefore the Infinite is also not
incomprehensible; it can be grasped in its content; only encompassed it cannot
be.'' (Fechner). Evolution cannot therefore be a
necessary fundamental idea within the finite without also
having validity for the Infinite. As we have already
said, that which is for us the highest necessity is perhaps
% --- p. 146 ---
\setcounter{footnote}{0}%
\opage{146}only a relative necessity, but yet a real necessity,
which we must let stand, and which only fancy
can set itself above. The Unconditioned would not deserve
its name if its essence were not life and evolution;
that is a truth to which the metaphysical investigations in
the idealistic systems have led. And the evolution
that must be ascribed to the Absolute cannot be an abstract,
ideal and timeless one, as in Hegel, the thinker who with greatest
energy has maintained the application of the concept of evolution to
the Absolute; if one recognizes, with the post-Hegelian philosophy,
that time, space and matter cannot be placed outside
the Absolute, but that ``the concept of nature links the consciousness of the world
to the consciousness of God'' (Weisse), then its
evolution must be seen as real, its perfection not as completedness
but as unlimited progress.\footnote[1]{Philosophien i Tydskland. p. 216 f. 262 f.} Admittedly
we lack means for a knowledge of the evolutionary process
of the Absolute, and Hegel could only by a speculative
illusion think he gave one; but that cannot exempt
us from going the way that logical consequence demands,
especially since consequence here is of such great importance for
the whole view of the world.

We have already seen earlier that Spencer himself does not
maintain this distinction, when he ascribes to the Absolute
the determinations force and cause, which can yet be derived only
from the relative. But he goes still further. Of the three
chief religious views, whose untenability he thinks
he has demonstrated, he can really be said to stand in decided
opposition only to atheism and vulgar theism:
he maintains, in opposition to the first, that the religious
consciousness is in the right as regards the central matter, and
against the second, that this central matter cannot find its
expression in the form of personality, which for him is a finite
form.\footnote[2]{We shall later come to mention that Spencer decidedly sets evolution in opposition to creation, so that the former makes the latter unthinkable.} To that extent he can be called a pantheist, although
% --- p. 147 ---
\setcounter{footnote}{0}%
\opage{147}he does not like such dogmatic expressions. And it must
also be maintained that the Absolute is for him not an impersonal,
but a superpersonal; he has seen that the debate
between pantheism and theism has demonstrated the falsity
of the opposition between the personal and the impersonal,
when the subject is the Absolute.\footnote[1]{Op. cit. p. 161---165; 213---214; 279; 313; 324---325.}

In his early work Spencer conceived evolution
teleologically and saw harmony and happiness as the content of
the divine idea that was to be realized through
it. Later also he maintains the beneficent results
to which evolution must lead; but since his standpoint is now
objective and universal, so that human evolution
becomes only an episode within the great world-drama (or
dramas), there can no longer be any question of such
a teleology. Hence people have complained of the dreariness of
this view, that Spencer --- after having shown us
how every process of evolution ends in equilibrium and
dissolution, whereupon new processes follow in the infinite, ---
``leaves us anxious, bent over the brink of the eternal
abyss, as with our thought we question the infinite
darkness.''\footnote[2]{\emph{Caro} in Revue de deux mondes. 15 octbr. 1873 p. 774.} If any philosophy, any science
can answer the questions that here arise, the answer can
only proceed from the insight on which we have laid so much weight in
dealing with Spencer, that the relative stands
in an inner relation to the Absolute; the value that evolution,
its struggles and its results have from the relative,
human standpoint cannot be lost in the
Absolute. ---

From Spencer's ``general philosophy'' we now go on to
his ``special philosophy,'' of which only biology and
psychology are available fully worked out. Yet we shall find in his
``Essays,'' in ``Social Statics'' and in the book ``The study of sociology,'' which appeared while this work
was being printed, contributions to his ethical and sociological view. ---


% --- p. 148 ---
\setcounter{footnote}{0}%
\opage{148}Spencer's whole system has in a way its germ in
the doctrine of the more recent physiology concerning life as a development from
homogeneity to heterogeneity, an articulation and individualization.
This doctrine he applied first to the
life of human society and next to existence as a whole,
and thereby won a deeper and more many-sided determination of
the general laws of evolution. These laws must now in turn
find their clearest expression in the evolution of living beings.
Life consists in a definitely formed combination of heterogeneous
changes, both simultaneous and successive, whereby
the inner relations are brought into correspondence with the outer.
Life is in its generality an adjustment of the inner
to the outer. On lifeless objects the surroundings make
their impression without these impressions having any
further consequence; in every living body, on the other hand,
a secondary change takes place as a consequence of the impression received,
a change whereby the living being becomes
able to meet a following change in the external world.
Thereby the adjustment to the outer relations takes place; the organism
undertakes inner changes in order to be able to meet new relations.
Life presupposes, then, an interaction between the
inner and the outer. In order that the living being may be able to
subsist under the outer relations in their multiplicity and
diversity, a system of corresponding
inner diversities must be developed. Function, activity is thus
the principal moment in the idea of life. Organic function precedes
organic structure. Function develops structure;
structure, on the other hand, gives function
its more definite, distinct form. An example of life
without proper organization we have in the so-called rhizopods,
which show no separation of parts or organs, but
nevertheless eat, sleep and move. They consist of a microscopic
mass of slime which constantly changes its form, in that it
shoots itself out into fine, irregular fringes, which soon
draw back again; these movements serve both for
the nourishment of the animal and to bring it forward in the water. Here
% --- p. 149 ---
\setcounter{footnote}{0}%
\opage{149}we have in the most elementary form the adjusting of inner
relations to outer, in which life consists.

What the living being takes up from the external world
it transforms into a form peculiar to it. When
the organism has been damaged, it has to a certain degree
the capacity to make itself whole again. In lower animals a new limb of the same form as the lost one
can at once grow out. Reproduction and the inheritance of determinate qualities show
us a similar capacity to shape the parts after a certain
pattern. Even inorganic matter has such a
capacity. The material atoms have their determinate conditions,
under which they come into equilibrium. The more composite
the molecule is, the more unstable is its equilibrium.
And since now the organic molecules are the most composite
of all, the equilibrium here is the most unstable. A
slight influence is enough to bring about disturbance. But
just as we already find in the crystal a capacity to
shape the atoms that are taken up or brought out of equilibrium
in a manner agreeing with the structure of the crystallizing aggregate as a whole,
so we must assume a similar
capacity in the organism. When the leg is torn off a lizard,
the germ of a new one at once grows forth, which, after having gone through
the same process of development as the earlier one, assumes
a similar form. Here the structure of the organism determines
the development at the particular point. But
this can happen only when the form for every kind of organism
is determined by a peculiarity in its elements
(the physiological units), in that these, like the inorganic
atoms, have a special structure, according to which they strive
to arrange themselves. The force by virtue of which inorganic atoms
are heaped up into a crystalline form peculiar to them
we call their polarity; by analogy with this one might call
the corresponding capacity in the elements of living bodies
their ``organic polarity.''

These elements, the physiological units, are on the one
hand different from the chemical units that compose
the organic bodies; for these are about the same
% --- p. 150 ---
\setcounter{footnote}{0}%
\opage{150}in all kinds and thus do not explain the differences between
them; --- on the other hand from the morphological units
or the cells; for in part the cell is not a universal form
(the rhizopods, for example, are not cellular), in part the very
origin of the cell presupposes precisely that organic polarity. The
physiological units are much more composite than the
chemical, so that a slight change in their composition
will bring with it a change in the form that the whole aggregate
assumes. Thereby it becomes possible that the outer
influences can serve for the development of organic life.
The great plasticity of the physiological units gives the impressions
their power, whereas their polarity determines the manner
in which these impressions are worked up and shaped, the direction
in which they come to act. The lowest organic
phenomena are as yet wholly without organization; one can
speak in their case of neither form nor type. Their shape
changes from one moment to the next, according
as the outer conditions bring it with them. An absolute beginning
of organic life we cannot point to. The results of chemistry
on the one hand and of biology on the other
seem to entitle us to the assumption that the organic
compounds have gradually developed,
and that they have gradually won greater and greater
firmness.

If the compensation that the organic body is able
to procure for the losses its functions cause went only
just up to the amount lost, there would take place no development or
adjustment to the external world. But when the function and
the loss it causes do not exceed certain limits, the loss
is more than made good by what is taken up. Movement
and activity are, both for the whole organism and for its
individual parts, more favorable to the accumulation of cell-tissue than
rest is. It is a well-known fact that the muscles grow by
exercise. And since the different functions stand in connection
with one another, an increase of the one entails
also an increase of the other. Muscular movement thus presupposes
the circulation of the blood, and both must therefore be developed together.
% --- p. 151 ---
\setcounter{footnote}{0}%
\opage{151}But from this it follows also that the adjustment to the external world
has its limits, particularly in the full-grown individual, and
that much is required for it to be transmitted to descendants.
This takes place only when a new state of equilibrium has been formed on account of the changes. Then
the sum of all earlier influences and experiences can
gain the upper hand over the individual, isolated influences,
and the modification that was at first only a temporary and
individual adjustment becomes fixed and hereditary.

By this way the origin of new species must be explained.
It is true that it cannot be directly proved that the species have arisen
through continued adjustments and the effects of these adjustments on
structure (the organic polarity); but we know at any rate so
much, that in this manner profound modifications take place,
varieties of the same species are formed, which can differ
greatly from one another; indeed, one can often not decide whether a
certain difference constitutes a difference of species or only denotes
a variety. This vagueness in the concept of species shows
us that it is only a relative concept.

It follows, moreover, from the foregoing development that
the species must be assumed to have arisen in the interaction
of living beings with the external world. The assumption that a species
should not change by being exposed for millions of years
to changed influences would be a denial
of the proposition of the persistence of force. All organisms are
constantly exposed to astronomical, geological, meteorological
and organic influences. Every new capacity, every
advance in organization that is acquired, carries with it the possibility
of new influences, so that the changes must
increase in geometrical progression. But the presupposition is
always that the changes are not too great or
too sudden. Only such changes develop as
do not abolish the general harmony between the influencing
forces and the functions of the organism.

Spencer has thus by his own way come to the
result that Darwin maintained in his famous work on the ``Origin of
Species.'' True, Spencer's biology is later than this work,
% --- p. 152 ---
\setcounter{footnote}{0}%
\opage{152}but he had already earlier, in various treatises
(see his ``Essays'') and in the first edition of his Psychology,
maintained the hypothesis of evolution, for which reason Darwin also in the preface
to his work names him among his predecessors. Yet
there is constantly a difference in standpoint between the two
investigators. Darwin lays, as is well known, the main weight on
the so-called selection of quality, i.e. on the fact that
the organic beings which cannot take up the struggle
with the external world must perish, and only those beings are preserved
and propagate which are strong enough, but especially
those which in one way or another have acquired a
superiority over the others. This Darwin also calls the
natural selection: the strongest lives, the weak go
under. Thereby the harmony between life and its conditions
is maintained and developed. But according to Spencer's view
Darwin has laid too predominant a weight on this indirect
development. He himself ascribes the greatest significance to the direct manner of development: the working of the functions forms
the structure, and this continued adjustment gradually passes
over into the nature of the living being, in that there is formed a
new kind of equilibrium (organic polarity) between the physiological
elements. So long as the organism is still imperfect
and passive, the selection of quality, the indirect
production of equilibrium (indirect eqvilibration),
plays, to be sure, the greatest part. But as activity develops,
the indirect process passes over into the direct.
The selection of quality of course continues to play a great part,
but becomes more and more a subordinate factor. When
a nervous and muscular system has developed, and with it
a capacity to vary actions according to circumstances,
then the direct production of equilibrium (direct
eqvilibration) becomes the most important. In the civilized races of men
finally, the indirect process strikes only those whom the order of society
\emph{as yet} is not able to protect.

The relation between the direct and the indirect process
may be stated briefly thus: An external influence will
kill the organism if it is too weak, --- and in that case
% --- p. 153 ---
\setcounter{footnote}{0}%
\opage{153}the indirect manner of development operates; --- if the organism
on the other hand is strong enough, it will adjust itself to the impression,
and then the direct manner of development operates, in that the adjustment
leads to a new state of equilibrium within the organism.

This theory, the theory of evolution, Spencer sets in
sharp opposition to the usual view, the theory of creation,
according to which every species is created separately after a
fixed type. Spencer subjects this latter view to a
crushing criticism. In the first place the idea
of a creation contains a thought impossible in itself. It posits
an absolute beginning, sets up a relation between Nothing and
Something, a relation which is not a relation. From the very
nature and fundamental law of our thought there follows thus the
invalidity of the hypothesis of creation. Those who conceive every species as created
separately do so without really thinking, without translating
the words into thoughts; they do not in reality believe it, but believe
that they believe it; for belief presupposes that one ideally constructs
the matter for oneself, but that is here impossible. Next it would
follow from this hypothesis that the Creator had adhered to the
type laid down, even where it was unsuitable, and where
other forms would have fitted better. The members of one and the same
class of animals live, after all, under many different
conditions. Thus the seal is a mammal, but lives in
the water. If one now assumes that it was created for its present
abode, then it is meaningless that its limbs have
the same fundamental form as those of land mammals, since they are after all supposed to
be organs of swimming, whereas it becomes easily explicable
if it descends from a mammal that little by little
became an aquatic animal in its mode of life. Finally it would follow from
the hypothesis of creation that certain species were created to harm others,
and that pain and distress were a part of the divine
plan. According to the hypothesis of evolution, on the other hand, there is brought about
in the course of evolution slowly but surely a higher
degree of happiness, in that all evils are only transitory. A greater correspondence
must continually be brought about between
the living beings and the external world. To be sure, one
may ask why all evils, all distress and pain could not
% --- p. 154 ---
\setcounter{footnote}{0}%
\opage{154}have been avoided; but there can no longer be any question
of their being inflicted with premeditation.

The hypothesis of evolution teaches us just as little as the hypothesis of creation
to know the real cause; the strife
between them concerns only the manner in which the unknown
cause works. We find now in every domain that the
anthropomorphic conception of the unknown cause
disappears; one gradually ceases to derive the phenomena
from the immediate intervention of a divine will,
although many who in other domains have given up the belief
in such an intervention still maintain it in the question
of the origin of species; --- and yet it has here no
better warrant than elsewhere. --- The hypothesis of evolution
has emerged at once in many different domains:
within astronomy, geology, physiology, politics and
philology. Just as according to Laplace the solar system has through
a process of evolution conditioned by natural causes attained
its present shape, so according to Lyell the causes still
working daily have through a series of imperceptible
changes produced the present appearance of the surface of the earth.
Physiology no longer sees in the germ a ready-made model of the
mature organism; one recognizes in politics Coleridge's words
that constitutions are not made but grow; and one no
longer regards the origin of language as a supernatural
act. ---

Now there is, to be sure, a standpoint midway between
both hypotheses; it is that taken up by \emph{Auguste Comte}\footnote[1]{Cours de philosophie positive. 2e éd. III p. 390 f.}
and with him many investigators of natural science, who hold to the factual, that a transition from one
species to another has not yet been demonstrated, and who therefore
will set up no theory at all of the origin of species.
Recently \emph{Charles Robin} has with strong emphasis
repeated Comte's objection to such hypotheses.\footnote[2]{Anatomie et physiologie cellulaires. Paris 1873. p. XXIV: ``If one leaves aside an anti-scientific intrusion of \opage{155}theological presuppositions in the biological problems, then it is by no means a lack of philosophical culture in the investigators of our time that in France diminishes the number of those who adhere to the doctrines of Lamarck and his successors; ... The reason is on the contrary the developed philosophical insight that the validity of a speculative view must be judged by the possibility of a positive control, and the lack of a real proof on this point, in that up to now such a proof has not been given for a single class of beings, any more than for a single kind of cells ... As yet only possibilities have been adduced as proof, on which few naturalists are agreed, and not real facts.'' --- Another of Comte's most famous disciples, \emph{Littré}, takes a more accommodating, yet cautiously waiting, attitude toward the hypothesis of evolution (La science au point de vue philos. p. 539-559).} ---
% --- p. 155 ---
\setcounter{footnote}{0}%
\opage{155}Spencer regards this skeptical and waiting attitude
as a state of ``unstable equilibrium.''  Involuntarily thought must
go in the one direction or the other: the origin of species
must after all be \emph{either} natural \emph{or} supernatural,
a third is not given. And when one in general
holds a rational view of the world, it follows of itself
which member of this dilemma one must seize. When Comte
declares the problem of the origin of species to be the fruit of
useless speculation, Spencer in opposition to this declares
that ``the part of biology which concerns itself with
the origin of species is the most important, in relation to which
all others are subordinate. For on the decision of biology
in this matter rests our conception of human
nature in past, present and future, our theories of the human
mind and of human society.''\footnote[1]{The classification of the sciences. p. 40.} ---

All morphological and physiological phenomena must
now be able to find their explanation according to the hypothesis of evolution.
The organic forms can no longer be regarded as
types that have hovered before the Creator, or as means he
has wished to employ for the good of the individual; they must be regarded as
necessary results of the interaction between the individuals
and the external world. The mechanical conception thus takes the
place of the typical and the teleological. Spencer sets up,
in order to demonstrate the possibility of this conception, a series of
% --- p. 156 ---
\setcounter{footnote}{0}%
\opage{156}detailed, in part extraordinarily acute hypotheses, which
we cannot dwell on here. His main thought is, on the
basis of comparative anatomy and physiology, to
demonstrate a gradual development parallel to the development
of the conditions. Even the highest and most composite structures
and functions must be explicable by this way. Already
the structureless protoplasm possesses, e.g., in a lower degree,
the same property that the nerve in the higher animals possesses.
The gelatinous mass of a rhizopod and an unfertilized
egg show movements which presuppose that an impression has
been propagated from one part of the mass to the other;
and through the nerveless body of a polyp we see a contraction evoked by the touch
of a feeler slowly propagate
and spread. By continued function the
special organ for this function must then be able to develop.

The degree of life depends on the degree in which the correspondence
between the organism and the external world is comprehensive
and exact. As a system of inner differences is developed,
corresponding to the outer, the organism becomes ever more
composite; but the connection between its parts also becomes
gradually more intimate. In an organism as developed as
the human, a slight impression can produce a
multitude of different effects, can evoke muscular movements
and contractions and arouse ideas, feelings and
resolutions. ---

As there cannot be pointed out any definite point where
the inorganic passes over into the organic, so we can
likewise posit no fixed limit between the physiological
and the psychological phenomena. Both classes
belong under biology and follow the general laws of life.
As geology is a special part of astronomy, and
biology of geology, so psychology is in turn a
subdivision of general biology. It differs
from the sciences on which it rests in that it does not,
like physics, consider the external phenomena in and for themselves,
or, like physiology, the internal in and for themselves, but investigates
the connection and relation between the internal and the
% --- p. 157 ---
\setcounter{footnote}{0}%
\opage{157}external. Objective psychology studies the functions by
which the organism is able to adjust itself to the external
events and relations, and the manner in which these
functions appear in the actions of individuals. Subjective
psychology stands in opposition to all other sciences.
It investigates the sense-perceptions themselves, the ideas,
the feelings and the expressions of will, which accompany the
visible adjustment of the inner relations to the outer, and considers
these phenomena of consciousness in their origin and
mutual connections. Its object, consciousness itself,
and its method, analysis of the self itself, whereby alone one
can find the laws of the changes in consciousness, place
it in the same relation to all other sciences as the
subjective stands in to the objective. Psychology (both
as objective and as subjective --- for the latter is a presupposition
for the former) has therefore a still greater claim to be
an independent science than any other whatsoever.

Mind and consciousness are something that has no likeness at all
to other things, are something unique in its kind. The substance
of the mind is that which persists through all changes in
consciousness, and whose modifications all these changes
are. Therefore Hume is wrong when he teaches that nothing
exists but impressions and ideas. Every single
sensation or idea could be absent,
but that which connects the sensations and the ideas
cannot be thought away; it is that which maintains the unity of the circle of ideas
in spite of all attempts to divide it,
--- the substance of the mind as opposed to all the changing
forms it can assume. But we cannot know
this substance. Every state of consciousness is a modification
of it, but in no state of consciousness is it wholly
present. Moreover, all knowledge rests on the opposition
between subject and object; but here the same would
have to be the subject and object of knowledge, which is impossible.
All knowledge rests further on a comparison,
but since the mind is sui generis, we have nothing to compare
it with.

% --- p. 158 ---
\setcounter{footnote}{0}%
\opage{158}Psychology can thus no more grasp the nature of the mind
in itself than chemistry can grasp the nature of matter
in itself. Yet it can give an explanation of the mental
phenomena, just as chemistry of the bodily. As
chemistry seeks the units, the indissoluble atoms, of which all
substances are composed, so psychology seeks the units of consciousness,
the ultimate elements, on whose combination
the life of consciousness rests. One can form an idea
of such a psychological unit (unit of consciousness) by
conceiving it in analogy with a strong and sudden
impression. In the case of such (a blow, an effect of light or a sound) one cannot
in the first moment determine of what kind it is;
the nature of the sensation is known only when the impression
has a longer duration. Of such shocks or blows
every sensation is composed, just as a tone is composed
of vibrations of the air, which become perceptible
singly only when the rapidity with which they follow one another
does not exceed 16 in the second. All differences within
the life of consciousness rest on the different manners and
degrees of the combination of these psychological units.

The investigation of the material world leads us finally
to an objective unit, a material atom; the investigation
of the mental world leads us finally to a subjective
unit, a mental element. The question now arises how
they are related to one another. One will in vain try
to reduce the one to the other. We can only
think the material by means of the mental. The
most elementary material property is resistance, but this
we understand only by laying our own feeling and force
over into the external things. Only in our own movement, our own
energy have we a standard for the external phenomena; we
know things only through their reaction on us, but
this reaction we apprehend only by means of our own
action. Without muscular tension we should not be able to apprehend
matter. (Cf. above p. 133). If we therefore had
the choice between reducing the mental elements to the
bodily or the bodily to the mental, we should have to choose
% --- p. 159 ---
\setcounter{footnote}{0}%
\opage{159}the latter. It is absurd to trace the known back
to the unknown; an intelligible hypothesis we get only
when we trace the unknown back to the known and
thus regard the objective, material elements as in
their essence of the same nature as the subjective
elements of our consciousness.\footnote[1]{This way is taken by \emph{Lotze}. See ``Philosophien i Tydskland'' p. 300-302.} But yet neither can such a hypothesis
be carried through. Just as we can apprehend matter only by
means of mental determinations, so we can only
apprehend mind by means of material determinations. We
express the mental phenomena by designations that are
taken from the material phenomena. We never get, then,
beyond the symbolic. The conditioned form under
which the essence reveals itself in the subject can no more
than the conditioned form under which it reveals
itself in the object be the unconditioned essence itself, which lies
at the ground of both. Only so much do we see, that the same
law of evolution that holds within the one series of
the forms of manifestation of the Unconditioned also holds within
the other.

Every sensation must, as we have seen, be regarded
as composed of elementary impressions. But in order that a
sensation may form itself, not only must the elementary
impressions, the psychological units, be gathered into a unity; this
unity becomes a sensation only by being a member in a
chain of states of consciousness; only by its relation to
other sensations, successive, simultaneous or both
together, can the single sensation exist. Here
we come to one of the decisive differences between
physiology and psychology. Physiology occupies
itself with phenomena that proceed simultaneously, whereas
the phenomena that psychology investigates form a series.
The transition from physiology to psychology takes place through
the increasing concentration of the nervous system, whereby the
scattered impressions are gathered and brought into connection with each
% --- p. 160 ---
\setcounter{footnote}{0}%
\opage{160}other. Only when a central organ has formed, to
which the impressions can be conducted, does sensation become possible.
Simultaneity then passes over into succession, and it is precisely
an uninterrupted series of changes in the common center
that constitutes consciousness. This denotes a higher
stage of development; for when the impressions grow in number and in
composition, the adjustment to them cannot take place when
they occur at once; they would then work against and confuse
one another. Only when they form a series can this be
avoided. Consciousness arises, then, by virtue of the general
law of the development of life: the adjustment of the inner
to the outer.

The single sensation thus exists only by being
compared with the other members in the series of sensations and
ideas. When it is brought into closer connection
with one or more of these members, there forms itself
again a higher unity, and so on to infinity. The mode of composition
is the same throughout the whole structure of the
mind. The mind does not consist in the mere collection of
sense-perceptions and sensations; these must be organized and
distinguished in fixed groups, to one of which each new sensation can
be referred. On this rests what we call
thought and knowledge. In this its development the
mind presents the same features that have been shown to hold
for development in general. From the obscure sensation of sense,
resting on few recurring impressions with
little diversity and little combination, such as one finds
in the lower forms of animals, which have as yet no definitely
formed nerves and nerve centers, we come step by
step, as the organization becomes more perfect and the senses
more numerous, to greater combinations and further-reaching
articulation, or we come here also ``from an indefinite
and incoherent homogeneity to a definite and coherent
heterogeneity.'' The external senses (sight, hearing
and skin-feeling) stand in these respects above those that concern
the interior of the organism (emotions and visceral feelings); but
above all sense-perception stands the understanding in clearness and definiteness.
% --- p. 161 ---
\setcounter{footnote}{0}%
\opage{161}Thought or understanding is one with the sharply marked articulation of ideas.  The development of thought is
therefore the development of the relations and combinations of ideas.

This development must take place in the same respects
in which the development of life proceeds in general.
No living being can subsist except on the condition
of a harmony between the organism and its surroundings.  The
higher and richer life is, the more manifold and the deeper is
this harmony.  What a difference between the intestinal worm in
an organic tissue and a Newton or a Shakespeare, whose
thought embraces the world!  And yet there is a continuous
transition between these extremes.  The lowest stage in
the series is that at which the living being is entirely homogeneous
with the medium in which it lives.  A higher stage
is reached already in the zoophytes, which contract on
contact and thus show a succession of sensation
and movement.  Gradually the harmony extends in
space; immediate contact is no longer required
for a movement to occur.  When the impression of light
is concentrated in the eye, orientation of wider range becomes possible.
The carrier pigeon finds its way through the air; the helmsman guides
the ship across the ocean; the astronomer explores the space of the universe.
So too with regard to time; it is thus,
to take a single example, one of the marks of civilization
that a longer span of time is taken into account in
its calculations.  Men plant trees that will benefit
only their descendants, erect buildings that are meant
to stand for centuries, and in astronomy and geology cast
their gaze over immense periods of time.  In a similar
way the harmony between the mind and nature grows
in specialty, in composition and in systematic connection.\footnote[1]{Spencer here recalls \emph{Leibnitz}.  The relation between these thinkers has been stated by \emph{Ribot} in a fitting manner.  He finds two main points in common in their doctrines.  ``First, the idea of continuity or of universal compenetration, which has the effect that all things are joined together, that all things are causes and effects of one another, and that the manner \opage{162}in which the human mind separates them is arbitrary, although it is necessary.  This idea is really one with the idea of progress, and is only another expression of it; for when everything is transformed and assumes other shapes, everything is connected; there is in nature no hiatus, no break in coherence.  Only the idea of progress is chiefly dynamic, the idea of continuity on the other hand static. --- It is well known that Leibnitz, in order to explain the universe, devised monads, a kind of metaphysical atoms, which possessed all possible degrees from simple antitypy to perfect apperception.  According to their nature they constitute lifeless matter or the living being, man, animal or angel.  And since nothing in the universe is isolated, the whole universe acts upon every single monad, and the monad therefore expresses the universe, becomes a mirror which in its own way reflects the universe.  Now strip this great idea of its metaphysical form, and an indisputable positive truth will appear.  Place different beings --- a stone, a tree, a dog, a savage, a European, Newton or Shakespeare --- in the same medium, and each of them will reflect it in its own way and according to its nature, one very little, another to a high degree.  Between the being and its medium there will be what Herbert Spencer calls correspondence, and the degree of life will be measured by the degree of correspondence, the ideal form of life being perfect correspondence.  The man who could reach this summit would completely mirror the reality of the universe; he would be a microcosm adequate to the macrocosm.  This idea of correspondence, which is a main point in Spencer's psychology, seems to me a translation into the language of experimental psychology of Leibnitz's words: `every monad is a mirror reflecting the universe.'\,'' La psychologie Anglaise contemporaine. p. 149-150.}


% --- p. 162 ---
\setcounter{footnote}{0}%
\opage{162}Since all mental activity consists in an adjustment
of the inner to the outer, and since the outer relations
through imperceptible transitions grow in number, complexity
and diversity, the transitions among the
inner relations must be equally continuous and imperceptible.
Therefore the usual psychological classifications can have
only superficial validity.  Instinct, reason, observation,
concept, memory, imagination, will, etc. --- these
so-called faculties can be nothing but different degrees and forms
of the adjustment of the inner relations to the outer.  The lowest form
of mental life is \emph{reflex movement}.  It stands
on the border of the purely physical.  A simple inner relation is here regulated by
a simple outer relation; several reflex movements can go on
% --- p. 163 ---
\setcounter{footnote}{0}%
\opage{163}simultaneously; they therefore do not take place with consciousness.  When
reflex movement becomes more complex, it passes over into
\emph{instinct.}  Through the frequent repetition of several
psychical states (impressions and the movement that follows upon them)
there forms an hereditary tendency to perform these
complex functions.  The simpler connections form
first, being more frequent than the more complex ones.
But the more complex instinct becomes, the less
certain it is; the functions then gradually lose their
automatic character.  Several different functions will
then under certain conditions have a tendency to come into
force; and through the struggle between these tendencies there forms the
train of ideas that makes up the content of \emph{memory.}
Memory can in turn pass over into instinct,
namely when it is sufficiently firm and sure; an example
of this is found in every acquired skill.  In the same way
instinct passes over into what we call \emph{reason,} without any
definite boundary being assignable.  As Stuart Mill\footnote[1]{And, for that matter, already \emph{Leibnitz}.  De anima brutorum § 14: Consecutiones empiricæ nobis sunt communes cum brutis, et in eo consistunt, ut sentiens ea qvæ aliqvoties conjuncta fuisse expertum est, rursus conjunctum iri exspectet. --- But Leibnitz does not carry his lex continui through here, since this empirical inference is to be absolutely (toto coelo) different from the inference of reason (consecutio rationalis), which rests on an insight into the necessity of the relation.} has
shown, the original mode of inference is entirely instinctive.
When a phenomenon presents itself, one infers
involuntarily that its usual properties and other accompaniments
must also be present.  It is an inference from one particular
to another particular, such as even animals and small children make.
From this one passes by degrees to the more complex
(inductive and deductive) forms of inference.  Rational
thinking, like memory, through repetition
and practice often goes back to the instinctive.
Thus we believe that we see the distance between things immediately,
although only through long practice have we learned to
% --- p. 164 ---
\setcounter{footnote}{0}%
\opage{164}combine the various single ideas of which the idea
of distance is composed.

\emph{Will} also arises through the growing complexity
of reflex movements.  The difference between a voluntary
and an involuntary movement consists in this, that the former only
occurs after it has been represented in consciousness.
Because of the complexity of the mental states the
action cannot be completed at once; a partial execution precedes
the full execution, so that consciousness gets
time to occupy itself with it and to weigh the possibilities.
The cessation of automatic action and the origin of will
coincide. --- The will is free only in the sense that
the action proceeds from the individual's own state of consciousness.
On the other hand the usual doctrine of the freedom of the will is meaningless,
because it presupposes that one has a faculty of
action which can govern the preceding state of consciousness.
But this depends on the earlier experiences
that the individual has either made himself or has inherited in his
constitution.  If there were such an incalculable factor
as the free will is according to the usual conception, then
psychology would be impossible as a science,
but the whole great evolution that we can trace through
the world of nature and of mind would be fruitless.  This
evolution goes in the direction of harmony between the inner
and the outer.  Life must become higher and felicity
greater, because the inner relations are determined by the outer and
adjust themselves to them.  But if the inner relations were in part
determined by another factor, the harmony could be disturbed
at any moment and progress hindered. ---

What has thus been set forth by the way of \emph{synthesis},
--- the growth of mental life out of the most inconspicuous germs
through continuous transitions, and in such a way that it consists
throughout of the same elements, only at different
stages of development and in different combinations, ---
shall now be tried by the way of \emph{analysis.}  Having seen
the course of development from the simplest elements to the highest
complexity, we now begin conversely with the highest
% --- p. 165 ---
\setcounter{footnote}{0}%
\opage{165}forms of thought and seek to reduce them in such a way that
one and the same element proves to lie at the ground of them all.

As the highest form of thinking Spencer regards
the \emph{quantitative} or mathematical.  It treats
the most complex relations and yet with perfect exactness.
It presupposes homogeneity in its objects
and investigates only the relation between their magnitudes.  At the
ground of all mathematical investigations and proofs lies
the axiom that two things or relations that are equal to one and
the same third are equal to one another.  This axiom
must be known by intuition, by an act of thought that cannot
be reduced to any simpler form.  On such an act of
thought rest, each for itself, the successive transformations of
an equation.  We posit the equality here for no other
reason than that we are immediately conscious of it. ---
When the quantitative relations do not admit of any exact determination,
we stand at the transition \emph{to qualitative}
\emph{thinking,} which concerns the connection between things and
properties.  It is not mere magnitudes that are investigated here;
rather it is sought what general determinations
hold for certain phenomena, or how certain general
determinations find their application in the particular
cases.  The form of qualitative thinking is the inference,
which differs from the mathematical proportion only
by the lack of exact determinations of magnitude
; the hitherto unknown fact that one posits by means of
an inference corresponds entirely to the unknown fourth term in
a proportion.  Both inductive and deductive
inference rest in the end on a comparison between relations,
and in the primitive form of inference (aboriginal reasoning),
which lies at the ground of both, the inference from one particular
to another particular, we have an immediate reference of
one phenomenon to another phenomenon.  The act of thought that
lies at the ground is thus here too an intuition, an immediate
conception of likeness.\footnote[1]{In his doctrine of inference Spencer differs \opage{166}from Mill especially in that he sharply distinguishes between logic and the psychological doctrine of the process of inference.  The former is a science of universal objective relations, ``the conditions under which we can assume connections in time and space'', whereas the latter describes the subjective process by which we know those necessary connections.  Logic is thus a statics, whereas the psychological doctrine of thinking is a dynamics.  From this it follows that Spencer does not, like Mill, move in a nominalist direction, just as he in general takes a higher standpoint than Mill's empirical-subjective one.  See ``Principles of psychology'' 2d ed. II. p. 87-98. The classification of sciences p. 12.}


% --- p. 166 ---
\setcounter{footnote}{0}%
\opage{166}All thinking, both quantitative and qualitative, thus rests
on \emph{classification}, the referring of related
phenomena to the same group.  Every recognition is
a classification, indeed even the very \emph{conception} (\emph{perception})
of an object is one.  One cannot see
a thing without referring it to a certain kind.  There lies in
immediate intuition itself an inference: this thing
looks thus (e.g. it is red and round, hangs on a
tree of a certain kind, etc.), --- consequently it is such a
thing (an apple).  Inference is an indirect,
perception a direct establishment of a relation between two things. ---
Perception or conception itself is obtained by interaction
between the thing and the knowing subject.  All
the properties of the thing can be traced back to the different
positions in which the subject and the thing stand to one another.  But
these positions we in turn grasp only by means of the movement
of our own limbs and their mutual relations.  By comparing
these movements the intuitions of space and \emph{of time} arise.
All impressions are felt as successive.  When
two or more impressions can only follow one another
in a certain order, we have the intuition of time; if on the other hand
it is possible to reverse this order, we get the idea of
coexistence, out of which the intuition of space develops.  The latter
is thus, no more than the intuition of time, a form of intuition
ready-made from the start, but presupposes the conception of coexistence,
which in turn has the conception of succession as its
presupposition. --- It is always, then, a conception
% --- p. 167 ---
\setcounter{footnote}{0}%
\opage{167}of succession and change, a conception of difference and
likeness between states of consciousness, that lies at the ground as
the element common to all forms of consciousness, from the lowest to the highest.

The basic law of mental development is one with
the law of the association of ideas.  The primitive and
the most essential association is that between every single
sensation and the kind or class of preceding sensations
which it resembles and to which it is therefore referred.  This
association takes place entirely automatically.  It is not an
arbitrary act of thought, which may occur or not; but it
is one with the recognition of a sensation we have formerly
had, and, as already shown, a sensation cannot at all
arise in the mind except in so far as it is brought
into connection with its predecessors of like nature.
Only when the determination goes into the specific does a
deliberation take place, and thereby we pass beyond the automatic.


All thinking, from the most abstract and complex
inference to the most elementary intuition, consists
in positing relations of likeness and unlikeness.  Of these
the conception of unlikeness (difference, change, transition)
is in turn the underlying one.  For the conception of likeness presupposes
the conception of unlikeness.  When something is seen to
be like something else, we have three states of consciousness, of
which the third would coincide with the first if
the second (the transition itself) did not lie between.
These three states can be expressed thus: a, x, a,
(the three states of the relation of unlikeness, on the other hand, by a, x, b).  But
then there is a relation of unlikeness between a and x, a relation of unlikeness
that cannot be described otherwise than as a
change in consciousness.  The analysis cannot be carried
further. ---

Having seen how Spencer by the way of synthesis
and of analysis develops the laws and elements of mental life,
we naturally come to raise the question how
he stands toward the two main philosophical
% --- p. 168 ---
\setcounter{footnote}{0}%
\opage{168}directions and their attempts to solve the problem of the nature of
knowledge.

Spencer takes a critical attitude toward both directions.
The one regards, either, with Leibnitz, the inner
relations as created from the start in harmony with the outer,
or, with Kant, as subjective forms of intuition and of thought,
in which we receive and by which we grasp all impressions.
But in the first case evolution becomes meaningless, since
the harmony is present from the start.  The foregoing investigations
have on the contrary shown us harmony as the result
of an evolution, by which the inner relations are adjusted to the
outer and thereby ever assume higher forms.  And when
Kant assumes original forms of intuition and of thought, he commits
the same error on which the whole imperfection of the older
psychology rests, namely to regard only
the consciousness of the fully developed and civilized man,
instead of following the life of consciousness back to its
primitive forms.  Then it would, as we have seen, appear
that space and time are not original forms, but that they,
like all other intuitions and concepts, spring from
the consciousness of difference and likeness.  This consciousness
thus remains as the sole original form, for
intuition, understanding and reason alike.  When Kant further
assumes the forms of knowledge to be merely subjective, so that
the thing in itself is either formless or has forms that
do not affect us at all, then consciousness is really
not determined at all by things, but develops purely from within,
and Kant's doctrine thus leads us back to Leibnitz's
pre-established harmony.  Both sin against the law of
evolution, which permits no separation of the inner from the
outer, but according to which thought grows in its interaction
with things.

So far empiricism is right.  All knowledge
derives from experience.  The outer relations produce the
inner.  This is seen from the most common facts
: Experience makes wise; habit is second nature; coherence
and agreement among the impressions produce firm
% --- p. 169 ---
\setcounter{footnote}{0}%
\opage{169}connections between the ideas.  But the hypothesis of experience
in its usual form is no less untenable than
the hypotheses of Leibnitz and Kant.  If the mind, as
Locke taught, were a blank tablet, whence then should
the faculty of organizing experiences come, and whence
should the differences of degree in this faculty among
different species of organisms and among different individuals
within the same species come?  Why could a cat then
not be turned into a dog by education?  And why was a nervous system
necessary at all?  Leibnitz and Kant
are right that without an original organization the faculty
of ordering and working up impressions and of performing corresponding
actions cannot be explained.  Experiences form only the
material which the individual works up.  Still more clearly does
this show itself in the feelings.  Feelings and passions manifest themselves
before so extensive an experience has been had that their development
can find its explanation in it.  Not to speak of the
manifold passions that already manifest themselves in the child,
the mightiest of all passions, the passion of love,
precedes absolutely all experience.

Spencer combats pure empiricism, as
Stuart Mill especially has consistently developed it.  He maintains
in particular that experience cannot ground itself, but
presupposes a criterion, a principle which stands above experience,
a principle by virtue of which every step in
the development of knowledge is taken, and without which nothing can be
done.\footnote[1]{See above p. 14; 34.}  This last, abstract criterion is: the impossibility
of the opposite.  What we cannot, despite every possible effort,
combine in our thought, we must reject.  Every belief
or assumption is the result of a test of strength between the
different combinations of ideas, a systematic
struggle which serves to show which ideas
are inseparably combined.  What we cannot think apart
\emph{must} belong together.  And by these necessary
% --- p. 170 ---
\setcounter{footnote}{0}%
\opage{170}connections of thought we then judge everything else. --- The fixed
and universal relations between things must in the course of
evolution form fixed and universal connections in the organism;
through the infinite repetition of absolute outer
uniformities there are engendered absolute inner connections, necessary
forms of knowledge.  In the course of evolution
there may of course arise false connections of thought, which are only
the product of a limited and one-sided interaction with
things; but these connections, which for a time may
seem inseparable, are dissolved in the further progress of evolution.
That does not shake the principle: there \emph{are} inseparable
connections of thought.  But it is another question \emph{which}
they are. --- The inseparable connections of thought express the net
result of all experience down to the present time,
perhaps of the experience of millions of generations, and thereby acquire
a value that far surpasses the value of a demonstration
resting only on individual experience.

The result is that all knowledge and all forms of consciousness
derive from experience, but not from individual
experience.  Each single individual does not begin absolutely afresh,
but receives as an inheritance from his forefathers an accumulated and
fixed experience, which has guided the organization of the race in a
certain direction and toward certain forms.  There is an a priori,
but only for the individual; what is a priori for the individual
is a posteriori for the race.  In the course of the race's evolution
the outer relations so act upon it that the general
forms under which things present themselves obtain a
corresponding expression in the inner organization, especially in the nervous system.
Definite nervous connections are formed, which
are potentially present before birth, which precede and
are independent of all individual experiences, and which automatically
manifest themselves already in the first acts of knowledge.


As for the actual point in dispute between the two philosophical
schools, Spencer sides with the idealistic
direction, in that he denies that consciousness is
passive, and that knowledge can be explained out of the individual's
% --- p. 171 ---
\setcounter{footnote}{0}%
\opage{171}experience.  In order that experience may become knowledge, it must be
formed and organized, and there is thus presupposed an organizing
force and forms in which the organization proceeds.
This agrees precisely with Kant's doctrine.  Spencer
departs from Kant here only in reducing all forms of intuition
and of thought to one original function, the conception of
difference and likeness.

But Spencer does not stop at the form under which
the problem has hitherto always been treated.  He
carries the investigation further by asking whence the mental
organization derives.  This question had earlier
not been drawn into the debate.  The discussion was closed
once the necessity of such an
original organization had been shown as a presupposition of all experience
and knowledge.  By showing this presupposition
one had determined the most essential mark of man.
In the same way one stopped in botany and zoology when
one had shown the essential marks of a
certain species, and did not ask about the origin of this species itself and of its marks.
But just as Darwin has brought
this question forward within natural science, so
Spencer (independently of and before Darwin) has brought it forward
within philosophy.  The law of evolution, which is grounded
in the nature of our consciousness, and which experience in every
domain confirms, must find its application here too.
The higher forms of consciousness have developed out of the lower;
the very life of consciousness has developed out of the lower
forms of life; and in the course of this evolution, which consists in
the adjustment of the inner to the outer, the outer universal
relations have impressed their forms upon the inner relations.
Spencer can therefore declare himself an adherent of empiricism,
if by experience one understands the experience of the whole
race.

To this it must nevertheless be remarked that Spencer's own criticism of
empiricism applies just as well when the question concerns
the experience of the race as when it concerns the experience of the individual.
A real necessary and universal can be reached at
% --- p. 172 ---
\setcounter{footnote}{0}%
\opage{172}no point whatever, when experience is taken in the sense of empiricism.
If there does not lie at the ground, also as regards the race
as a whole, an appropriating and organizing activity,
it cannot be explained how the forms and
functions that form in the course of evolution could
become elements of our nature.  If one says: a priori for the individual,
one must also say: a priori for the race, for
what is in question here is a fundamental relation following from the nature of the thing, not an accident that could stand
in one way for the individual, and in another for the race.  ``The race'' is moreover only a relative
determination; it can in turn be conceived as an individual in relation
to a more inclusive race, and so on to
infinity.

That Spencer has not clearly seen that what is at stake here
is a question of principle has its ground in his standpoint
on the border of empiricism.  And yet it lay close for him
to notice it.  He certainly still stands on the ground of
``Association psychology''; but the associations of ideas
no longer stand with him as mechanical connections, but
appear as forms and results of the fundamental mental activity.
Mechanical psychology is here as it were
melted.  Activity and passivity merge into one, in that every
single sensation becomes possible only through that fundamental activity.
But herewith the principle is given: without inner
activity no appropriation of the outer.  This leads us
straight to the process on which, according to Spencer, the evolution of life
rests, the adjustment of the inner to the outer.  It is
the inner that adjusts itself; hence it must be conceived as
active, and the outer then supplies the aim and manner for the unfolding
of the activity.  But the inner always retains priority,
however far back we go.

What is now in the last instance this inner, which all knowledge
and experience presuppose, and which is thus a
principle beyond which we cannot go?  It can, according to Spencer's
own doctrine, be nothing other than the Absolute, the mysterious
power that reveals itself in all things.  Only through it
% --- p. 173 ---
\setcounter{footnote}{0}%
\opage{173}are life and consciousness possible.  It is therefore also the
last source of truth, the last guarantee of the validity of the
functions of knowledge.  The highest necessity
to which thought leads us is only an expression of our highest mental
power, and every power is again only a form of the
absolute power.  The Absolute is thus, as among the post-Hegelian
philosophers Weisse especially has maintained,\footnote[1]{Philosophien i Tydskland. p. 187---89.} the last
ground of the certainty of thought.  Spencer has himself
expressed this conviction when he regards the progressing thinker
as an organ of the absolute cause and derives his belief
from this source (see above p. 131).  One sees how
the deficiency in Spencer, pointed out earlier, in the conception of
the Unconditioned in its relation to the phenomena, becomes of
far-reaching significance also in the treatment of the special problems.

Apart from what is deficient in the treatment, it must be
acknowledged as a great merit that Spencer has initiated
a consideration of the problem of knowledge from this new side.
Just as surely as man is a relative and finite being,
the question must necessarily also be asked about the
origin of the determinations of his being.  Only a mystical and hyperidealistic
view can reject this question.  The Absolute and
that which is valid in itself is not shaken thereby; for the conception we
can have of it can after all only be relative, cannot
itself be absolute; and it is only about the origin of this peculiar
relative mode of conception that the question is raised.

Spencer's doctrine of the origin of the forms of knowledge
from the inherited results of the race's experience is a hypothesis, just
like Darwin's doctrine of the origin of species as a whole.  This
must be firmly held.  His criticism of empiricism, on the other hand,
is no hypothesis, but the result of a rational
psychological analysis.  By that hypothesis he himself believes
that he belongs to empiricism; by his criticism he has on the other hand
in principle gone beyond it. ---


% --- p. 174 ---
\setcounter{footnote}{0}%
\opage{174}There is a great class of psychological phenomena
which Spencer has treated only very briefly; these are the feelings,
which, as already remarked, are perhaps the most difficult
subject within psychology.  Spencer gives here only a
sketch (Princ. of psychol. II. p. 508 ff.), with the special
aim of showing the psychological foundation for ethics
and sociology.  Here as on all points the idea of evolution
is the ruling one.

At the lowest stage the feelings are immediately tied
to sensory sensations; but gradually they move away from
this point of attachment and assume a more ideal character.
This development moves through a series
of stages.  Freedom and mobility become greater; the
original automatic sensations, the moods that
immediately force themselves upon one, give way to a receptivity
for new and more comprehensive points of view.  But above all
evolution leads to the sympathetic feelings gaining the upper hand
over the egoistic.  The struggle for existence made
primitive man (the aboriginal man) a robber.  He
had, in order to survive in face of his many enemies, to be equipped
with savage instincts.  A developed sympathy would
not only have been impossible, but also ruinous to
him.  But indirectly this savage life already paves the way
for sympathy.  If egoism is necessary for the preservation of the individual,
then sympathy is necessary for the preservation of the species,
especially among weaker beings; therefore we see
herbivorous animals in particular live in company.  The constant companionship
with kindred beings, common fear and common joy nourish
sympathy and gradually make it a lasting element
in the nature of individuals.  Besides the relation between beings
of the same species, the relation between the two
sexes and between parents and children now also works in the same direction.  The
causes pointed out here lead from the dominion of the egoistic
feelings to a higher stage, where the ego-altruistic\footnote[1]{The expression altruism Spencer has borrowed from Comte to denote the opposite of egoism.}
feelings rule.

% --- p. 175 ---
\setcounter{footnote}{0}%
\opage{175}Sympathy does not, namely, at once displace egoism;
on the contrary, it itself first works to develop and maintain
new forms of egoism.  Sympathy is called forth by
life together with others.  But this union calls forth
precisely a stronger and sharper opposition to those
who stand outside this life together.  This shows itself in the relation of one nation,
one religion, one estate to the
others. --- Within society too egoism works in
close connection with sympathy.  For the individual gets
the motives for his acting from the authority which society
exercises over him.  Fear of the chieftain's power and of the
wrath of the tribal gods brings the individual to do what
counts as duty, which is what in the tribe's
history has proved to be useful and necessary.  To
such actions praise and reward are attached, to their
omission punishment.  Hero worship thus becomes a necessary
stage in evolution, an expression of a close
fusion of egoistic and altruistic feelings.  Even
if the reward or punishment that is expected is at first assumed
to come in another world, the motives are nevertheless worldly,
though the worldliness is what has aptly been
called ``other-worldliness.''  Actions are not regarded
according to their inner nature, apart from nearer or more remote
regard to the acting self.  The ego-altruistic feelings
thus mark a period of transition.  At the present stage of evolution
they are the only ones that can guide the conduct of men.
The inner, according to the general law of life, always adjusts itself
to the outer and cannot subsist in any other
way.  Every age, every state of society
therefore develops the feelings and ideas that
suit its conditions.  This is seen especially in the religious ideas.
These are necessarily anthropomorphic, and
the human forms that are transferred to the divine
are naturally precisely those that at a given time and place rule
in consciousness, because they are called forth by the interaction of
inner and outer relations.

% --- p. 176 ---
\setcounter{footnote}{0}%
\opage{176}We are still in the period of transition from savagery and
barbarism to real social life.  We often indeed think that
barbarism lies far behind us, and that our period is that of civilization;
but all earlier ages have thought the same,
and there was scarcely any barbarism in the past that
does not in one form or another still exist.  The life of robbers and
warriors, which formerly ruled, nourished especially the
egoistic feelings, whereas industrial life, especially through
the division of labor, awakens and nourishes the feeling of belonging together
and of community.  The great and savage crimes
therefore become rarer.  Sympathy prevents the direct
infliction of evil; the pain of others is felt more strongly than
before, --- though not so strongly that it is not often brought about
by roundabout ways.  Human nature still has
elements in it that derive from the savagery of primitive man.
This is seen, among other places, in children.  Their egoism, their
inclination to appropriate everything, their mistreatment
of lower beings are traits that suited the state
of the savage, but which must disappear as
humane social life develops.  The immoral rests,
like evil in general, on a discord
between life and its conditions; the evils, moral as well
as physical, which we suffer are signs that our nature has only
imperfectly adjusted itself to the conditions.

Man's original conditions brought it about that he had to
sacrifice the welfare of others for his own; his present conditions
demand more and more that he shall have only such wishes
whose satisfaction does not hinder the welfare of others.  Therefore
the egoistic feelings are gradually pushed back, and the
altruistic develop.  Fear, ambition and striving
for approval are no longer sufficient motives; the
moral sense has now developed in such a way that regard
for the welfare of others becomes the predominant one.  But altruism
always presupposes egoism, in that the altruistic feelings
consist in sympathetic awakenings of the egoistic feelings.
Sympathy in itself is therefore not sufficient as a principle of morality.
Adam Smith precisely overlooked that sympathy always
% --- p. 177 ---
\setcounter{footnote}{0}%
\opage{177}presupposes that one has oneself experienced that for which one has
compassion with another.  Pity, e.g., presupposes
that one has oneself felt pain.  One of the most complex
altruistic feelings is the feeling of justice.  It consists
in a sympathetic awakening of the drive to assert one's own
rights.  Whoever does not himself know and strive after
personal freedom cannot and will not recognize that of others; hence
is explained the tyranny of servile natures.  And on the other
side an absolute altruism, even if it were possible,
would presuppose that there were always some who were egoistic enough
to take advantage of the self-surrender of others and the resulting
self-annihilation. --- Evolution must move toward a
state where everyone works for the welfare and right of others
as for his own.  The ego-altruistic feelings correspond to
the transitory states in society, the altruistic to
the abiding ones, because they specify the necessary
conditions of social life.

Psychological development itself thus leads us to
the ethical principle.  This principle cannot be isolated
from the principle of life in general.  Morality is transcendental
physiology.  The course of life, the adjustment of the inner to
the outer, leads humanity from barbarism to moral
social life, and the moral law is the law that would
hold in the perfect society.  Therefore it is no
proof against an ethical system that it does not agree with
reality; it would on the contrary be a proof against it
if it agreed with reality, since the law
for the ideal state cannot be fulfilled under the
imperfect states, in the transition from animality and
barbarism to humanity.  But even if the ethical ideal cannot
be realized under the given conditions, it is nevertheless
of great importance to know and hold fast to it, so that evolution
may proceed in the direction of it and not away from it.
It was therefore a great error when thinkers formerly sought
the foundation of the moral principle in a state of nature that preceded
social life.  Such a state is precisely
in complete opposition to the ethical principle.  To
% --- p. 178 ---
\setcounter{footnote}{0}%
\opage{178}find this one must not look backward, but forward.  Absolute
ethics follows from the basic laws of life, from the conditions
of the highest development of human nature.  (Absolute
morality = conformity to the laws of complete life).

But on the other hand we must nevertheless constantly take
account of the actual state and character of humanity,
of all the imperfections in the present conditions, which
hinder the realization of absolute ethics.  We must proceed not
only a priori and deductively, but also a posteriori and inductively.\footnote[1]{Essays II. p. 210-211.}  The principle of utility then supplements here the
ideal moral principle.  Advancing civilization necessarily rests
on a series of compromises between old
and new, between the ideal and the practical.  In our
striving to do what is just fitting for the present
time, thought must nevertheless always be directed toward
what is best viewed abstractly.  The principle of utility is a necessary
supplement, though not the highest principle.

It is not the knowledge of the useful that calls forth
the moral feelings, but conversely the moral feelings
that make the knowledge of the useful possible.  The principle of utility
therefore always presupposes the deeper motives that spring
from the nature of the race, in which the individual has a share.  Utilitarianism,
like empiricism in general, sets the individual as an isolated
being of understanding outside the race, the evolving
race, and therefore cannot grasp what is original in the nature of the individual.
To learn that sympathy is beneficial one must first
act sympathetically.  Only by realizing this do we discover the
true motive, which lets the action be performed for its own
sake, whereas the utility theory lets the end fall outside
the action.

The utility theory moreover leaves undecided what utility
or happiness really is.  Happiness must indeed be regarded
as a divine idea that is to be realized.  But it
can consist only in the exercise and development of our faculties.  It
% --- p. 179 ---
\setcounter{footnote}{0}%
\opage{179}is our duty to develop our faculties, and we have a right to
the freedom by which this development alone becomes possible.
The utility theory thus here too points back to a more
comprehensive view, to the general idea of evolution.\footnote[1]{Stuart Mill relates (Moral p. 82) that Spencer has informed him that he does not wish to be regarded as an opponent of the utility theory, although he has combated it both in his ``Social Statics'' and in his Psychology.  The utility theory is indeed also, as the foregoing exposition shows, a \emph{part} of Spencer's ethical doctrine; but according to Bentham and Mill it was to be the \emph{whole} of morality.  Besides in this question of principle there is also this difference between Spencer and the Utilitarians, that the latter proceed purely empirically, whereas Spencer connects ethics, like sociology, with the general laws of life and thereby gains the possibility of a deductive method within ethics.  This is connected with his systematic tendency and comprehensive outlook.}

Not only in utilitarianism, but in Comte too Spencer finds
an overvaluation of knowledge at the expense of feeling.  In opposition to Comte's doctrine (in his ``Philosophie
positive''), that ``Ideas govern the world,'' that ``the whole
social mechanism rests in the end on opinions,'' Spencer
teaches that the world is governed by the feelings, for which the ideas
serve only as guides, and that the social mechanism almost
exclusively has the characters as its foundation.  ``Men's
drives are in the main inherited, but their
convictions are in the main self-acquired and depend
on the surrounding conditions.'' ``In reality it depends
on the character of the people and the state of society which ideas
are to spread.  The chief cause of social progress
is therefore the change in the nature of men which
is brought about by the constant discipline of social life, bringing it
more and more into agreement with social relations.''
According to Comte it is ``intellectual anarchy,'' according to Spencer
the moral oppositions, that bring about the political
crises of our time.\footnote[2]{The classification of the sciences. p. 37-38.}

This conception springs from Spencer's psychological
basic view.  We have seen that he especially maintains
% --- p. 180 ---
\setcounter{footnote}{0}%
\opage{180}the originality of the feelings, because they, even less than the forms
of knowledge, can be the fruit of the single individual's
experience and labor.  It is therefore especially through
the feelings that the individual is affected by the earlier experiences of the race,
and the qualities and traits of character that develop
through them are too strong to be changed by arbitrary measures and influences.  Spencer draws
from this important sociological consequences.  In his opinion
one overvalues the effects that external reforms, especially
in education and instruction, can have.  Above all
it is a matter of changing the basic qualities, the elements of
characters --- but these are changed only through the
practical struggle and interaction with the conditions of life.  The development of
society is a natural process, whose results one
cannot hasten in an artificial way.  It is an illusion,
which the Radicals share with the Conservatives, that legislation
is omnipotent.  The laws must indeed emerge from human nature
and can thus not be exalted above the intellectual
and moral imperfections from which this still suffers.
Nor can education be reformed so long as the parents
are still so far behind.  ``All superstition dies
slowly, and we fear that this belief in the omnipotence of government
will form no exception.''

``A scientific study of sociology leads to
a clearer judgment of the different political and religious
parties.  The conception which social science grounds
and develops is at once radical and conservative --- radical
in a degree that surpasses all that ordinary radicalism
grasps, and conservative in a degree that surpasses the
ideas of present-day conservatism.  When one has
appropriated the truth that societies are products of
evolution, which at different times and in different
places assume different forms in structure
and function, --- then arises the conviction that what
to our thoughts and feelings stands as extremely bad
institutions yet suited certain conditions that made better institutions
impossible; and from this there springs a
% --- p. 181 ---
\setcounter{footnote}{0}%
\opage{181}tolerant interpretation of the tyrannies of the past, at which even
the most zealous Tory of our day would be scandalized.  On the
other hand there springs from the observation of how the
processes that have brought things to their present state
always go on: not with a decreasing rapidity,
which might indicate that they were approaching their conclusion, but with
an increasing rapidity, which indicates long duration
and immense transformations, --- there springs from this observation
the conviction that the distant future holds
in its bosom forms of social life that are higher than we can
conceive; and this belief goes far beyond what the
Radical believes, who strives after a change that can
be compared with organizations that already exist.
And while this conception of society leads us to see
that immense changes can be brought about in the slow
course of things, so it also leads us to see that within
short periods only slight changes are possible.'' ``\footnote[1]{The study of sociology. London 1874. p. 399-400.}

One sees that Spencer does not, like Mill, step down into
the arena to take part in the struggles of the time, but keeps to
a purely theoretical and scientific standpoint.  His
striving aims precisely at freeing sociology from all
subjective and party points of view, something that naturally, as
he himself acknowledges, is possible only approximately, since
the investigator is himself a member of the society whose conditions of life
and course of development, whose natural history he investigates.\footnote[2]{As an inductive foundation for ``Principles of Sociology'' Spencer, in association with several collaborators, has begun to publish a ``descriptive sociology,'' a tabular survey of the different civilized societies, their forms and their development, which is to facilitate a comparative study of the sociological phenomena. --- The work cited above, ``The study of sociology,'' presents the general task of sociology and the subjective and objective difficulties of its treatment.}
That this striving of his after scientific
objectivity and his skepticism toward proposals of reform has not
its ground in lack of insight into the defects of our
present social conditions is seen not only from what has just been cited,
% --- p. 182 ---
\setcounter{footnote}{0}%
\opage{182}but also from utterances such as these: ``Ethically viewed there has
never been any other justification for the submission of the many
to the few except that it promoted the welfare of the many.
--- Just as the old theory of political government
has been transformed so that the ruler, instead of
being the nation's proprietor, is regarded as its servant, so
the old theory of industrial and social government
will undergo a change that will make the ruling classes
feel that their interests are only subordinate in relation to
the interests of the masses whose labors they direct.''\footnote[1]{Ibid. p. 256.}

According to Spencer ethics has an absolute and ideal character.
It is therefore ignorant of evil and its
consequences.  The ethical law, which is one with the law of
evolution, itself leads precisely beyond evil.  Progress is
a necessity; civilization proceeds as a natural process.
Imperfection means nothing other than a
discord with the conditions, which consists either in
one's having faculties that neither suit the conditions nor can
find application under them, or in one's lacking
the faculties that are required in order to subsist.  But from the general
law of life, that the inner adjusts itself to the outer and
thereby develops itself, it follows that the superfluous faculties die
away for lack of exercise, and that the necessary powers
gradually develop.  The realization of the ideal is therefore
logically certain.  Evil withers away, the good lives.\footnote[2]{Already Condorcet and I. G. Fichte have expressed this thought of ``the evanescence of evil.''}  One
must only let reality rule, let it exercise its education
and cut away the weak and the corrupt.

The transition period from animality to humanity must necessarily
be unhappy, for unhappiness is the necessary
consequence of a discord between the inner organization and the
outer conditions.  These sufferings must be endured, although they
can of course be mitigated by the activity of philanthropy.
Attempts artificially to withdraw men
% --- p. 183 ---
\setcounter{footnote}{0}%
\opage{183}from the consequences of the circumstances serve only to protract the matter
and at the same time to deprive them of the gain the sufferings
were to bring: education and development.\footnote[1]{``If, instead of allowing men to come into contact with the real conditions of their position, we place them under artificial conditions, under false circumstances, they will adjust themselves to these, and will then in the end have to make a new adjustment to the real relations.'' Social Statics. p. 387. Cf. Spencer's work ``Education, intellectual, moral and physical.'' p. 113 ff.}  It is only
shortsightedness and softness that, over the momentary
evils, forget the great, universal end.  Even in the
greatest pain and distress we descry a struggling love,
a power that works toward health, goodness and happiness.  Everyone
must work along according to his faculties in the conviction
that nothing is lost; even a purely momentary mood of
nobility and benevolence adds an atom to the ideal edifice
on which work is being done. ---

Since Spencer's more special treatment of ethics and
sociology is not yet available, we shall not examine more closely
or criticize the general features of his practical
philosophy, which have now been brought forward.  It is clear that
his position toward empiricism in the practical domain
corresponds to that which he takes toward it in the theoretical respect.
The ethical law is maintained in its ideality, but is seen
at the same time as the highest natural law, the highest form
of the general law of life.  Hereby the ideal gets the
firm and real grounding which idealism has so often
let it lack.  It is the principle of evolution and
the law that holds for the working of this principle that also
unite the ideal and the real.  And the conception of the
ethical ideal becomes possible only through the infinitizing
power which consciousness possesses, but which empiricism wholly
subordinates to passive receptivity. --- Here too it
shows itself, then, that Spencer's view closes itself only in a
consistent manner, if we are right in what in the
present exposition and critique has repeatedly been
emphasized,
% --- p. 184 ---
\setcounter{footnote}{0}%
\opage{184}that the ideal element in it is the predominant one,
is the principle that rules and works in the empirical
reality. ---

\closinghead{Conclusion.}


The road we have traveled has led us to a
definite goal.  Empiricism, with its conception of consciousness
as passively receptive and with its dissolution of
the unity of mind into atoms, which only through the
laws of the association of ideas got a certain, though accidental connection, --- this
doctrine, which was an inheritance from the 18th century, we have
seen develop through its own consistency and through debate
with the critical and idealistic philosophy to a point
where not only the active unity of consciousness and of mind is recognized,
but where inner necessity takes the place
of outer accident.

The atoms into which the life of consciousness dissolved for the
empirical ``Association-psychology'' were the single sensations
or rather elements of sensation.  But it
proved that no single sensation was possible without an
elementary mental activity that grasped its difference from
other sensations and in one way or another combined
it with them.  The last factor to which the analysis
led was thus this very ideal activity, on which
the unity and coherence of the mind rest; it lies at the
ground of all sensation and experience, which only through it
are possible.

English philosophy has thus in its own peculiar
way worked its way forward to the great fundamental thought
which Descartes, Leibnitz, Kant and Hegel have developed,
and which modern natural science has so strikingly and grandly
confirmed, that the idea of subjectivity is the last to which
we come, and that in it the principle of the objective must
% --- p. 185 ---
\setcounter{footnote}{0}%
\opage{185}be sought.  But it has hitherto been the defect and one-sidedness
of continental philosophy that it has developed
the idea of subjectivity predominantly in an a priori and dogmatic
manner.  Thereby one cut oneself off not only from the a posteriori
verification, without which no scientific truth can
be maintained, but one also came to conceive the ideal
principle itself abstractly and isolated from real reality.  This
came clearly to light at the turning point which occurred
with Hegel's system, the grand attempt at an absolute
idealism, a presentation of the Idea in its absolute being
in and for itself (theologically expressed: of God as he was before
the creation of the world).

The strength of English philosophy lies in its sharp
and sure eye for the conditions, which has indeed often veiled
the view of the Unconditioned, but which in return has
made a more real and concrete conception possible.  It has
seen that when evolution does not take place in an ideal world of ether,
but in the world of reality and of relativity, then
a decisive weight, a principled accent, must be laid
on the single conditions and circumstances.  The mind
develops only through the content it receives.  Only
through a long struggle, a laborious process of development,
does the received content become a part of the nature of the mind.

The law of evolution thus shows us the Unconditioned in
its unity with the conditions, the inner in its necessary
relation to the outer.  Its most definite and pregnant
expression the idea of evolution finds in the law of life, as
Spencer has formulated it: the inner adjusts itself to the
outer and develops itself through this adjustment.  The
validity of the ideal principle shows itself precisely through its dependence
or apparent disappearance.

Only through the idea of evolution does English philosophy gain
coherence and completion.  We have sought to show
what contradictions arose, especially in Stuart Mill,
when this idea was not taken as the ground.  The single
real moment, or in the domain of psychology the single
sensation, remains, when it is conceived in isolation, dead and
% --- p. 186 ---
\setcounter{footnote}{0}%
\opage{186}ful; in order to have real validity and reality it must
be seen as a member of the whole coherence; but a member of the coherence
it becomes only through a process, and the law
holding for this process thus becomes the last ground of explanation
to which we can come.  By this law are reconciled
the determinations that, taken immediately and in isolation, conflict
with one another.  Thus the laws of the association of ideas conflict,
when they are conceived as mechanical and accidental, with the
unity and energy of consciousness; but this conflict is removed when
one sees them as laws of the development of the life of consciousness, as
expressions of the conditions by which the unity is developed and
enriched and the energy assumes higher forms.

It follows of itself that in most of the English
philosophers, even in Spencer, who has precisely made
evolution his fundamental idea and, supported by it, combats
pure empiricism, we must find traces and after-effects of
the empirical-mechanical mode of conception that forms the starting point
of English philosophy in our century.
This the foregoing exposition has repeatedly
emphasized, and we shall not dwell on it here, where
we wish only briefly to bring out the general results
we can draw from the exposition and critique now concluded.

English philosophy in our time is in its character
akin to the post-Hegelian philosophy in Germany and
has essential points of contact with it.  In England the
movement goes from empiricism and realism back to idealism;
in Germany one had speculative idealism
as a historical presupposition, but sought now to take a new
starting point in experience and in the real given by it,
in order by this road to return to the ideal view.
In both directions the experimental
starting point is thus recognized as essential, although it is recognized that
this starting point points back to deeper presuppositions
and is itself a member of a larger coherence determined by these presuppositions.\footnote[1]{In view of this the statement in my book ``Philosophien i Tydskland \opage{187}efter Hegel'' p. 323 must be corrected; it hits Mill, but not English philosophy as a whole.}
 One may say that idealism has defeated
% --- p. 187 ---
\setcounter{footnote}{0}%
\opage{187}realism, but that just as conversely realism has defeated
idealism.  From both sides we are pointed back to one and the
same fundamental idea.

The method in English philosophy is essentially inductive.
The logic of induction is indeed also precisely owed to this
philosophy.  The same method to which natural science
owes its greatness is also recognized more and more as
essential for philosophy.  Since philosophy, as the general
science, is to bring together the results of the special sciences,
this can happen only by the same method as is applied
within each of these themselves.  The last, strict
criterion for what is to be recognized as real can
be sought only in exact experience.  It is Stuart Mill's
great merit that he has not only formulated the concrete
methods of induction, but also, more penetratingly
than any other logician, maintained the necessity of verification.
He does not overlook the significance of deduction, but
insists that its results must always be tested in experience.
What he does not see, on the other hand, is that induction
itself presupposes principles and axioms, without which
it could build up nothing.  What the real is,
only experience can teach us; but that through experience we
reach the knowledge of the real is grounded in our mental
and bodily organization.  The inner is here too
the principle of the outer; it must indeed in experience adjust itself
to the outer, behave receptively toward it; but this
it can do only by virtue of its ideal originality; the way
in which it appropriates the outer rests on its own
previously given nature; and through the adjustment it develops
itself.  In its method as in its fundamental idea the English philosophy of
our time thus leads us beyond empiricism; and
in the question of method too it is by the law of evolution
that progress is made.

If it stands firm that we have a true knowledge
% --- p. 188 ---
\setcounter{footnote}{0}%
\opage{188}of the real only when we can verify it, i.e. demonstrate
its confirmation by facts given in experience,
then it is also shown that we cannot gain a real
knowledge of the absolute principle.  Our development of concepts
cannot be verified on this point.  Kant's ``Kritik eller
spekulativen Theologie'' stands so far in force even
to this day.  But when the analysis of the nature of our consciousness
and of the results of experience leads us to the presupposition
of such an absolute principle (Spencer); when the source
of the idea of the Infinite lies in the original activity of our mind
(Bain); when the idea of cause is an expression of the fact that we
with necessity see an inner unity and coherence
as the ground of phenomenal multiplicity and movement
(Hamilton); --- then this is not shaken by the impossibility
of verification.  We cannot deduce the factual
and the conditioned from the idea of the Unconditioned; but we can
develop this idea in such a way that it is seen as the highest
presupposition, toward which all knowledge points back, and
as the ideal toward which all development in the world of nature and of
mind points forward. ---

But it is not only by its main ideas and its
method that English philosophy is of great importance,
its merit rests not least on the series of
detailed investigations it has undertaken in the domain of logic and psychology.
Progress in the domain of general ideas
is always conditioned by progress in the treatment
of the concrete problems.  Stuart Mill's formulation
of the experimental method, Bain's descriptive psychology,
Hamilton's acute theory of the principle of causality,
Spencer's psychological and metaphysical analyses are models
of scientific investigation.  This tendency toward
detailed investigations English philosophy also shares with
the post-Hegelian philosophy in Germany; we recall only
the physiological and psychological investigations of Fechner and Lotze
and the investigations in the philosophy of religion of Feuerbach and Weisse.
The more one recognizes the significance of induction and verification,
% --- p. 189 ---
\setcounter{footnote}{0}%
\opage{189}the greater weight must also, in the nature of the case,
be laid on such special investigations.

We stand here at a point where the difference between
the schools loses its significance.  Through detailed investigations
as through the appropriation of the content won by
empirical science there forms gradually a common ground
for all philosophical inquiry.  So far philosophy
is on the point of entering what Comte calls the positive
stage.  That this entrance cannot, however, as Comte thought,
take place by giving up the deeper grounding and the comprehensive
view, which are reached only by the road of logical and metaphysical
analysis, is a truth by whose recognition
the English philosophers distinguish themselves from French positivism.
But a sharper consciousness of that common
ground, a more and more consistent carrying out of the
philosophical principles will, as history, and especially the history of
philosophy in our time, shows, also lead to greater
and greater agreement in the results.  Both in the fundamental ideas
and in the method we have thus found convergence
between the main philosophical tendencies of our time.  The schools
lose their significance the more, on the one side, the significance
of the results of exact experience, on the other side the
subjective character of speculative deductions and the impossibility
of their verification come to appear to consciousness.

The post-Hegelian philosophy in Germany had as
its historical presupposition absolute idealism and could
win conditions for further advance only by a break with it.  English philosophy, on the other hand, has developed
continuously from the last century into
this one and has paved its way to new results by
following the course it once entered upon while at the same time
being influenced by the results of continental philosophy.

Both directions point us, each from its own side, toward
the invisible center where truth dwells, which indeed
will never be grasped by finite thought, but which
nevertheless is the goal toward which it constantly strives, and
toward which it moves in infinite approximation.  It
% --- p. 190 ---
\setcounter{footnote}{0}%
\opage{190}has been the intention of the now concluded account of
these important tendencies within philosophical inquiry
in our time to show that work is also now being done and effort made toward
this goal.  The emphasis does not lie on the results won at a
given time; these constantly point beyond
themselves.  The law of evolution teaches us, on the contrary, that truth
lies in the eternal striving, and it is the calling of science and of
philosophy to awaken and nourish this striving, to
lead it forward to new tasks and new labors.

\end{document}
